
\include{Roach}

%                                  elementary-math-from-an-advanced-standpoint_roach
%\documentclass[12pt]{amsart}
%\documentclass{report}
%\usepackage{amssymb,latexsym}
%\usepackage{amsmath}
%\usepackage{amsfonts}
%\usepackage{graphics,graphicx}
\setlength{\topmargin}{-10pt}
\setlength{\voffset}{-40pt}
\setlength{\textheight}{700pt}
%\usepackage{amsscr}
%\usepackage[amsscr]{eucal}

%\newtheorem{theorem}{Theorem}
%\newtheorem{lemma}{Lemma}
%\newtheorem{proposition}{Proposition}
%\newtheorem{definition}{Definition}
%\newtheorem{corollary}{Corollary}
%\newtheorem{notation}{Notation}
%\newtheorem{remark}{Remark}
%\newtheorem{assertion}{Assertion}
%\newtheorem{claim}{Claim}

\newcommand{\abs}[1]{\left|#1\right|}
\newcommand{\norm}[1]{\left\|#1\right\|}
\newcommand{\p}[1]{\left(#1\right)}
\newcommand{\tail}[1]{#1^{\p{N}}}
\newcommand{\trunk}[1]{#1^{\left[N\right]}}
\newcommand{\dg}{^{\circ}}
\newcommand{\nl}{\newline}
\newcommand{\fl}{\flushleft}
%\newcommand{\tild}{\texttt{\symbol{126}}}
\newcommand{\tild}{\sim}
\newcommand{\up}{\wedge}
\newcommand{\dn}{\lor}
\newcommand{\lra}{\leftrightarrow}
\newcommand{\ra}{\to}
\newcommand{\qqu}{\qquad \qquad}
\newcommand{\onto}{\dashv}
\newcommand{\power}{\Pi}
\newcommand{\univ}{\textbf{U}}
\newcommand{\ns}{\emptyset}
\newcommand{\of}{$(F, +, \cdot , <)\,$}
\newcommand{\zp}{Z^{+}}
\newcommand{\fld}{ $(F,+,\cdot)$}
\newcommand{\po}[1]{\dfrac{\pi}{#1}}
\newcommand{\lpo}[1]{\frac{\pi}{#1}}
\newcommand{\q}{\indent}


%\pagestyle{myheadings}
%\thanks{I would like to thank Charles A. Coppin and William T. Mahavier for their help in the publishing of this book.  I would, also, like to give special thanks to Sam Creswell for doing the typesetting and preparing this book for publication.}
\begin{document}
%\acknowledgment{I would like to thank Charles A. Coppin and William T. Mahavier for their help in the publishing of this book.  I would, also, like to give special thanks to Sam Creswell for doing the typesetting and preparing this book for publication.}
%\begin{center}
%\begin{Huge}
%Algebra and Trigonometry \\
%from an \\
%Advanced Standpoint \\
%Here \\

%\end{Huge}

%$ $\\

%$ $\\

%\begin{Large}
%by Everett D. Roach
%\end{Large}
%\end{center}

%\vfill
%Mathematics Department\\

%University of Dallas\\

%Irving, Texas\\
\affiliation{
Mathematics Department\\
University of Dallas\\
Irving, Texas 75060\\
everett\_r@verizon.net}


%e-mail : everett\_r@verizon.net\\

%\copyright 1970 by Everett D. Roach
%\newpage

\title{Elementary Mathematics from an Advanced Standpoint}

\author{Everett D. Roach}
\maketitle
\thispagestyle{empty}
\vspace{3in}

\begin{Large}
\begin{center}
To my wife, Irene
\end{center}
\end{Large}
\newpage
\thispagestyle{empty}
\begin{large}

\textbf{ Preface }\nl


This book has been written for a two semester-hour course
to precede calculus at the University of Dallas. I tried to make most of the material easy enough for the average students, but parts of it challenging enough for the brightest students.  The topics discussed have been taken from the areas of mathematical logic, modern algebra, and trigonometry.\nl

There are two primary goals of this book.  The first is to introduce the undergraduate to the process of mathematical reasoning early in his college studies. The second is to develop the skills needed to understand calculus. It is recommended that the students try to prove the theorems for which the proofs are given before reading the proofs.  It is also recommended that to use the book as an algebra and trigonometry text one should start with Section 1.4 Sets and Relations.\nl

%I would like to express my extreme appreciation to Francis L. Hardesty and Joseph R. L. Simkins, without whose help in editing and printing this book would not have been published.\nl

%I would like to thank Charles A. Coppin and William T. Mahavier for their help in the publishing of this book.  I would, also, like to give special thanks to Sam Creswell for doing the typesetting and preparing this book for publication.\nl

%\hfill E.D. Roach

%\hfill Irving, Tex.
%\vfill
Acknowledgment for the 1970 printing of this book.\nl

I would like to express my appreciation to Francis L. Hardesty and Joseph R. L. Simkins for the typing, editing and printing of this book.  Without their help this book would not have been published.\nl

Acknowledgment to the 2013 online printing of this book.\nl

Thanks go to Charles A. Coppin and William T. Mahavier for making all the arrangements necessary for the online publishing of this book and to James Harvey for helping with the proofreading.  I would, also, like to give special thanks to Sam Creswell for doing the typesetting, editing and preparing this book for publication.\nl

\hfill E.D. Roach

\hfill Irving, Tex.

%\acknowledgment{I would like to thank Charles A. Coppin and William T. Mahavier for their help in the publishing of this book.  I would, also, like to give special thanks to Sam Creswell for doing the typesetting and preparing this book for publication.}
\newpage
\thispagestyle{empty}
\begin{center}
\textbf{Table of Contents}
\end{center}

$ $ \nl


\fl \textbf{Chapter 1 Logic and Sets} \newline

\fl 1.1 Truth Tables  \hfill 1 \newline
1.2 Statement Calculus \hfill 4 \newline
1.3 Predicate Calculus \hfill 11 \nl
1.4 Sets and Relations \hfill 16 \nl
1.5 Functions \hfill 19 \nl
1.5 Algebra of Sets \hfill 21\nl

\textbf{Chapter 2 Fields} \nl

\fl 2.1 Properties of a Field \hfill 26 \nl
2.2 Fractions  \hfill 30 \nl
2.3 Exponents  \hfill 31 \nl


\textbf{Chapter 3 Ordered Fields} \nl

\flushleft 3.1 Inequalities \hfill 33\nl
3.2 Intervals  \hfill 37 \nl
3.3 Absolute Value \hfill 39\nl
3.4 Mathematical Induction \hfill 44\nl
3.5 Algebra of Functions \hfill 51\nl
3.6 Inverse Functions \hfill 55\nl
3.7 Arithmetic \hfill 57\nl
3.8 The Binomial Theorem \hfill 58 \nl

\textbf{Chapter 4 The Real Number System} \nl

\flushleft 4.1 The Completeness Axiom \hfill 62\nl
4.2 Logarithmic Functions \hfill 65\nl
4.3 Exponential Functions \hfill 67\nl
4.4 Radicals \hfill 69\nl
4.5 Factoring \hfill 71\nl
4.6 Polynomials \hfill 73\nl

\textbf{Chapter 5 Trigonometry} \nl

\fl 5.1 Six Fundamental Identities \hfill 78\nl
5.2 The Sine and Cosine Addition Formulas \hfill 81\nl

\newpage
\thispagestyle{empty}

\textbf{Chapter 5 Continued}\nl

\fl 5.3 Applications of the Addition Formulas \hfill 83\nl
5.4 The Double and Half-Angle Formulas \hfill 85\nl
5.5 Finding the Value of a Trigonometric Function of a Number \hfill 88\nl
5.6 The Sum, Difference, and Product of Sines and Cosines \hfill 91\nl
5.7 Inverse Trigonometric Functions \hfill 95\nl

\textbf{Chapter 6 Complex Numbers and Vector Spaces} \nl

\fl 6.1 Complex Numbers \hfill 98\nl
6.2 Generalized Exponents \hfill 101\nl
6.3 Vector Spaces \hfill 104\nl
6.4 Measure of an Angle \hfill 106\nl
6.5 Triangles \hfill 108\nl
6.6 Lines \hfill 110\nl
6.7 Matrices \hfill 113  \nl
6.8 Systems of Linear Equations \hfill 117\nl

\newpage
%\thispagestyle{}
\setcounter{page}{1}
\textbf{Chapter 1}\nl
\lhead{Section 1.1}
\fl \textbf{Logic and Sets} \nl
%page1

Mathematics is a study of axiom systems. It is an art in which statements called theorems are proven as a logical consequence of a set of assumptions called axioms. Since theorems are proved using logic~ it is useful to discuss symbolic logic first. In Section 1.1, we will introduce logical connectives intuitively using truth tables. \nl
\fl \textbf{1.1 Truth Tables}\nl

A definition is the naming of a concept.
\fl \textbf{Definition 1.1.1.} A \textit{statement} is a group of words which is true or false but not both. \nl
\fl \textbf{Definition 1.1.2.} If $p$ is a statement, then the \textit{negation} of $p$, denoted by $\tild p$, read ``not $p$", a statement which is true if and only if $p$ is false. \nl
\fl \textbf{Definition 1.1.3.} If each of $p$ and $q$ is a statement, then the \textit{conjunction} of $p$ and $q$, denoted by $p$ $\wedge$ $q$, read ``$p$ and $q$", is a statement which is true if and only if $p$ is true and $q$ is true.\nl

\fl \textbf{Definition 1.1.4.} If each of $p$ and $q$ is a statement, then the \textit{disjunction} of $p$ and $q$, denoted by $p$ $\dn$ $q$, read ``$p$ or $q$", is a statement which is true if and only if $p$ is true or $q$ is true or both. \nl

%\fl \textbf{Definition 1.1.5.} If each of $p$ and $q$ is a statement, then the implication $p$ $\to$ $q$, read ``$p$ implies $q$", or ``if $p$, then $q$", is a statement which is true if and only if $( \tild  p) \dn q$ is true. \nl

\fl \textbf{Definition 1.1.6.} If each of $p$ and $q$ is a statement, then the \textit{biconditional} $p$ $\leftrightarrow$ $q$, read ``$p$ if and only if $q$", or ``$p$ is equivalent to $q$", is a statement which is true if and only if $p$ and $q$ are both true or both false. \nl

\fl \textbf{Definition 1.1.7.} The symbols $\tild ,\up, \dn, \to$ , and $\leftrightarrow$ are called \textit{connectives.} \nl

\fl \textbf{Definition 1.1.8.} Statements which contain 2 or more connectives are called \textit{composite statements.} All other statements are called \textit{component statements.}

\newpage
%page2

\fl \textbf{ Definition 1.1.9.} Suppose each of $p$ and $q$ is a statement. $\tild$ is called the \textit{main connective} of $\tild p$. If $*$ is one of the connectives $\up, \dn, \to  or \leftrightarrow$, then $*$ is called the main connective of $p* q$. \nl
\fl \textbf{ Example 1.1.1.} Suppose each of $p$ and $q$ is a statement. \nl
 1) $\tild$ is the main connective of $\tild(p \dn q).$\nl
 2) $\up$ is the main connective of $( \tild P) \up ( \tild q)$.  \nl
 3) $\leftrightarrow$  is the main connective of $(p \up q) \leftrightarrow q$.  \nl

\fl \textbf{Definition 1.1.10.} If $p$ is a statement which is contained in a statement $q$, then $p$ is called a \textit{substatement} of $q$. \nl

\fl \textbf{Example 1.1.2.} The statements $p \dn q, p \ra q, (p \ra q) \up q$, and $(p \dn q) \lra ((p \ra q) \up q)$ are all substatements of $(p \dn q) \lra ((p \ra q) \up q)$ \nl

\fl Definitions  1.1.2 -1.1.6 are summarized by the following table, called a truth table. \nl
$ $ \nl
\begin{center}
\textbf{Table 1.1} \nl
%P  q  t'Jp  r-pl\q  pvq  p-+q  -------~ p~qi
$ $ \nl
\begin{tabular}{ | l | c | c | c | c | c | c |}
\hline
$p$ & $q$ & $\tild$ $p$ & $p$ $\up$ $q$ & $p$ $\dn$ $q$ & $p$ $\ra$ $q$ & $p$ $\lra$ $q$ \\ \hline
T & T & F & T & T & T & T\\ \hline
T & F & F & F & T & F & F\\ \hline
F & T & T & F & T & T & F\\ \hline
F & F & T & F & F & T & T\\ \hline
\end{tabular} \nl
\end{center}
$ $ \nl

The heading of a truth table contains one or more composite statements with their component statements listed to their left. If there are $n$ component statements, there are $2^{n}$ ways of assinging truth values to the component statements. Each of these ways is called a truth case. There are $2^{n}$ rows in the body of the table, one row for each truth case. For each truth case, the truth value of each composite substatement appearing in the heading is placed below its main connective in the row containing that truth case. \nl

For example, we have an $F$ in row 3 column 1 and a $T$ in row 3, column 2 of Table 1.2. This means that if $p$ is false and $q$ is true, then $(p \up q ) \ra (p \dn q)$ is true.


\newpage
%page3
\begin{center}
\textbf{Table 1.2}\nl
\begin{tabular}{ | c | c | c | c | c | c | c |}
\hline
$p$ & $q$ & [($p$ $\up$ $q$)& $\ra$ & ($p$ $\dn$ $q$)] & $\lra$  & $\tild$ $p$ \\ \hline
T & T & T & T & T & F & F\\ \hline
T & F & F & T & T & F & F\\ \hline
F & T & F & \textbf{T} & T & T & T\\ \hline
F & F & F & T & F & T & T\\ \hline
\end{tabular}\nl
$ $\nl
\end{center}

\fl \textbf{Definition 1.1.11.} A \textit{tautology} is a composite statement which is true regardless of the truth value of its component statements, i.e., it is true in each truth case.
From Table 1.2, we see that $(p \up q) \ra (p \dn q)$ is a tautology. \nl

\fl \textbf{Exercises} \nl
Denoting the statement ``It will rain." by ``$r$" and the statement ``I will wash my car." by ``$w$", express each of the following composite statements symbolically. \nl

1. If I wash my car, it will rain. \nl

2. I will wash my car or it will rain. \nl

3. It will rain if and only if I wash my car. \nl

4. If it rains, I will not wash my car. \nl

5. It is not true that I will wash my car and it will rain. \nl

6. I will wash my car and it will not rain. \nl


Let ``s" stand for the statement ``I will study mathematics." and

``$e$" stand for the statement ``I will enjoy mathematics.". Express each
of the following symbolic statements in English. \nl
\begin{tabular}{ l l}
7. $s \ra e$ & 10. $\tild  (e \ra s)$ \\
8. $s \up e$ & 11.  $s \lra e$ \\
9. $\tild s \ra e$ & 12. $\tild s \up \tild e$ \\
\end{tabular}\nl

Which is the main connective in each of the following statements?  \nl

13. $\tild (p\up q)$ \nl
14. $p \ra (p \dn q)$ \nl
15. $(p \ra q) \lra  ( \tild q \ra \tild p)$ \nl
16. $(p \ra q) \up \tild q$ \nl
17. $(p \lra \tild q) \dn (p \up q)$ \nl

\newpage
%page4
\fl \textbf{Construct a truth table for each of the following statements.} \nl

\fl 19. $[(p \up q) \dn p] \ra \tild q$ \nl

\fl 20. $(p\up q) \dn ( \tild p \up \tild q)$ \nl

\fl 21. $(\tild p \ra q) \dn (p \up q) $ \nl

\fl \textbf{Prove that each of the following is a tautology.} \nl
\fl 22. $\tild (p \up q) \lra ( \tild p \dn  \tild  q)$ \nl

\fl 23. $r \up (p \dn q) \lra [(r \up p) \dn (r \up q)] $ \nl

\fl 24. $(p \ra q) \lra ( \tild q \ra \tild p) $ \nl

\lhead{Section 1.2}
\fl \textbf{1.2 Statement Calculus} \nl

In this section we will use the axiomatic method to formalize the
concepts introduced in Section 1.1. That is, we will take some terms
to be undefined. All other terms will be defined using the undefined
terms. We then choose certain statements, which are usually intuitively
obvious, as axioms. Other statements called theorems are proved as a
logical consequence of the axioms. \nl

There are two main reasons for using this approach. (1) The axioms
are, in general, easier to work with than the definitions they come from,
if there are any. (2) There are meanings which could be assigned to the
undefined terms other than the obvious ones. This means that the results
can be applied to more than one system. For example, the following
axiom system could be used to design switching circuits. \nl

There are many possible choices for the undefined terms. We will
take $\tild , \dn$, and \textbf{statement} to be undefined. \nl

\fl \textbf{Axiom 1.2.1.} If each of $p$ and $q$ is a statement, then $\tild p$, read
``not $p$", and $p \dn q$, read ``$p$ or $q$", are statements. \nl

\newpage
%page5
\fl \textbf{Definition 1.2.1.} Suppose each of $p$ and $q$ is a statement. \nl
a) $p \up q$, read ``$p$ and $q$", stands for $\tild (\tild p \dn \tild q)$. \nl
b) $p \ra q$, read ``$p$ implies $q$", stands for $(\tild p) \dn q$. \nl
c) $p \lra q$, read ``$p$ is equivalent to $q$", stands for $(p \ra q) \up (q \ra p)$. \nl

\fl \textbf{Axiom 1.2.2.} If each of $p, q$, and $r$ is a statement, then each of
the following is an axiom. \nl
a) $p \lra q$\nl
b) $(p \dn q) \lra (q \dn p)$ \nl
c) $[p \dn (q \dn r)]  \lra [(p \dn q) \dn r] $ \nl
d) $[p \dn (q \up r)] \lra [(p \dn q) \up (p \dn r)]$ \nl
%e) $(p \dn p) \lra p$ \nl
e) $p \lra \tild \tild p$ \nl
f) $(p \up q) \ra q $ \nl
g) $p \ra (q \ra (p \up q))$ \nl

\fl \textbf{Definition 1.2.2.} Suppose $p$ is a statement. $r$ is called the \textit{conclusion} of $p$ if and only if (1) there is a statement $s$ (called the \textit{hypothesis} of $p$) such that $p$ is $s$ $\ra$ $r$ or (2) $p$ is $r$, in which case $p$ has no hypothesis. \nl

\fl \textbf{Definition 1.2.3.} A statement $p$ is called a \textit{theorem} if and only if there is a list of statements (called a \textit{proof} of $p$) such that the conclusion of $p$ is the last statement in the list and if $s$ is a statement in the list, then one of the following holds: \nl
\fl J1. $s$ is the hypothesis of p. \nl
\fl J2. $s$ is an axiom. \nl
\fl J3. $s$ is a previously proved theorem. \nl
\fl J4. $s$ is a statement resulting from the application of a definition to a
preceding statement in the list. \nl

\fl J5. (Modus Ponens) There is a statement $r$ such that $r$ and $r$ $\ra$ $s$ are preceding statements in the list. \nl

\fl J6. (Equivalence) There exist statements $r, t$, and $q$ such that $r$ $\lra$ $t$ is a theorem or preceding statement in the list, $q$ is a preceding statement in the list, $r$ is a substatement of $q$, and $s$ is the result of re­placing $r$ in $q$ by $t$.\nl

\newpage
%page6

We will use Hyp, Mp, and Eq to indicate the use of Jl, J5, and J6, respectively. We will also adopt the convention that $p, q, r,$ and $s$ stand for statements and that the operations are to be performed in the order $\tild, \up , \dn , \lra$ ,  unless otherwise indicated by parenthesis. The order of $\up$ and  $\dn$ will always be denoted by parenthesis. \nl

For example $ \tild p \up q \ra p \lra p \ra \tild p \dn q$ means $[(( \tild p) \up q) \ra p] \lra p \ra ((\tild p) \dn q)]$. \nl

\fl \textbf{Theorem 1.2.1.} If $p$ $\ra$ $q$ and $q$ $\ra$ p are theorems, then $p$ $\lra$ $q$ is a theorem. \nl
\fl Proof. \nl
\begin{tabular}{ l l }

1. $p$ $\up$ $q$ & Theorem \\
2. $q$ $\ra$ $p$ & Theorem \\
3. $(p \ra q) \ra [(q \ra p) \ra ((p \ra q) \up (q \ra p))]$ & Ax 1.2.2(g) \\
4. $(q \ra p) \ra ( (p \ra q)  \up (q \ra p))$ & MP 1,3 \\
5. $(p \ra q) \up (q \ra p)$ & MP 2,4 \\
6.  p $\lra$ q & Def 1.2.1(c) \\
\end{tabular}
\fl \textbf{Theorem 1.2.2} If $p$ and $q$ are theorems, then $p$ $\lra$ $q$ is a theorem. \nl
 Proof. $p ,q$ and $q ,p$ are proofs of $p$ $\ra$ $q$ and $q$ $\ra$ $p$, respectively.
 By Theorem 1.2.1, $p$ $\lra$ $q$ is a theorem. \nl

\fl \textbf{Definition 1.2.4.} Let $t$ denote a theorem and $f$ denote $\tild t$. \nl

\fl \textbf{Theorem 1.2.3.} If $p$ $\lra$ $t$ is a theorem, then $p$ is a theorem. Proof. \nl
\begin{tabular}{l l}
1. $p$ $\lra$ $t$ & Theorem  \\
2. $(p$ $\ra$ $t$) $\up$ ($t$ $\ra$ $p$) & Def 1.2.1(c)\\
3. $[(p \ra t) \up (t \ra p)] \ra ( t \ra p)$ & Ax 1.2.2(f)  \\
4. $t \ra p$ & MP 2,3  \\
5. $t$ & Theorem \\
6. $p$ & MP4,5 \\
\end{tabular}
\fl \textbf{Theorem 1.2.4.} ($p$ $\up$ $q$) $\lra$ ($q$ $\up$ $p$). Proof.
\begin{tabular}{l r l l l l}
1.& $p \up q$  &$\lra p \up q$ &&&Ax 1.2.2(a)\\
2.    &&$\lra \tild (\tild p \dn \tild q)$ &$ $&$ $& Def 1.2.1(a)\\
3.    &&$\lra \tild (\tild q \dn \tild p)$ &$ $&$ $& Eq, Ax 1.2.2(b)\\
4.     &$q \up \ p$ &$ $&$ $&& Def 1.2.1(a)\\
\end{tabular}
\newpage
%page7
\fl \textbf{Theorem 1.2.5.} If $p$ $\lra$ $q$ is a theorem, then $q$ $\lra$ $p$ is a theorem. Proof. \nl
\begin{tabular}{ l l l l}
1. $p \lra q$ &&&Theorem \\
2. ($p$ $\ra$ $q$)  $\up$ ($q$ $\ra$ $p$) &$\quad $&$ $& Def 1.2.1(c) \\
3. ($q$ $\ra$ $p$) $\up$  ($p$ $\ra$ $q$) &$ $&$ $& Th 1.2.4  \\
4. $q$ $\lra p$  &$ $&$ $& Def 1.2.l(a) \\
\end{tabular}\\

\fl \textbf{Theorem 1.2.6.} Each of the following is a theorem. \nl
SC1. $\tild (p$ $\dn$ $q$) $\lra (\tild$ $p$ $\up \tild q$) \nl
SC2. $\tild(p \up q)$ $\lra$ $( \tild p \dn \tild q$) \nl
SC3. $(p \up q) \lra (q \up p)$ \nl
SC4. $[p \up(q  \up r)] \lra [ (p \up q) \up r]$ \nl
SC5. $[p \up (q \dn r)] \lra (p \up q) \dn (p \up r)$ \nl
SC6. $p \ra (p \up p)$ \nl
SC7. a)$(p \up p) \lra p$ \qquad \qquad b)$(p \dn p) \lra p$ \nl
%SC8. $\tild  f \lra t$ \nl
SC8. $(p\up q) \ra p$ \nl
\begin{tabular}{ l l}
SC9. $\tild f \lra t$  & \\
SC10. $p \ra p$&\\
SC11. a) $\tild p \dn p$ & b) $\quad p\dn \tild p$  \\
SC12. a) $(\tild p \up p) \lra f$ &b) $p\up \tild p \lra f$\\
SC13. $p \ra (t \up p)$    &\\
SC14. $a) (p \up t)  \lra p$ & $b)\quad (t \up p) \lra p$  \\
SC15. $a) p \dn f \lra p$ & $b)\quad f \dn p \lra p$ \\
SC16. $a) p \dn t$ & $b)\quad  t \dn p $ \\
SC17. $a) (p \up f) \lra f$ & $b) \quad (f \up p) \lra f$ \\
\end{tabular}\nl
SC18. $[p \dn (q \up \tild p)] \lra (p \dn q)$ \nl
SC19. $[p \up (q \dn \tild p)] \lra (p\up q)$ \nl
SC20. $[p \dn (p \up q)] \lra p$  \nl
SC21. $[p \up (p \dn q)] \lra p$ \nl
SC22. $q \ra (p \ra q)$ \nl
SC23. $\tild p \ra (p \ra q)$ \nl
SC24. $p \ra (p \dn q)$ \nl
SC25. $(p \up q) \ra (p \dn q)$ \nl
SC26. $(p \ra q) \lra ( \tild q \ra \tild p)$ \nl
SC27. $(p \ra q) \lra  [(p \up \tild q) \ra f]$ \nl
SC28. $(p \ra q) \lra [(p \dn q) \lra q)]$ \nl
\newpage
%page8
SC29. $(p \ra q) \lra [(p \up q) \lra p)]$  \nl
SC30. $(p \lra q) \lra ( \tild p \lra \tild q)$ \nl
SC3l. $[p \ra (q \ra r)] \lra [(p \up q) \ra r]$  \nl
SC32. $[p \ra (q \ra r)] \lra [(q \ra p) \ra (q \ra r)]$ \nl
SC33. $(p \ra q) \ra (p \up r \ra q  \up r)$\nl
SC34  $(p \lra q) \ra (p \up r \ra q \up r)$\nl
SC35. $(p \ra q) \ra (p \dn r \ra q \dn r)$\nl
SC36. $(p \lra q) \ra (p \dn r \lra q \dn r)$\nl
SC37. $(p \lra \tild p) \lra f$\nl
SC38. $ (p\ra q)\lra (\tild p \dn q \lra t)  $\nl

Proof of SC1. \nl
\begin{tabular}{l l l }
1. $\tild (p \dn q)$ &$\lra \tild (p \dn q)$ & \text{Ax 1.2.2(a)} \\
2. &$\lra  \tild (\tild  \tild p \dn \tild \tild q)$ & \text{Eq, Ax 1.2.2(e)} \\
3. &$\lra \tild p \up \tild q$ & \text{Def 1.2.1(a)}\\
\end{tabular}\nl

Proof of SC2. \nl
\begin{tabular}{ l l l}
1.  $\tild (p \up q)$ &$\lra \tild (p \up q)$ &Axiom 1.2.2(a)\\
2.  &$\lra  \tild \tild  ( \tild p \dn \tild q)$  &\text{Def 1.2.1(a)}\\
3.  &$\lra \tild p \dn \tild q$  &\text{Eq, Ax 1.2.2(e)} \\
\end{tabular}\nl

Proof of SC5. \nl
\begin{tabular}{l l l}
1. $[p \up (q \dn r) ]$&$\lra [p \up ( q \dn r)]$& Axiom 1.2 2(a) \\
2. &$\lra \tild [\tild p \dn \tild (q \dn r)]$&Def 1.2.1(a)\\
3. &$\lra \tild [\tild p \dn (\tild q \up \tild r)]$&Eq, SC1\\
4. &$\lra \tild [(\tild p \dn \tild q) \up (\tild p \dn \tild r)]$&Eq, Ax 1.2.2(d)\\
5. &$\lra \tild [\tild (p \up q)\up \tild (p \up r)]$&Eq, SC2\\
6. &$\lra \tild \tild [(p \up q)\dn (p \up r)]$&Eq, SC1\\
7. &$\lra (p \up q) \dn (p \up r)$&Eq, Ax 1.2.2(e)
\end{tabular}\nl

Proof of SC6. \nl
1. $p$ \nl
2. $p \ra (p \ra (p \up p))$\nl
3. $p \ra (p \up p)$\nl
4. $p \up p$\nl
\newpage
%page9


Proof of SC7(a)\nl

1. $p \ra (p \up p)$\nl
2. $p \up p \ra p$\nl
3. $(p \up p) \lra p$\nl

Proof of SC7(b)\nl
\begin{tabular}{l r l}
1.&$(p \dn p)$&$\lra \tild \tild (p \dn p)$\\
2.&&$\lra \tild (\tild p \up \tild p)$\\
3.&&$\lra \tild (\tild p)$\\
4.&&$\lra p$\\ 
\end{tabular}\nl

Proof of SC10. \nl
1.  $p$\nl

Proof of SC11(a) \nl
1. $p \ra p$\nl
2. $\tild p \dn p$\nl

Proof of SC12(b) \nl

\begin{tabular}{l r l}
1.&$p \up \tild p $&$\lra p \up \tild p$\\
2.&&$\lra \tild (\tild p \dn \tild \tild p)$\\
3.&&$\lra \tild(\tild p \dn p)$\\
4.&&$\lra \tild t$\\
5.&&$\lra f$\\
\end{tabular}\nl

Proof of SC13 \nl

1.  $t$\nl
2.  $t \ra (p \ra (t \up p))$\nl
3.  $p \ra (t \up p)$\nl


Proof of SC14(b) \nl
%Let $r$ stand for $(t \up p) \ra p$ and let $s$ stand for $p \ra (t \up p)$.\nl

1.  $(t \up p) \ra p$\nl
2.  $p \ra (t \up p)$\nl
3.  $(t \up p ) \lra p $\nl

\newpage
%page10

Proof of SC15(a)\nl
\begin{tabular}{l r l}
1.&  $p\dn f$ &$\lra p \dn f$\\
2.&& $\lra \tild \tild p \dn f$\\
3. && $\lra \tild \tild p \dn \tild t$\\
4. &&$\lra \tild (\tild p \up t)$\\
5. && $\lra \tild \tild p$\\
6. && $\lra p$\\
\end{tabular}\nl

Proof of SC16(a)\nl
\begin{tabular}{l r l}
1.& $p\dn t$ &$\lra p \dn t$\\
2. &&$\lra (p \dn t) \up t$\\
3. &&$\lra (p \dn t) \up (p \dn \tild p$) \\
4. &&$\lra p \dn (t \up \tild p)$\\
5. &&$\lra p \dn \tild p$\\
6. &&$\lra t$\\
\end{tabular}\nl

Proof of SC20 \nl
\begin{tabular}{l r l}
1.  &$[p \dn (p \up q)]$& $\lra [p \dn (p \up q)]$ \\
2.	&&$\lra (p \up t) \dn (p \up q)$ \\
3.  &&$\lra p \up (t \dn q)$\\
4.  &&$\lra p \up t$\\
5.  &&$\lra p$\\
\end{tabular}\nl

\textbf{Exercises}\nl

Give justifications of each statement in the proofs of the following theorems.\nl

\begin{tabular}{l l l}
1. Theorem SC6& 2. Theorem SC7(b)&3. Theorem SC10\\
4. Theorem SC11(a)&5. Theorem SC12(b) &6. Theorem SC13\\
7. Theorem SC14(b)&8. Theorem SC15(a)&9. Theorem 16(a)\\
&10. Theorem SC20 &\\
&&\\
11. Prove SC8&12. Prove SC17(a)&13. Prove SC18\\
14. Prove SC19&15. Prove SC21&16. Prove SC22\\
17. Prove SC23&18. Prove SC24&19. Prove SC25\\
\end{tabular}\nl



%\qquad Prove Theorems SC8, SC17(a), SC18, SC19, SC21, SC22, SC23, SC24, and SC25.\nl

\newpage
%page11
\fl \textbf{1.3 Predicate Calculus}\nl

\lhead{Section 1.3}
We will begin by giving intuitive definitions of the concepts used in this section. A variable is a symbol which may stand for any object. A formula is a statement which involves at least one object or variable. For example, ``$x$ is the father of $y$" is a formula in which $x$ and $y$ are variables. A formula may be thought of as a collection of statements, one statement for each assignment of the variables to objects. In the above example, one such assignment would lead to the statement ``Tom is the father of John."\nl

\fl \textbf{Definition 1.3.1.} If $p$ is a formula, $x$ is a variable, and $t$ is an object, then: \nl
 a) $\forall x~ p$ means ``For each $x, p$." The symbol $\forall$ is called the \textit{universal quantifier.} \nl
 b) $\exists x~ p$ means ``There exists an $x$ such that $p$." The symbol $\exists$ is called the \textit{existential quantifier.} \nl
 c) $x$ is \textit{free} in $p$ if and only if the truth value of $p$ is not the same value for each object $x$. \nl
 d) $p$ is \textit{true} if and only if $p$ is true regardless of which objects its variables stand for. \nl
 e) $p(t/x)$ is the result of replacing each $x$ in $p$ by $t$. \nl

Each theorem in this section can be proven from the preceding definitions. A much simpler approach would be to take $\tild ,\dn , \forall$ ,  \textbf{object}, \textbf{variable}, \textbf{formula}, and \textbf{free} as undefined and prove the results of this section from Axioms 1.3.1. and 1.3.2. We will use the latter approach. \nl

\textbf{Axiom 1.3.1.} If each of $p$ and $q$ is a formula, each of $x$ and $y$ is a variable, and $t$ is an object, then: \nl
 a) $\tild p, p$ $\dn$ $q$, and $p(t/x)$ are formulas. \nl
 b) $\forall x~ p$ is a formula in which $x$ is not free. \nl
 c) If $x$ is not free in $p$ and $q$, then $x$ is not free in  $\tild p, p$ $\dn$ $q$, and
$\forall y ~p$.\nl

\textbf{Definition 1.3.2.} Suppose each of $p, q$, and $x \in A$ is a formula, $x$ is a variable, and $t$ is an object.
\newpage
%page12
a) $p \up q$, read ``$p$ and $q$", stands for $\tild(\tild p \dn \tild q$).\nl
b) $p \ra q$, read ``$p$ implies $q$", stands for $(\tild p) \dn q$ \nl
c) $p \lra q$, read ``$p$ is equivalent to $q$", stands for $(p \ra q) \up (q \ra p$). \nl
d) $\exists x~p$  stands for $\tild \forall x \tild p$. \nl
e) $\forall x \in A, p(x)$ is read ''for each $x$ in $A, p$ of $x$ and stands for $\forall x(x \in A \ra p)$. \nl
f) $\exists x \in A \ni $p(x) is read ``there exists an $x$ in $A$ such that  $p$ of $x$" and stands for $\exists x(x \in A \up p)$. \nl

\textbf{Axiom 1.3.2.} If each of $p,q$, and $r$ is a formula, and $t$ is an object, then each of the following is an axiom: \nl
a) $p \lra p$ \nl
b) $(p \dn q)$ $\lra$ $( q \dn p)$ \nl
c) $[p\dn (q\dn r)] \lra [(p \dn q) \dn r] $\nl
d) $[p \dn (q \up r)] \lra [(p \dn q) \up (p \dn r)]$ \nl
e) $(p \dn p) \lra p$ \nl
f) $p \lra \tild \tild p$ \nl
g) $(p \up q) \ra q$ \nl
h) $p \ra (q \ra (p \up q))$\nl
i) $\forall x~ p \ra p(t/x)$ \nl
\qqu $(\tild  p) (t/x) \lra \tild p(t/x)$ \nl
\qqu $(p \dn q) (t/x) \lra p(t/x) \dn q(t/x)$ \nl
j) $\forall x~ p \ra p$ \nl
k) $\forall x(p \dn q)\lra (\forall x~ p) \dn q$, if $x$ is not free in $q$. \nl

\textbf{Definition 1.3.3.} Suppose p is a formula. $r$ is called the \textit{conclusion} of $p$ if and only if (1) there is a formula $s$ (called the \textit{hypothesis} of $p$) such that $p$ is $s \ra r$, or (2) $p$ is $r$, in which case $p$ has no hypothesis. \nl

\textbf{Definition 1.3.4.} A formula $p$ is called a \textit{theorem} if and only if there is a list of formulas (called a \textit{proof} of $p$) such that the conclusion of $p$ is the last formula in the list and if $s$ is a formula in the list, then one of the following holds: \nl
J1. $s$ is the hypothesis of $p$. \nl
J2. $s$ is an axiom. \nl
J3. $s$ is a previously proven theorem. \nl
J4. $s$ is a formula resulting from the application of a definition to a preceding formula in the list. \nl
J5. (Modus Ponens) There is a formula $r$ such that $r$ and $r \ra s$ are pre­ceding formulas in the list. \nl
\newpage
%page13
J6. (Equivalence) There exist formulas $r, t$, and $q$ such that $r \lra t$ is a
theorem, $q$ is a preceding formula in the list, $r$ is a subformula of $q$,
and $s$ is the result of replacing $r$ in $q$ by $t$.\nl
J7. (Generalization) There is a preceding formula $r$ in the list such that $s$ is $\forall x~ r$ and $x$ is not free in the hypothesis of $p$. \nl

\qqu Each axiom of the statement calculus is an axiom of the predicate calculus with the \textbf{word statement} replaced with the \textbf{word} \textbf{formula}. Therefore, each theorem of the statement calculus corresponds to a theorem in the predicate calculus. \nl

\textbf{Theorem 1.3.1.} If ($p \up s$) $\ra$ $r$ and $p$ $\ra \exists$ $x ~s$ are theorems and $x$ is not free in $p$ or $r$, then $p \ra r$ is a theorem. \nl
Proof \nl
\begin{tabular}{ l l}
1. $p$ & Hyp \\
2. $p \ra \exists  x~ s$ & Theorem \\
3. $\exists x~ s$ & MP 1,2 \\
4. $p \up s \ra r$ & Theorem \\
5. $p \ra(s \ra r )$ & Eq, SC 31 \\
6. $s\ra r$ & MP 1,5 \\
7. $\forall x (s \ra r$) & Gen \\
8. $\forall x(\tild s\dn r)$ & Def 1.3.2(b) \\
9. $\forall x( \tild  s\dn r) \ra ( \forall x \tild s) \dn r$ & Ax 1.3.2(j) \\
10. ($\forall x \tild s) \dn r$ & MP 8,9  \\
11. $(\tild  \tild  \forall x \tild s) \dn r$ & Ax 1.3.2(e) \\
12. $(\tild \exists x~ s) \dn r$ & Def l.3.2(d) \\
13. $\exists x~ s \ra r$ & Def 1.3.2(b) \\
14. $r$ & MP3,13 \\
\end{tabular}\nl

In an application of Theorem 1.3.1, we begin by justifying $p$ with Hyp and proceed to prove $\exists x~ s$. Thus we have $p\ra \exists x~ s$. Then we write Th 1.3.1 for the justification of $s$ (to indicate that Th 1.3.1 will be applied) and proceed to prove $r$. Since both $p$ and $s$ are logical consequences of $p \up s$, we have in effect proved $p \up s \ra r$. If $x$ is not free in $p$ or $r$, we can conclude, from Theorem 1.3.1, that $p \ra r$. We will illustrate the procedure in the proof of PC11. \nl

\textbf{Theorem 1.3.2.} If each of $p, q, x \in A$  ,and $r$ is a formula, each of $x$ and $y$ is a variable, and $t$ is an object or variable, then: \nl


\newpage
%page14
\begin{tabular}{ l l}
PC1. $\forall x~ p \ra p$ &\\
PC2. a) $p(t/x) \ra \exists x~ p$   &b) $p \ra p$ \\
PC3. $\forall x~ p \ra \exists x~ p$ &\\
PC4. $\tild \forall $ x~ p $\lra \exists x \tild p$ &\\
PC5. $\tild \exists x~ p \lra \forall x \tild p$ &\\
PC6. $\tild \forall x \in A, p(x) \lra  \exists x \in A \ni \tild p(x)$&\\
PC7. $\tild \exists x \in A \ni p(x) \lra \forall x \in A ~ \tild p(x)$ &\\
PC8. $\forall x \forall y ~p \ra \forall y \forall x~ p$&\\
PC9.  $\exists x \exists y ~p \ra \exists y \exists x ~p$&\\
PC10. $\exists x \forall y~p \ra \forall y \exists x~ p$&\\
PC11.  $\exists x [p \up q] \ra [ \exists x ~p \up \exists x~ q]$&\\
PC12.  $\forall x ~p \dn \forall x ~q \ra \forall x [p \dn q]$&\\
PC13.  $\forall x [p \up q] \lra [\forall x ~ p \up \forall x~ q]$ &\\
PC14. $\exists x [p \dn q] \lra [\exists x ~ p \dn \exists x ~q]$& \\
PC15. $\forall ~x[p \ra q] \ra [\forall x~p \ra \forall x ~q]$&\\
PC16.  If $x$ is not free in $p$, then :&\\
a) $p \lra \forall x~p$&\\
b) $p \lra \exists x~p$&\\
c) $\forall x[p \up q] \lra p \up \forall x~q$&\\
d) $\forall x [p \dn q] \lra p \dn \forall x~q$&\\
e) $\exists x [ p \up q] \lra p \up \exists x~q$&\\
f) $\exists x[p\dn q] \lra p \dn \exists x~q$&\\
\end{tabular}\nl

%Proof of Theorem PC2(a).\nl
%\begin{tabular}{ l l}
%1. $\forall x \tild p \ra \tild p(t/x)$ & Ax 1.3.2(h)\\
%2. $\tild \tild p(t/x) \ra \tild \forall x~\tild p$ & Eq, SC26  \\
%3. $p(t/x) \ra \tild\exists x \tild p$ & Ax 1.3.2(e) \\
%4. $p(t/x) \ra \exists x~p$ & Def 1.3.2(d) \\
%\end{tabular}\nl

Proof of PC3. \nl
1. $\forall x~p$\nl
2. $\forall x~p \ra p$\nl
3. $p$\nl
4. $p \ra \exists x~p$\nl
5. $\exists x ~p$\nl

\newpage
%page 15
Proof of PC4.\nl

1.  $\tild \forall x~p \lra \tild \forall x~p$\nl
2.  \qqu $\lra \tild \forall x \tild \tild p$\nl
3.  \qqu $\lra \exists x \tild p$\nl

Proof of PC8.\nl

1. $\forall x\forall y ~p$\nl
2. $\forall x\forall y ~p \ra \forall y~p$\nl
3. $\forall y~p$\nl
4. $\forall y~p \ra p$\nl
5. $p$\nl
6. $\forall x~p$\nl
7. $\forall y \forall x ~p$\nl

Proof of PC11. \nl
\begin{tabular}{ l l}
1. $\exists x~(p \up q)$ & Hyp \\
2. $p \up q$ & Th1.3.1 \\
3. $p$ & MP2, SC8 \\
4. $\exists x~p$ & MP3, PC2(b)\\
5. $q$ & MP2, Ax 1.3.2(f)\\
6. $\exists x ~q$ & MP5,PC2(b) \\
7. $\exists x~p \up \exists x~q$ & MP 4,6,Ax1.3.2(g)\\
\end{tabular}\nl

\textbf{Exercises.}\nl

Justify each step in the proofs of Theorems \nl

1. PC3\nl
2. PC4\nl
3. PC8\nl

Prove each of the following theorems:\nl

\begin{tabular}{l l}
4. PC5  &12. PC15\\
5. PC6 &13. PC16(a)\\
6. PC7 &14. PC16(b)\\
7. PC9 &15. PC16(c)\\
8. PC10 &16.  PC16(d)\\
9.  PC12 &17.  PC16(e)\\
10.  PC13 &18. PC16(f)\\
11. PC14 &\\
\end{tabular}\nl
%12. PC15 13. PC16(a) 14. PC16(b)
%15. PC16(c) 16.  PC16(d) 17.  PC16(e) 18. PC16(f)

\newpage
%page16
\textbf{1.4 Sets and Relations} \nl

\lhead{Section 1.4}
 During the last one hundred years the concepts of mathematics have been made more precise through the use of sets. In fact, every concept in mathematics can be expressed in terms of sets. We will consider a set to be a collection of objects. If $S$ is a set and $x$ is an object, then we will use $x \in S$ to denote the statement ``$x$ is an element of $S$" and use $x \notin S$ to denote  $\tild (x \in S)$. \nl

To specify a set, we must tell which objects it contains. There are two commonly used ways of doing this. The first method of specifying a set is to list its objects within braces. For example, the \textbf{set} containing the objects 1, 2, and 3 is denoted by $\{1, 2, 3 \}$.  Three dots are sometimes used to denote missing elements. For example,
$\{2, 4, 6, ...~ , 100\}$ denotes the set of all even positive integers less than or equal to 100, and $\{1, 3, 5, 7, ... \}$ denotes the set of all odd positive integers. \nl

The second method is better suited for specifying large sets without a pattern. If $p(x)$ is a formula, it is customary in mathematics to use $\{x \mid p(x) \}$ to denote the set to which $x$ belongs if and only if $x$ is an object such that $p(x)$ is true. There are some formulas for which no such set exists. For example, if $A= \{x \mid x \notin A\}$, then $x \in A \lra x \notin A$.  By SC37, $x \in A \lra x \notin A$ is a false statement. Therefore, $A = \{x \mid x \notin A\}$ is not a set. To avoid this difficulty we introduce the following axiom. \nl

\textbf{Axiom 1.4.1.} If $p(x)$ is a formula for which there is a set $B$ which contains all elements $x$ for which $p(x)$ is true, then there is a set (called the truth set of $p(x)$ and denoted by $\{ x \mid p(x)\}$ ) which contains exactly those elements for which $p(x)$ is true,
i. e. $y \in \{x \mid p(x)\} \lra p(y).$\nl

\textbf{Example 1.4.1.} Suppose $R$ denotes the set of all real numbers. Since ($x$ is an integer $\ra x \in R$),
we have by Axiom 1.4.1, $x\in \{x \mid x \text{~is~ an ~integer}\} \lra x$ is an integer.
In other words $\{x \mid x \text{~is~ an~ integer}\}$ denotes the set of all integers. \nl

\textbf{Definition 1.4.1.} If each of $x$ and $y$ is an object, then $x = y$ means $x$ and $y$ are names for the same object. \nl

\textbf{Definition 1.4.2.} The statement that a set $A$ is a \textit{subset} of a set $B$, written $A \subseteq B$, means that every element of $A$ is an element of $B$, i.e., $A \subseteq B \lra \forall x (x \in A \ra x \in B)$.\nl

\newpage


\textbf{Definition 1.4.3.} The statement that a set $A$ is a \textit{proper subset} of $B$, written $A\subset B$, means $A \subseteq  B$ and $A \neq B$.\nl

\textbf{Definition 1.4.4.} If each of $A$ and $B$ is a set, then: \nl
a) $A = B$ means $A \subseteq B$ and $B \subseteq A$, i.e., $A = B \lra  \forall x(x \in  A \lra x \in B)$.\nl
b) $A \neq B$ means $\tild (A = B)$.\nl

\textbf{Example 1.4.2.}

1.  $\{1, 2\} \subset \{1, 2, 3\}$\nl

2.  $\{1, 2, 2\}=\{1, 2\}$\nl

3.  $\{1,2,3, ... \} \subseteq \{1,2,3, ... \}$\nl

4.  $\{x \mid x > 2\} \subseteq \{x \mid x > 0\}$\nl


\textbf{Definition 1.4.5.} If each of $a, b, c$ and $d$ is an object, then: \nl
1) $\{\{a\},\{a,b \}\}$ is called an \textit{ordered pair} and denoted by $(a,b)$. If $(a,b)$ is an ordered pair, then $a$ is called its \textit{abscissa} and $b$ is called its \textit{ordinate.} \nl
2) $((a,b),c)$ is denoted by $(a,b,c)$.\nl
3) $((a,b,c),d)$ is denoted by $(a,b,c,d)$\nl

\textbf{Theorem 1.4.1.} If each of $(a,b)$ and $(c,d)$ is an ordered pair, then $(a,b)=(c,d)$ if and only if $a =c$ and $b = d$.\nl

\textbf{Definition 1.4.6.} If each of A and B is a set, then $\{(x,y)  \mid x \in A \text{~and~} y \in B\}$ is called the \textit{cartesian product} of A and B and denoted by $A \times  B$, i.e., $\forall x \forall y((x,y) \in A \times B \lra x \in A \up Y \in B)$.\nl

\textbf{Example 1.4.3.} If $A= \{1,2\}$ and $B= \{3,4,5\}$, then $A \times B =\{(1,3), (1,4), (1,5), (2,3), (2,4), (2,5)\}$. \nl

\textbf{Definition 1.4.7.} $\tild$ is said to be a \textit{binary relation} on a set $A$ if and only if $\tild \subseteq A \times A$. $\tild$ is said to be a binary relation if and only if $\tild$ is a binary relation on some set. We will write $x \tild y$ for $(x,y) \in \tild $. \nl

\textbf{Example 1.4.4.} $\{(1,2), (3,1), (2,3)\}$ and $\{(1,1), (2,2)\}$ are binary relations on $\{1,2,3\}$.\nl

\newpage
%page18
\textbf{Definition 1.4.8.} A binary relation $\tild$ on a set $A$ is said to be an
\textit{equivalence relation} if and only if the following formulas are true. \nl
 1) If $x \in A$, then $x \tild x$. (reflexive)  \nl
 2) If $x,y \in A$ and $x \tild y$, then $y \tild x$. (symmetric) \nl
 3) If $x, y, z \in A, x \tild y$ and $ y \tild z$, then $x \tild z$ . (transitive) \nl

\textbf{Theorem 1.4.2.} If $A$ is a set, then $=$ is an equivalence relation on $A$.\nl

\textbf{Exercises} \nl
1. If no set can contain itself, is $\{x \mid x \text{~is~ a~ set~}\}$ a set?  \nl
2. If $A$ is a set, which of the following are true? \nl
a) $\{1,2\} \subseteq \{1,2,3\}$ \nl
b) $\{1,2\} \subset \{ 1,2,2\}$ \nl
c) $\{(1,2)\} \subseteq \{1,2\} \times \{2,3\}$ \nl
d) $A=\{x \mid x \in A\}$ \nl
e) $\{\{1\}\} \subseteq \{1, 2, 3\}$ \nl
f) $\{\{a\} \} = (a,a)$ \nl
g) $A \in \{A\}$ \nl
h) $a \in A \ra \{a\} \subseteq A$ \nl

3. If $A= \{1,2,3\}$ and $B= \{0,1 \}$, find $A \times B$.  \nl

4. 	Which of the following are binary relations on {1,2,3}? \nl

a) $\{(1,2),(2,3)\}$ \nl
b) $\{(1,3),(3,4),(2,1)\}$  \nl
c) $\{(x,y) \mid x,y \in \{1,2,3\} \text{~and~} x=y\} $ \nl
d) $\{(x,y) \mid  x \text{~ and~} y \text{~are ~numbers~ and ~}x < y\}$ \nl

5. 	Which of the following are equivalence relations on {1,2,3}? \nl

a) $\{(1,1) (2,2)\}$\nl
b) $\{(1,1), (2,2), (3,3)\}$\nl
c) $\{(1,1), (2,2), (3,3), (1,3)\}$\nl
d) $\{(1,1), (2,2), (3,3), (1,3), (3,1)\}$\nl
e) $\{(1,1), (2,2), (3,3), (1,2), (2,3). (2,1), (3,2)\}$\nl

6.  If ``$p(x)$" stands for the formula ``$x$ likes calculus," express each of the following statements in English.\nl

\begin{tabular}{l l}
a) $p$(He) & c) $p$(John)\\
b) $p$(She)& d) $p$(Everybody)\\
\end{tabular}\nl

7. If $S = \left\lbrace n \mid n \text{~is ~an ~integer ~greater ~than ~2} \right\rbrace,$ express each of the following statements in English.\nl

\begin{tabular}{l l}
a) $1 \in S$ & c) $k + 1 \in S$\\
b) $k \in S$ & d) $h + k \in S$\\
\end{tabular}\\
\newpage
%page19
\textbf{1.5 Functions}\nl

\lhead{Section 1.5}
 Of all the sets we will study probably none is more important than those we call functions.\nl
 
\textbf{Definition 1.5.1.} $f$ is said to be a \textit{function} if and only if $f$ is a set of ordered pairs such that no two pairs of $f$ have the same first term. \nl

\textbf{Example 1.5.1.} $\{(1,2)\}, \{(1,2),(5,13),(3, 8)\}$, and $\{((1,1),2), ((1,2),3), ((2,1),3), ((2,2),4)\}$ are functions. $\{(1,2), (3,5), (1,-2)\}$ and $\{ 1, 2\}$ are not functions. \nl

\textbf{Definition 1.5.2.} Suppose $f$ is a function. $\{x \mid \text{~there ~is ~a ~y~} \text{~such ~that~} (x,y) \in f\}$ is called the \textit{domain} of $f$ and denoted by $D_{f}$.  $\{y \mid \text {~there ~is ~an~} x \text{~such ~that~} (x,y) \in f\}$ is called the \textit{range} of $f$ and denoted by $R_{f}$. \nl

\textbf{Example 1.5.2.} If f= $\{(1,2), (-3,1), (2,4), (3,6)\}$, then $D_{f}= \{1, -3, 2, 3\}$ and $R_{f}= \{2, 1, 4, 6\}$.  \nl

\textbf{Definition 1.5.3.} If $f$ is a function and $x \in D_{f}$, then the only element $y$ such that $(x,y) \in f$ is called ``$f$ of $x$" and denoted by $f( x)$, i.e., $y = f(x) \lra (x,y) \in f.$ \nl

\textbf{Example 1.5.3.} $f =\{(1,2), (-1,3), (2,-4) \}$, then $f( 1) =2, f( -1)=3$, and $f( 2)=-4$.\nl

\textbf{Definition 1.5.4.} Suppose $A$ and $B$ are sets. The statement that $f$ is a function from $A~into~B$, written $f: A \ra B$, means $f$ is a function with domain $A$ and range a subset of $B$. The statement that $f$ is func­tion from $A~onto~B$, written $f:A \dashv B$, means $f:A \ra B$ and $B$ is the range of $f$.\nl

\textbf{Example 1.5.4.} If $f= \{(1,2), (3,-1), (-1,1)\}$, then $f: \{1, 3, -1\} \ra \{ -1, 0, 1, 2\}$ and $f: \{1, 3, -1 \} \ra \{-1, 1, 2\}$  are true and $f : \{1, 0, -1, 3 \} \ra \{-1, 1, 2\}$ and $f : \{1, 3, -1\} \onto \{-1, 0,1, 2\}$ are false. \nl

\textbf{Definition 1.5.5.} $\bigstar$ is said to be a \textit{binary operation} over a set $A$ if and only if $\bigstar$ is a function and the domain of $\bigstar$ contains $A \times A$. We will write $x \bigstar y$ for $\bigstar(x,y)$. \nl

\textbf{Example 1.5.4.5} $\{((1,1)2),((1,2),3),((2,1),3),((2,2),4)\}$ is a binary operation over $\{1, 2\}$ .

\newpage
%page20

\textbf{Exercises}\nl

1. Which of the following are functions? \nl
\indent a) $\{ 1, 2 \}$\nl
\indent b) $\{(-1, 2), (-3, 4), (2, -1)\}$\nl
\indent c) $\{(3, -1),(1, -4),(3,6) \}$\nl
\indent d) $\{((1,2),3),(4,5)\}$\nl
\indent e) $\{1, -2, (1,3)   \}$\nl

2. Suppose $f= \{(-1,2), (0,-1), (1,3), (-2,-1)\}$. \nl
a) Find $D_{f}$. \nl
b) Find $R_{f}$. \nl
c) Find $f(-1), f(0), f(1)$, and $f(-2)$ . \nl
d) List the elements in $\{x \mid f(x)= -1\}$ and $\{x \mid f(x+2)=3\}$\nl

3. Suppose $f = \{(x,y) \mid x \text{~is ~a number~ and~} y=x^{2}+2x-1\}$ . \nl
a) Find $D_{f}$. \nl
b) Find $R_{f}$. (Hint: $f(x)=(x+1)^{2} -2$) \nl
c) Show that if $x$ and $h$ are numbers, then $[f(x+h)-f(x)]/h=2x + 2 + h$ \nl

4. 	Which of the following are binary operations over {1, 2, 3} ?\nl
a) $\{((1,1),2), ((1,2),3), ((1,3),4), ((2,1),3), ((2,2),4), ((2,3),5),$
$((3,1),4),((3,2) 5)((2,1),3)\}$\nl
b) $\{((1,2),3), ((2,1),3)\}$\nl
c) $\{((x,y),z) \mid ~x~ \text{and} ~y~ \text{~are ~numbers~ and~} z = x \cdot y\}$ \nl

5 .Suppose $A=\{1,2,3\}$ and $B=\{4,5,6\}$. For each of the following denitions of $f$ determine whether $f :A \ra B$ or $f :A \onto B$ is true.  \nl
a) $f= \{(1,4), (2,6), (3,5)\}$\nl
b) $f=\{ (1,5), (2,6), (3,5)\}$ \nl
c) $f =\{(1,4), (2,3), (3,5)\}$\nl
d) $f= \{(x,y) \mid x \in \{1,2, 3\} \text{~and~} y = x + 3\}$ \nl

6. If $\bigstar$ is a binary operation over $\{1,2,3\}$ defined by the following table, find $a$ for each of the following. \nl


%6.  If ``$p(x)$" stands for the formula ``$x$ likes calculus," express each of the following statements in English.\\

%\begin{tabular}{l r}
%a) $p$(He & c) $p$(John)\\
%b) $p$(She) &d) $p$(Everybody)\\
%\end{tabular}\\

%7. If $S = \left\lbrace n \mid n \textbf{~is ~an ~integer ~greater ~than ~2} \right\rbrace,$ express each of the following statements in English.

\begin{tabular}{| c | c | c | c |}
\hline
$\bigstar$ & 1 & 2 & 3\\ \hline
1 & 2 & 3 & 1 \\ \hline
2 & 3 & 1 & 2 \\ \hline
3 & 1 & 2 & 3 \\ \hline
\end{tabular} \nl

a) $2\bigstar 3 = a$\nl
b) $a\bigstar 2 = 2$\nl
c) $\forall x \in  \{1,2,3\}, x\bigstar a = x $\nl
d) $1\bigstar a = 3 $\nl
e) $2\bigstar a = 3 $\nl
f) $3\bigstar a = 3 $\nl

\newpage
%page21
\textbf{1.6 Algebra of Sets}\nl

\lhead{Section 1.6}
The algebra of sets is very closely related to the statement calculus. In fact, by an appropriate substitution of symbols in each theorem in the statement calculus, we can obtain a theorem about sets. \nl

\textbf{Definition 1.6.1.} If each of $A$ and $B$ is a set, then $\{x \mid  x \in A \up x \in B\}$ is called the \textit{intersection} of $A$ and $B$ and is denoted by $A \cap B$, i.e., $\forall x (x \in A \cap B) \lra x \in A \up x \in B)$ \nl

\textbf{Example 1.6.1.} $\{1, 2, 3\} \cap \{2, 3, 4\} =\{2, 3\}$. \nl

\textbf{Definition 1.6.2.} If each of $A$ and $B$ is a set, then $\{x \mid x \in A \dn  x \in B \}$ is called the \textit{union} of $A$ and $B$ and is denoted by $A \cup B$, i.e.   $\forall x (x \in (A \cup B) \lra x \in A \dn x \in B)$ \nl

\textbf{Example 1.6.2.} $\{ 1, 2, 3\} \cup  \{2, 3, 4\} = \{ 1,2,3,4\}$.\nl

\textbf{Definition 1.6.3.} If each of $A$ and $B$ is a set, then ·$\{x \mid x \in A \up x \notin B\}$ is called $A~minus~B$ and denoted by $A - B$, i.e., $\forall (x(x \in (A-B) \lra x \in A \up  x \notin B)$. \nl

\textbf{Example 1.6.3.} $\{1, 2, 3\}-\{3, 4, 5\} = \{1, 2\}$. \nl

\textbf{Definition 1.6.4.} The set which contains no elements is called the \textit{empty set} and denoted by $\emptyset$, i.e., $\forall x(x \notin \emptyset)$.\nl

\textbf{Definition 1.6.5.} The nonempty set of all elements under consideration is called the \textit{universal set} and sometimes denoted by \textbf{U}, i. e., $\forall x(x \in \textbf{U}).$ \nl

\textbf{Definition 1.6.6.} If \textbf{U} is a universal set and $A \subseteq \textbf{U}$, then $\textbf{U} - A$ is called the \textit{complement} of $A$ and denoted by $A^{'}$ . \nl

\textbf{Definition 1.6.7.} Suppose $G$ is a collection of sets. $\{x \mid \text{~there ~is ~a ~}g \text{~in~} G \text{~such ~that~} x \in g\}$ is called the \textit{union} of $G$ and denoted by $\cup G$. \nl
$\{x \mid \text {~for~ each~} ~g~ \text{~in~ }G, x \in g\}$ is called the \textit{intersection} of $G$ and denoted by
$\cap G$. \nl

\textbf{Definition 1.6.8.} If $A$ is a set, then $\{B \mid B \subseteq A\}$ is called the \textit{power set} of
$A$ and is denoted by $\power A$, i.e., $ \forall B(B \in \power A \lra B \subseteq A)$. \nl

\textbf{Example 1.6.4}. $ \power \{ 1, 2, 3\} = \{ \emptyset, \{ 1 \}, \{2\}, \{3\}, \{ 1, 2\}, \{1, 3\} , \{ 2, 3\}, \{1, 2, 3\} \}.$ \nl

\newpage
% p 22
\textbf{Section 1.6} \nl

If $p$ is a statement containing finitely many symbols and each symbol in $p$ denotes a subset of $\univ$ or is one of $=, \cup , \cap ,~ ', \emptyset , \text{~or~ } \univ $, then the truth or falsity of $p$ can be determined from S1 through S6. Any system having these properties is called a Boolean Algebra. \nl

\textbf{Theorem 1.6.1.} If $\univ$ is a universal set and each of $A, B,$ and $C$ are subsets of $\univ$, then: \nl
\begin{tabular}{l l}
S1.  a) $A \cup B = B \cup A$ &  b) $A \cap B=B \cap A$ (Commutative Laws)\\
S2.  a) $A \cup (B\cup C)=(A \cup B)\cup C$ & b) $A\cap (B\cap C)= (A\cap B)\cap C$ $(*)$ \\
%( Associative Laws) \\
S3. a) $A\cup (B\cap C)=(A\cup B)\cap (A\cup C)$ & b) $A\cap (B\cap C)=(A\cap B)\cup (A\cap C)$ $(*)$\\
%(Distributive Laws) \\
S4. a) $A \cup \emptyset = A$ & b) $A \cap \univ = A$
(Identity elements) \\
S5. a) $A \cup A^{'} = \univ$ & b) $A \cap A^{'} = \emptyset $
(Complement)\\
S6. $\univ \neq \emptyset $&\\
\end{tabular}\nl
$(*)$ Note that S2 is called the Associative Laws and S3 is called the Distributive Laws.\nl
Proof of S3(a). Suppose $\univ$ is a universal set and $A, B$, and $C$ are subsets of $\univ$.\nl

\begin{alignat}{2}
A \cup (B \cap C) &= \{x \mid x \in A \dn x \in B \cap C\}\\
&= \{x \mid x \in A \dn (x\in B \up x \in C\}\\
&= \{x \mid (x \in A \dn x \in B) \up (x \in A \dn x \in C)\}\\
&= \{x \mid x\in A \cup B \up x \in A \cup C\}\\
&= (A \cup B) \cap (A \cup C)
\notag
\end{alignat}

Proof of S4(b). Suppose $\univ$ is a universal set and $A \subseteq \univ$  \nl
$\forall x(x \in A \ra x \in \univ )$ \nl
$\forall x(x \in A \up x \in \univ  \lra x \in A)$ \nl
$ A \cap \univ = A $ \nl

Proof of S5(a). Suppose $\univ$ is a universal set and $A \subseteq \univ $. \nl
$\forall x (x \in A \dn x \in A^{'}) $\nl
$\forall x(x \in A \dn A{'})) $ \nl
$\forall x(x \in \univ)$ \nl
Therefore, \nl
$A \cup A^{'} = \univ$ \nl

\newpage
%page 23

%=$\{x \mid x \in A \dn (x \in \univ \up x \notin A) \}$ \nl
%=$\{ x \mid (x \in A \dn x \in \univ ) \up (x \in A \dn x \notin A)\}$ \nl
%=$\{x \mid x \in A \cup \univ \up t\}$ \nl
%=$\{x \mid x \in \univ \up t\}$ \nl
%=$\{x \mid x \in \univ \}$ \nl
%=$\univ$ \nl

\textbf{Theorem 1.6.2.} If $\univ$ is a universal set and $A, B,$ and $C$ are subsets of $\univ$, then: \nl

S7. a) $A \cup A=A$ and b) $A \cap A =A $\nl
S8. a) $A \cup \univ = \univ$ and b) $A \cap \emptyset = \emptyset$\nl
S9. a) $A \cup (A \cap B) = A$ and b) $A \cap (A \cup B)=A$ \nl
S10 a) $A \cup (B \cap A^{'})= A\cup B$ and b) $A \cap (B \cup A^{'}) =A \cap B$ \nl
Proofs of S7(a) and S8(a). Suppose $\univ$ is a universal set and $A \subseteq \univ$.

Proof of S7(a)\nl
$A = A \cup \emptyset$\nl
$ = A\cup (A\cap A^{'})$  \nl
$ = (A\cup A) \cap (A \cup A^{'})$ \nl
$ = (A \cup A) \cap \univ $ \nl
$ = A \cup A$\nl

 Proof of S8(a)\nl
 $\univ = A \cup A^{'}  = (A \cup A) \cup A^{'}  = A \cup (A \cup A^{'}) =A \cup \univ $\nl

\textbf{Theorem 1.6.3.} If $\univ$ is a universal set and $A$ and $B$ are subsets of $\univ$ such that $A \cup  B = \univ$ and $A \cap B = \emptyset$ , then $A^{'} = B$. \nl
Proof. Suppose $\univ$ is a universal set, $A,B \subseteq \univ, A \cup B= \univ$, and $A \cap B = \emptyset $. $A^{'} = \univ \cap A^{'} = (A\cup B)\cap A^{'} =(A\cap A^{'}) \cup (B\cap A^{'})$
$=\emptyset \cup(B \cap A^{'}) =(A \cap B) \cup(A^{'} \cap B) = (A\cup A^{'})\cap B = \univ \cap B = B $\nl

\textbf{Theorem 1.6.4.} If $\univ$ is a universal set and $A$ and $B$ are subsets of $\univ$, then:\nl
 S11. $(A^{'})^{'} = A $\nl
 S12. a) $U^{'}= \emptyset$ \qqu b) $\emptyset ^{'} = \univ $ \nl
 S13. a) $(A \cup B)^{'} = A^{'}\cap B^{'}$  \qqu b) $(A\cap B)^{'} = A^{'} \cup B^{'}$ \nl

 Proof of S11. Suppose $\univ$ is a universal set and $A \subseteq \univ$. Since $A \cup A^{'} = \univ$ and $A \cap A^{'} = \emptyset $, we have $A^{'} \cup A = \univ$ and $A^{'} \cap A= \ns$. By Theorem 1.6.3, $(A^{'})^{'} = A$. \nl

 Proof of S13(a). Suppose $\univ$ is a universal set and $A,B \subseteq  \univ$. $(A\cup B) \cap (A^{'} \cap B^{'}) =[A\cap (A^{'} \cap B^{'})] \cup [B \cap  (A^{'} \cap B^{'})] $

\newpage
%page 24
$\qqu \qqu = [(A \cap A^{'}) \cap B^{'}] \cup [(B \cap B^{'}) \cap A^{'}]$ \\
$\qqu \qqu =(\ns \cap B^{'}) \cup (\ns \cap A^{'})$\\
$\qqu \qqu = \ns \up \ns = \ns$ \\

$(A \cup B) \cup (A^{'}\cap B^{'}) = [(A \cup B) \cup A^{'}] \cap [(A \cup B) \cup B^{'}]$\\
$\qqu \qqu \quad = [(A \cup A^{'}) \cup B] \cap  [A \cup (B\cup B{'})]$\\
$\qqu \qqu \quad =(\univ \cup B)\cap (A \cup \univ)$ \\
$\qqu \qqu \quad =\univ \cap \univ = \univ$ \\

By Theorem 1.4.3, $(A \cup B)^{'}=A^{'} \cap B^{'}$\nl

\textbf{Theorem 1.6.5.} If $\univ$ is a universal set and $A, B$ and $C$ are subsets of $\univ$, then: \nl
\begin{tabular}{l l}
S14. a) $A\subseteq B \lra A\cup B = B$&  b) $A \subseteq B \lra A \cap B = A$ \\
S15. a) $A \subseteq B \lra A^{'} \cup B= \univ$ \qqu & b) $A \subseteq B \lra A\cup B^{'} = \ns$\\
S16. a) $A\cap B \subseteq A$  & b) $A \subseteq A\cup B$ \\
S17. a) $A\subseteq B \ra (A \cup C)\subseteq (B \cup C)$ & b) $A \subseteq B \ra (A\cap C) \subseteq (B \cap C)$\\
\end{tabular}
S18. $A\subseteq B \lra B^{'} \subseteq A^{'}$ \\
S19. $\ns \subseteq A$ \nl

Proof of S14(a). Suppose $\univ$ is a universal set and $A,B \subseteq \univ$. \nl

\begin{tabular}{r l}
$A \subseteq B$ &$\lra \forall x [x \in A \ra x \in B ]$ \\
&$\lra \forall x [x \in A \dn x \in B \lra x \in B]$ \\
&$\lra \forall x [ x \in (A \cup B) \lra x \in B]$\\
&$\lra A \cup B = B$ \\
\end{tabular}\\

Proof of S14(b).  Suppose $\univ$ is a universal set and $A, B \subseteq \univ$.\nl
\begin{tabular}{ r l}
$A\subseteq B$ & $\lra \forall x[x\in A \ra x\in B]$ \\
&$\lra \forall x [x\in A \up x\in B \lra x\in A]$ \\
&$\lra \forall x[x\in A\cap B \lra x\in A] $    \\
&$\lra A\cap B = A $ \\
\end{tabular}\\

By Theorem 1.2.1, $ a \subseteq b \lra A \cap B = A$.\nl

$\qqu$ Proof of S15(a).  Suppose $\univ$ is a universal set and $A, B \subseteq \univ$.\nl
\begin{tabular}{ r l}
$A \subseteq B$ & $\lra \forall x[x\in A \ra x\in B]$\\
&$\lra \forall x[x\in A{'} \dn x\in B \ra x \in \univ] $\\
&$\lra \forall x[x\in A^{'}\cup B \lra x\in \univ]      $\\
&$ \lra A^{'}\cup B = \univ $\\
\end{tabular}\\

By Theorem 1.2.1, $ A \subseteq B \lra A^{'}\cup B = \univ$
\newpage
%page25

\textbf{Exercises}
1. If $\univ = \{0, 1, 2,3,4,5,6,7, 8\}, A= \{1, 2,3,4 \}, B= \{3, 4,5, 6\},\text{~ and~} C= \{2, 3, 6, 7\}$, specify the elements in each of the following: \nl

a) $A-B$\nl
b) $A \cap B$\nl
c) $(A \cup B)^{'}$\nl
d) $A^{'} \cap B^{'}$\nl
e) $B \cup C$\nl
f) $A \cap C$\nl
g) $A \cap (B \cup C)$\nl
h) $(A \cap B) \cup (A\cap C)$\nl
i) $A\cup (A \cap C)$\nl

2. Prove Sl(a).\nl
3. Prove S2(a).\nl
4. Prove S3(b).\nl
5. Prove S4(a).\nl
6. Prove S5(b).\nl
7. Prove S6(b).\nl
8. Prove S7(b).\nl
9. Prove S8(b).\nl
10. Prove S9(a).\nl
11. Prove S10(a).\nl
12. Prove S12(a).\nl
13. Prove S13(b). \nl
14. Prove S15(b).\nl
15. Prove S16(a). \nl

\newpage
%page26
\textbf{Chapter 2} \nl

\lhead{Section 2.1}
\textbf{Fields}\nl

In the physical sciences knowledge can be expanded by laboratory
experiments. Since mathematics is a creation of the mind, the only way
to extend the body of mathematical knowledge is by the process of reasoning. \nl

One possible way of developing the number system would be to assume every statement that is intuitively obvious about numbers and proceed by the process of reasoning to arrive at statements which are not obvious. This process has two definite disadvantages. (1) The statements which are obvious to one person are not necessarily obvious to another.  (2) Eventually someone is going to prove a statement false which is obviously true. The reason for this is that intuition is not an accurate measure of the truth or falsity of a statement.
This suggests another approach to the number system. By a careful choice of assumptions we can minimize the above mentioned drawbacks. The set of axioms we will use is that of a complete ordered field. These axioms are intuitively obvious to most people and no one has ever proven one of them false as a consequence of the others. This set of axioms does have the disadvantage of having to prove statements true which are intuitively obvious. But this objection is insignificant compared to the objections to the first method discussed.
Fields are important since the real number system, the complex number system, and vector spaces are all defined in terms of fields. \nl

\textbf{2.1 Properties of a Field} \nl

\textbf{Definition 2.1.1.} $(F, +,\cdot)$ is called a \textit{field} if and only if $F$ is a nonempty set, $+$ is a binary operation over $F$ (called addition), and $\cdot$ a binary operation over $F$ (called multiplication), such that the following statements are true. \nl

A1. If $x,y \in F$, then $(x+y) \in F$. (Closure)  \nl
A2. If $x,y \in F$, then $x + y = y + x$. (Commutative Law) \nl
A3. If $x,y,z \in F$, then $(x+y)+z=x+(y+z)$. (Associative Law) \nl
 \newpage
% page27

A4. There is an element $0$ in $F$ (called the identity for $+$) such that if $x \in F$, then $x + 0 = x.$ \nl
A5. If $x \in F$ , then there is an element $-x$ in $F$ (called the additive inverse
of $x$) such that $x+(-x)=0.$ \nl

M1. If $x,y \in F$, then $x\cdot y \in F$. (Closure) \nl
M2. If $x,y \in F$ ,then $x \cdot y = y \cdot x$. (Commutative Law) \nl
M3. If $x,y,z \in F$ , then $(x \cdot y)\cdot z = x \cdot(y\cdot z)$. (Associative Law)\nl
M4. There is an element $1$ in $(F-\{ 0\})$, ( called the identity for $\cdot$ ), such that if $x \in F$, then $x\cdot 1 =x $.\nl
M5.  If $x \in F$ and $x \neq 0$, then there is an element $x^{-1}$  in $F$ (called the multiplicative inverse of $x$) such that $x\cdot x^{-1} =1.$ \nl

D.	If $x,y,z \in F$, then $x\cdot(y + z) =(x\cdot y)+(x\cdot z)$. (Distributive Law) \nl

\textbf{Theorem 2.1.1.} If $(F,+ ,\cdot)$ is a field, then the following statements are true. \nl
a) If $x \in F$, then $0 + x = x.$\nl
b) If $x \in F$ ,then $(-x) + x = 0$. \nl
c) If $x \in F$, then $1\cdot x=x.$\nl
d) If $x \in F$ and $x \neq 0$, then $ x^{-1}\cdot x = 1.$ \nl
e) If $x,y,z \in F$, then $(x+y)\cdot z=(x\cdot z) + (y\cdot z).$\nl

\textbf{Theorem 2.1.2.} If $(F,+,\cdot)$ is a field, $x,y,z \in F$, and either $x + z = y + z$ or $z + x = z + y$, then $x = y.$ \nl
Proof. Suppose $(F,+,\cdot)$ is a field, $x,y,z \in F$, and $x + z = y + z$. \nl
\begin{tabular}{r l l}
$(x+z)+(-z)$&$=(y+z)+(-z)$ & Def of =, A5 \\
$x+(z+(-z))$&$=y+(z+(-z))$ & A3 \\
$x+0$&$= y+0$ &A5 \\
$x$& $= y$ & A4 \\
\end{tabular}\\
By a similar argument, we see that if $z + x = z + y$, then $x = y.$\nl

\textbf{Theorem 2.1.3.} If $(F,+,\cdot)$ is a field, $x,y,z \in F, z\neq  0$, and either $x\cdot z=y\cdot z$
or $z\cdot x = z\cdot y$ , then $x = y.$ \nl

\textbf{Theorem 2.1.4.} If $(F,+,\cdot)$ is a field, $x, y \in F$, and $x + y = x$, then $y=0.$ \nl

\textbf{Theorem 2.1.5.} If $(F,+,\cdot)$ is a field, $x,y \in F, x \neq 0$, and $x\cdot y = x$, then $y = 1.$ \nl

\textbf{Theorem 2.1.6.} If $(F,+,\cdot)$ is a field, $x,y \in F$, and $x + y = 0$, then $y = -x$.\nl
 Proof. Suppose $(F,+,\cdot)$ is a field, $x,y \in F$, and $x + y = 0$.
 \begin{tabular}{r l l}
 $x + (-x)$ &$= 0$ &A5  \\
 $x + y $ & $= x + (-x)$ & Def of = \\
 $y$ &$= -x$ & Th. 2.1.2\\
\end{tabular}
\newpage
%page28
\textbf{Theorem 2.1.7.} If $(F,+,\cdot)$ is a field and $x\in F$, then $x\cdot 0=0$ and $0\cdot x = 0$.\nl
 Proof. Suppose $(F,+,\cdot)$ is a field and $x \in F.$ \nl

\begin{tabular}{r l l l}
$0 + 0$&= 0 &&A4\\
$x\cdot (0 + 0)$&$=x\cdot 0$ &&Def of=\\
$x\cdot 0 + x\cdot 0$&$=x\cdot 0 + 0$ &&D,A4\\
$x\cdot 0$&$=0$ &&Th 2.1.2\\
$0\cdot x $&$= 0$ && M2\\
\end{tabular}\nl


\textbf{Theorem 2.1.8.} If $(F,+,\cdot)$ is a field, $x,y \in F$, and $x\cdot y = 1$, then $y = x^{-1}$ \nl
\textbf{Exercises}


1. Prove Theorem 2.1.1.

2. Prove Theorem 2.1.3.

3. Prove Theorem 2.1.4.

4. Prove Theorem 2.1.5.

5. Prove Theorem 2.1.8.


6. Prove that $(F,+,\cdot)$ is a field if $F = \{a,b\},a \neq b$,and $+$ and $\cdot$ are defined
by the following tables.\nl
\begin{center}
\begin{tabular}{|c|c|c|}\hline
$+$&$a$&b\\\hline
$a$&$a$&$b$\\\hline
$b$&$b$&$a$\\\hline
\end{tabular}
$\qqu \qqu$
\begin{tabular}{|c|c|c|}\hline
$\cdot$&$a$&b\\\hline
$a$&$a$&$a$\\\hline
$b$&$a$&$b$\\\hline
\end{tabular}\nl
\end{center}


\textbf{Theorem 2.1.9.} If $(F,+,\cdot )$ is a field and $x \in F$, then $-(-x) = x$.\nl
 Proof. Suppose $(F,+,\cdot )$ is a field and $x \in F$.
\begin{tabular}{r l l}
$(-x) + x$& $= 0$ & Th 2.1.1 \\
$x$& $= -(-x)$ & Th 2.1.6 \\
\end{tabular} \nl

\textbf{Theorem 2.1.10.} If $(F,+,\cdot)$ is a field, $x \in F$, and $x \neq 0$, then $(x^{-1})^{-1} = x.$ \nl

\textbf{Theorem 2.1.11.} If $(F,+,\cdot )$ is a field, and $x,y \in F$, then $x\cdot (-y) = -(x\cdot y).$\nl Proof. Suppose $(F,+,\cdot )$ is a field and $x,y \in F.$\nl
\begin{tabular}{r l l l}
$x\cdot y + x\cdot (-y)$&$=x\cdot (y + (-y))$&&D \\
&$=x\cdot 0$ &&A5\\
&$=0$ &&Th 2.1.7 \\
$x\cdot (-y)$&$=-(x\cdot y)$ &&Th 2.1.6
\end{tabular}\nl


\textbf{Theorem 2.1.12.} If $(F,+,\cdot )$ is a field and $x,y \in F$, then $(-x)\cdot y = -(x\cdot y)$ and $(-x)\cdot (-y) = x\cdot y.$

\newpage
%page29

\textbf{Theorem 2.1.13.} If $(F,+,\cdot )$ is a field and $x \in F$, then $(-1)\cdot x = -x.$\nl

\textbf{Theorem 2.1.14.} If $(F,+,\cdot )$ is a field, then $0 = -0.$\nl
Proof. Suppose $(F,+,\cdot)$ is a field.

\begin{tabular}{r l l l}
$0 + 0$&$= 0$ &&A4\\
$0$ &$=-0$&& Th 2.1.6\\
\end{tabular}

\textbf{Theorem 2.1.15.} If $(F,+,\cdot)$ is a field, then $1 = 1^{-1}.$\nl

\textbf{Definition 2.1.2.} If $(F,+,\cdot )$ is a field and $x,y \in F$, let $x - y$ denote $x + (-y)$, and if $y \neq 0$, let $x\div y$ denote $x\cdot y^{-1}.$\nl
Note: $x\cdot y^{-1}$ is also commonly denoted by $x/y$ and $\dfrac{x}{y}.$\nl

\textbf{Example 2.1.1.} If $(F,+,\cdot)$ is a field and $a,b,c\in F$, then $c\cdot(a-b)=c\cdot a - c\cdot b.$\nl
Proof. Suppose $(F,+,\cdot)$ is a field and $a,b,c \in F$.
\begin{tabular}{r l l l}
$c\cdot (a-b)$ &$=c\cdot (a+(-b))$ &&Def 2.1.2 \\
&$=c\cdot a + c\cdot (-b)$&& D \\
&$=c\cdot a + [-(c\cdot b)]$ &&Th 2.1.11 \\
&$=c\cdot a - c\cdot b$ &&Def 2.1.2 \\
\end{tabular}\nl


\textbf{Exercises} Assuming $(F,+,\cdot)$ is a field, $a,b,x,y,z \in F$, prove each of the following. \nl
1. If $a \neq 0$, then $-a \neq 0$.\nl
2. If $x + y = z$, then $y = z - x.$\nl
3. $a -0 = a.$\nl
4. $0 - a = -a.$\nl
5. $a \div 1 = a.$ \nl
6. If $a \neq 0$, then $0 \div a =0.$\nl
7. If $a \neq 0$, then $a \div a = 1.$\nl
8. If $x\cdot y = z$, and $x \neq 0$, then $y =z\div x.$ \nl
9. If $a\div b = 1$,then $a = b.$\nl
10. If $a = b, a \neq 0$, and $b \neq 0$, then $1\div a = 1 \div b.$\nl
11. If $a = -b,$ then $-a = b.$\nl
12. If $a = 1\div b$, and $a\neq 0$, then $b = 1\div a.$\nl
13. Theorem 2.1.10. \nl
14. Theorem 2.1.12. \nl
15. Theorem 2.1.13. \nl
16. Theorem 2.1.15. \nl
\newpage
%page30
\textbf{2.2 Fractions }\nl

\lhead{Section 2.2}
\textbf{Theorem 2.2.1.} If $(F,+,\cdot)$ is a field and $x,y \in F$, then $-(x+y) = -x - y$.\nl
Proof. Suppose $(F,+,\cdot )$ is a field and $x,y\in F.$\nl
\begin{tabular}{r l l l}
$(x+y)+(-x-y)$& $=(x+y)+(-x+(-y))$ &&Def 2.1.2 \\
&$=(x+y)+((-y)+(-x))$ &&A2 \\
&$=[(x+y)+(-y)] +(-x)$ &&A3\\
&$=[x+(y+(-y))] +(-x)$ &&A3\\
&$=(x+0)+(-x)$&& A5\\
&$=x+(-x)$&& A4\\
&$=0$ && A5\\
$-x - y$& $=-(x+y)$&& Th 2.1.6 \\
\end{tabular}\nl

\textbf{Theorem 2.2.2.} If $(F,+,\cdot)$ is a field, $x,y \in F, x \neq  0$, and $y \neq 0$ , then
$(x\cdot y)^{-1} =x^{-1}\cdot y^{-1}.$\nl

\textbf{Theorem 2.2.3.} If $(F,+,\cdot)$ is a field, $u,v,x,y\in F, u \neq 0$, and $v \neq 0$, then $(x\div u)\cdot (y \div v) =(x\cdot y) \div  (u\cdot v).$\nl
Proof. Suppose $(F,+,\cdot)$ is a field, $u,v,x,y\in F, u \neq 0$, and $v \neq 0$.
\begin{tabular}{r l l l}
$(x\div u)\cdot(y\div v)$ & $=(x\cdot u^{-1})\cdot(y\cdot v^{-1})$ &&Def 2.1.2 \\
&$=[(x\cdot u^{-1})\cdot y]\cdot v^{-1}$ &&M3\\
&$=[x\cdot (u^{-1} y)]\cdot v^{-1}$ &&M3 \\
&$=[x\cdot (y\cdot u^{-1})]\cdot v^{-1}$ &&M2\\
&$=[(x\cdot y)\cdot u^{-1}]\cdot v^{-1}$ &&M3 \\
&$=(x\cdot y)\cdot (u^{-1}\cdot v^{-1})$ &&M3 \\
&$=(x\cdot y)\cdot (u\cdot v)^{-1}$ &&Th 2.2.2 \\
&$=(x\cdot y)\div (u\cdot v)$ &&Def 2.1.2 \\
\end{tabular}\nl

\textbf{Theorem 2.2.4.} If $(F,+,\cdot)$ is a field, $x,y,z\in F, y \neq 0$, and $z \neq 0$, then $x \div y =(x\cdot z) \div (y\cdot z)$ and $x \div y =(z \cdot x) \div (z \cdot y).$ \nl

\textbf{Theorem 2.2.5.} If $(F,+,\cdot)$ is a field, $x,y \in F, x \neq 0$, and $y \neq 0$, then $(x \div y)^{-1}= y\div x.$\nl

\textbf{Theorem 2.2.6.} If $(F,+,\cdot)$ is a field, $u \in F$, and $v,x,y \in (F - \{0\})$, then
$(u \div v) \div (x \div y) $
$= (u \cdot y) \div  (v \cdot x).$ \nl

\textbf{Theorem 2.2.7.} If $(F,+,\cdot)$ is a field, $u,v,x,y \in F, u \neq 0,$ and $v \neq 0$, then $(x\div u) + (y\div v) =(x \cdot v +u\cdot y)\div (u\cdot v)$ .\nl

\newpage
%page31
\textbf{Exercises} \nl

Assuming $(F,+,\cdot)$ is a field and $a,b,c,d \in F$, prove each of the following. \nl

1. If $b \neq 0$, and $c\neq 0$ , then $a \div (b\div c) =(a\cdot c)\div b$. \nl

2. If $b\neq 0$~ and $c \neq 0$, then $(a \div b)\div c =a\div (b\cdot c)$.\nl

3. If $b\neq 0$, and $d\neq 0,$ then $(a\div b)=(c\div d)$ if and only if $a\cdot d=b\cdot c.$ \nl

4. If $\neq 0$, and $b \neq 0$, then $(1\div a)+(1\div b)=(b + a)\div (a\cdot b)$. \nl

5. If $a \neq 0$, and $b \neq 0$, then $1\div (a^{-1}\cdot b^{-1})= a\cdot b.$ \nl

6. If  $a \neq 0$ and $b \neq 0$, then $(a + b)\div (a^{-1} + b^{-1}) = a\cdot b$. \nl

7. If  $a \neq 0$ and $b \neq 0$, then $[(a^{-1} + b^{-1})\cdot (a+b)^{-1}] \div (a^{-1}\cdot b^{-1}) = 1$. \nl

8. If   $a \neq 0$ and $b \neq 0$, then $[(a^{-1} + b^{-1}) \div (a^{-1} - b^{-1})]\cdot [( b- a)\div (a+b)] =1.$ \nl

9. Theorem 2.2.2.

10. Theorem 2.2.4.

11. Theorem 2.2.5.

12. Theorem 2.2.6.

13. Theorem 2.2.7.\nl

\textbf{2.3 Exponents}\nl

\lhead{Section 2.3}
\textbf{Theorem 2.3.1.} If $(F,+,\cdot )$ is a field, $x\in F$, and $x \neq 0$, then $(-x)^{-1} = -(x^{-1}).$\nl

\textbf{Theorem 2.3.2.} If $(F,+,\cdot )$ is a field, $x,y\in F$, and $y \neq 0$, then $(-x)\div y =-(x\div y)$ and $x\div (-y) = -(x\div y).$\nl

\textbf{Theorem.2.3.3.} If $(F,+,\cdot )$ is a field, $x,y \in F$, and $x\cdot y=0$, then $x =0$ or $y = 0.$\nl Proof. Suppose$ (F,+,\cdot)$ is a field, $x,y \in F$, and $x\cdot y = 0$.\nl
 Suppose $y \neq 0.$\nl

\begin{tabular}{r l l}
$(x\cdot y)\cdot y^{-1}$ & $=0\cdot y^{-1}$ & Def of=, M5 \\
$x\cdot (y\cdot y^{-1})$ &$= 0$ &M3, Th 2.1.7 \\
$x\cdot 1$ &$= 0$ &M5 \\
$x$&$=0$& M4 \\
\end{tabular}
\newpage
%page32

\textbf{Theorem 2.3.4.} If $(F,+,\cdot)$ is a field and $x,y\in F$, then $x = y$ if and only
if $x - y = 0$.
Proof. Suppose $(F,+,\cdot )$ is a field and $x,y\in F.$
\begin{tabular}{r l l r l l}
$x$ &$=y$ &Hyp &$x - y$ &$= 0$ &Hyp\\
$x+(-y)$&$=y+(-y)$& Def of=,A5 &$x+(-y)$&$=y+(-y)$&Def 2.1.2,A5\\
$x-y$&$=0$& Def 2.1.2,A5 &$x$&$=y$ & Th 2.1.2 \\
\end{tabular}\nl

\textbf{Definition 2.3.1.} If $(F,+,\cdot)$ is a field, let 2=1+1, 3=2+1, 4=3+1, 5=4+1, 6=5+1, 7=6+1,   8=7+1, 9=8+1, and 10=9+1.\nl

\textbf{Definition 2.3.2.} If $(F,+,\cdot)$ is a field and $x\in F$, let $x^2$ denote $x\cdot x, x^{3}$ denote $x^{2}\cdot x,~x^{4}$ denote $x^{3}\cdot x$, and $x^{5}$ denote $x^{4}\cdot x.$\nl

\textbf{Theorem 2.3.5.} If $(F,+,\cdot )$ is a field and $x,y \in F,$ then $(x-y)\cdot(x+y)=x^{2} - y^{2}.$ \nl

\textbf{Theorem 2.3.6.} If $(F,+,\cdot )$ is a field and $x,y \in F$, then $x^{2} =y^{2}$ if and only if $x = y$ or $y = -x.$\nl
Proof. Suppose $(F,+,\cdot )$ is a field and $x,y\in F.$\nl
\begin{tabular}{r l l}
$x^{2} =y^{2}$& $\lra x^{2} - y^{2} =0$ &Th 2.3.4 \\
&$\lra(x-y).(x+y)=0$ & Th 2.3.5 \\
&$\lra x-y=0$ or $x+y=0$ &Th 2.3.3, Th 2.1.7 \\
&$\lra x=y$ or $y=-x$ & Th 2.3.4, Th 2.1.6, A5\\
\end{tabular}\nl

\textbf{Exercises}\nl

 Assuming $(F,+,\cdot )$ is a field, prove each of the following.

1. If $a,b \in F$, then $(a+b)^{2} = a^{2} + 2\cdot a\cdot b + b^{2}.$ \nl

2. If $a,b,x \in F$, then $(x+a)\cdot (x+b) = x^{2} + (a+b)\cdot x + a\cdot b.$\nl

3. $4=2^{2}.$\nl

4. $(-1)^{-1} = -1. $\nl

5. If $a\in F$, then $a^{4}=(a^{2})^{2}.$\nl

6. 5=3+2.\nl

7. If $a\in F$ and $a^{2} = 4$, then $a = 2$ or $a = -2.$\nl

8. Theorem 2.3.1.\nl

9. Theorem 2.3.2.\nl

10. Theorem 2.3.5.\nl
\newpage
%page33

\textbf{Chapter 3}\nl

\lhead{Section 3.1}
\textbf{Ordered Fields}\nl

From problem 6 of Section 2.1, we see that a field can have as few as two elements in its basic set. This indicates that a field does not completely characterize the real number system. By adding an order relation to a field we insure the fact that its basic set contains infinitely many elements.\nl

\textbf{3.1 Inequalities} \nl

\textbf{Definition 3.1.1.} \of  is said to be an \textit{ordered field} if and only if $(F,+, \cdot)$ is a field, $<$, read ``is less than", is a binary relation on $F$, and the following properties are true. \nl

O1. 	If $x,y \in F$, then exactly one of the following statements is true: $x < y, x = y, $ or $y < x.$ (Trichotomy Law) \nl

O2. 	If $x,y,z \in F, x<y,$ and $ y < z,$ then $x < z.$ (Transitive Law) \nl

O3. 	If $x,y,z \in  F$ and $x < y,$ then $x + z < y + z.$ \nl

O4. 	If $x,y,z\in F, x < y$, and $0 < z,$ then $x\cdot z < y\cdot ·z.$\nl

\textbf{Definition 3.1.2.} If \of is an ordered field, $x,y\in F$, then \nl
a) $x > y,$ read ``$x$ is greater than $y$", means $y < x.$ \nl
b) $y \geq x,$ read ``$y$ is greater than or equal to $x$", means $y > x$ or $y = x.$ \nl
c) $x \leq y,$ read ``$x$ is less than or equal to $y$", means $x < y$ or $x = y$. \nl
d) $x < y <z$ means $x < y$ and $y < z.$ \nl
e) $x \leq y \leq z$ means $x\leq y$ and $y\leq z.$ \nl
f) $x > y > z$ means $x > y$ and $y > z.$ \nl
g) $x\geq y \geq z$ means $x \geq y$ and $y \geq z.$ \nl

\textbf{Theorem 3.1.1.} If \of is an ordered field, $u,v,x,y\in F, u < v$ and $x < y,$ then $u + x < v + y.$ \nl

\textbf{Theorem 3.1.2.} If \of is an ordered field and $x,y\in F$, then the following four statements are equivalent.\nl
 a) $x < y$ \qquad b) $0 < y - x$ \qquad c) $x - y < 0$ \qquad d) $-y < -x$ \nl

\textbf{Theorem 3.1.3.} If \of is an ordered field, $x,y,z\in F, x < y,$ and $z < 0$,~then $z\cdot x > z\cdot y$ and $x\cdot z > y\cdot z.$\nl
Proof. Suppose \of is an ordered field, $x,y,z\in F, x < y,$ and $z < 0.$\nl

\newpage

\begin{tabular}{l l}
$0 < y - x$& Th 3.1.2\\
 $z\cdot (y - x) < 0\cdot (y - x)$&O4 \\
 $z\cdot y - z\cdot x < 0$ &Ex 2.1.1, Th 2.1.7 \\
 $z\cdot y <z \cdot x$& Th 3.1.2 \\
 $z\cdot x > z\cdot y$&Def 3.1.2 \\
 $x\cdot z > y\cdot z$ &M2 \\
\end{tabular}\nl

\textbf{Theorem 3.1.4.} If \of  is an ordered field, $x \in (F-\{0\})$, then $ 0 <x^{2}$. \nl
 Proof. Suppose \of  is an ordered field, $x \in (F-\{0\})$.
 By O1, we have $x < 0$ or $0 < x.$\nl

\qqu  Case 1. $x < 0$ \nl
\begin{tabular}{l l}
$x\cdot x > 0\cdot x$ & Th 3.1.3\\
$x^{2} > 0$&  Def 2.3.2, Th 2.1.7\\
\end{tabular}\nl

\qqu Case 2. $x > 0$\nl
\begin{tabular}{l l}
$x\cdot x >0\cdot x$& O4\\
$x^{2} > 0$ &Def 2.3.2, Th 2.1.7\\
\end{tabular}\nl

\textbf{Theorem 3.1.5.} If \of is an ordered field, $u,v,x,y\in F, 0 \leq u < v,$ and $0\leq x < y,$ then $u\cdot x < v\cdot y.$\nl
Proof. Suppose  \of is an ordered field, $u,v,x,y\in F, 0 \leq u < v,$ and $0\leq x < y,$\nl

\begin{tabular}{l l}
$0 < v$& O2\\
$v\cdot x < v\cdot y$& O4,M2 \\
\end{tabular}\nl

Case 1. $(0 < x)$ \qqu Case 2. $(x = 0)$
\begin{tabular}{l l l l l l}
$u\cdot x <v\cdot x$ &O4& & & $u\cdot x = 0 =v\cdot x$& Th 2.1.7 \\
$u\cdot x < v\cdot y$& O2 &&&$u\cdot x < v\cdot y$ &Def of= \\
\end{tabular}\nl

For the purposes of the examples and exercises, we will assume the results of Section 3.7. \nl

\textbf{Example 3.1.1.} Assuming \of ~is an ordered field and $x\in F$, solve the inequality $5x - 2 < 3x + 4,$ i.e., find a simpler expression which is equivalent to the given expression. Solution.
Suppose \of ~is an ordered field and $x\in F.$\nl

$5x - 2 < 3x + 4$\nl
$( 5x-2) + (-3x + 2) < (3x + 4) + (-3x+2)$ \nl
$2x < 6$\nl
$x < 3$\nl

\newpage
%page35

\textbf{Exercises}\nl

Assuming \of is an ordered field and $a,b,c,d\in F$, prove each of the following. \nl
1. If $a < 0$, then $-a > 0$. \nl
2. If $a > 0$, then $-a < 0$. \nl
3. $a^{2} =0$ if and only if $a = 0$. \nl
4. If $a\geq 0$ and $b\geq 0$, then $a\cdot b\geq 0$. \nl
5. If $a\geq 0$ and $b\leq 0$, then $a\cdot b\leq 0.$ \nl
6. If $a\leq b$ and $c\leq d,$ then $a\cdot c\leq b + d.$ \nl

7. 	The following four statements are equivalent.\nl
 a) $a\leq b$\qquad b) $a-b\leq 0$\qquad  c) $-b\leq -a$\qquad d) $0\leq b-a$\nl

8. If $a\leq b$ and $c \geq 0,$ then $a\cdot c \leq b\cdot c.$\nl

9. If $a\leq b$ and $c\leq 0$, then $a\cdot c\geq b\cdot c.$\nl

10. Theorem 3.1.1. \nl

11. Theorem 3.1.2. \nl

Assuming \of  is an ordered field and $x\in F$, solve the following inequalities.\nl

12. $4x - 1 < 2x + 3 $\nl
13. $3x - 5 > 7x + 3$\nl
14. $-3x +4 \leq 2x - 6$\nl


Assuming \of is an ordered field, find the largest element $m$ of $F$ such that for each $x$ in $F$: \nl
15. $m\leq x^{2} - 4x + 5$\nl

16. $m\leq 2x^{2} + 4x - 3 $\nl


\textbf{Theorem 3.1.6.} If \of  is an ordered field, then $0 < 1.$\nl
 Proof. Suppose \of  is an ordered field. By M4 and O1,$1 < 0$ or $0 < 1.$ Suppose $1 < 0.$ By Th 3.1.3,we have $1·\cdot 1 >0\cdot 1.$ By M4 and Th 2.1.7, we get $1 > 0.$   This is a contradiction. Therefore, $1 > 0.$ \nl

\textbf{Definition 3.1.3.} If \of  is an ordered field and $x,y\in F$, then\nl
a) $x$ and $y$ are said to have the \textit{same sign} if and only if either ($x < 0$ and $y < 0$) or ($x > 0$ and
$y > 0$).\nl
b) $x$ and $y$ are said to have \textit{different signs} if and only if either ($x < 0$ and $y > 0$) or ($x > 0$ and $y < 0$).\nl

\textbf{Theorem 3.1.7.} If \of  is an ordered field and $x,y\in F$, then

\newpage
%page36

$x\cdot y > 0$ if and only if $x$ and $y$ have the same sign. $x\cdot y < 0$ if and only if $x$ and $y$ have different signs.\nl
 Proof of (a). Suppose \of  is an ordered field and $x,y\in F.$ \nl
\quad Suppose $x$ and $y$ have the same sign. Thus ($x > 0$ and $y > 0$) or ($x < 0$ and $y <0$ ). If $x > 0$ and $y > 0,$ then by O4 and Th2.1.7, we have $x\cdot y > 0.$ If $x < 0$ and $y < 0$, then, by Th 3.1.3 and Th 2.1.7 ,we have $x\cdot y > 0$. In either case $x\cdot y > 0.$\nl
$ $\quad By a similar argument, we see that if $x$ and $y$ have different signs, then $x\cdot y < 0.$\nl
$ $\quad Suppose $x\cdot y > 0.$ By Th 2.1.7 , we see that $x\neq 0 \neq y.$ Hence $x$ and $y$ have the same sign or different signs. If $x$ and $y$ have different signs, then $x\cdot y < 0.$ Therefore $x$ and $y$ have the same sign. \nl
By a similar argument, we have that if $x\cdot y < 0,$ then $x$ and $y$ have different signs. \nl

\textbf{Theorem 3.1.8.} If \of  is an ordered field and $x\in F$, then $x$ and $x^{-1}$ have the same sign. \nl

\textbf{Theorem 3.1.9.} If \of  is an ordered field, $x,y\in F, x \geq 0,$ and $y\geq 0,$ then $x^{2} < y^{2}$ if and only if $x < y.$\nl
 Proof. Suppose \of is an ordered field, $x,y\in F, x \geq 0,$ and $y\geq 0$\nl
If $x < y,$ then, by Th 3.1.5, we have $x^{2} < y^{2}.$\nl
Suppose $x^{2} <y^{2}.$ By O1, we have $x < y, x = y,$ or $y < x.$ If $y < x,$ then, by Th 3.1.5, we obtain $y^{2} < x^{2}$, which is a contradiction. If $x = y$, then $x^{2} = y^{2},$ which is also a contradiction. Therefore $x  < y.$ \nl

\textbf{Example 3.1.2.} Assuming \of is an ordered field and $x\in F$, solve the inequality $x^{2} + 4x - 5< 0.$\nl Solution. Suppose \of is an ordered field and $x\in F$.\nl
\begin{tabular}{l l}
$x^{2} + 4x - 5 < 0$&\\
&$\lra (x+5)\cdot(x-1) < 0$\\
&$\lra (x+5 <0$ and $x-1>0)$ or $(x+5>0$ and $x-1 < 0)$\\
&$\lra (x<-5$ and $x> 1)$ or $(x>-5$ and $x< 1)$\\
&$\lra -5<x<1$\\
\end{tabular}\nl

\textbf{Example 3.1.3.} Assuming \of is an ordered field and $x\in F,$ solve the inequality $4\div(x-2) < 2.$\nl
 Solution. Suppose \of is an ordered field and $x\in F,$

\newpage
%page37
\begin{tabular}{l l}
%$4\div(x-2) < 2$&\\
$4\div(x-2) < 2$&$\lra 4\div (x-2) < 2$ and $(x-2>0$ or $x-2<0)$\\
&$\lra (4 < 2\cdot (x-2)$ and $x - 2 >0)$ or $(4>2\cdot (x-2)$ and $x - 2 < 0)$\\
&$(4 < 2x - 4$ and $x > 2)$ or $(4 > 2x - 4$ and $ x < 2)$\\
&$(4 < x$ and $x > 2)$ or $(4 > x$ and $x < 2)$ \\
&$\lra x > 4$ or $x < 2$\\
\end{tabular}\nl

\textbf{Exercises}\nl
 Assuming \of is an ordered field, prove each of the following. \nl

1. If $a\in F$ and $a > 1,$ then $a^{2} > a.$ \nl

2. If $a\in F$ and $0 < a < 1$, then $a^{2} < a.$\nl

3. If $a,b\in F$ and $a < b,$ then $a < (a+b)\div 2 < b.$ \nl

4. Theorem 3.1.8. \nl

Assuming \of  is an ordered field and $x\in F$, solve each of the following inequalities.

5. $x^{2} -x - 12 > 0$\nl

6. $x^{2} - 2x - 15 < 0$\nl

7. $x^{2} > 0$\nl

8. $-1 < 3x + 5 < 2 $\nl

9. $2x - 3 < 4x + 5 <-3x + 2$\nl

10. $x\div(x-1) > 0 $\nl

11. $0 < 1\div  x^{2} < 9 $\nl

12. $0 < 1\div(x-2) < 4$ \nl

\lhead{Section 3.2}
\textbf{3.2 Intervals }\nl

\textbf{Definition 3.2.1.} Suppose \of is an ordered field, $a,b\in F$, and $a < b.$\nl
 a) The \textit{open interval} from $a$ to $b$ is denoted by $(a,b)$ and defined by
$(a,b)=\{ x \mid a < x < b \}.$\nl
 b) The \textit{closed interval} from $a$ to $b$ is denoted by [$a,b$] and defined by
$[a,b] =\{x \mid a \leq x \leq b \}.$ \nl
c) There are two \textit{half-open intervals} from a to b denoted by $(a,b]$ and $[a,b)$ and defined by $(a,b] = \{x \mid a < x \leq b\}$ and $[a,b) = \{x \mid a \leq x < b\}.$\nl
d) Each of the following sets are called \textit{infinite intervals.} \nl

\newpage
%page38
\begin{tabular}{l l}
$(a,\infty)= \{x \mid a < x\}$& $(-\infty,a)= \{x \mid x < a\}$ \\
$[a,\infty) = \{x \mid a \leq x\}$& $(-\infty, a] = \{ x \mid x \leq a \}$ \\
$(-\infty, \infty) = F$& \\
\end{tabular}\nl

\textbf{Example 3.2.1.} Assuming \of is an ordered field, express the truth set of $[9\div(x-2)] > 3$ in terms of intervals.\nl

Solution. Suppose \of is an ordered field.\nl
$ \{x \mid [9\div (x-2)] >3\}$ \nl
$=\{x \mid 9\div (x-2)) > 3$ and $(x-2 > 0 $ or $x-2< 0)\}$ \nl
$= \{(9\div (x-2)$ and $x - 2 > 0)$ or $(9 <3\cdot (x-2)$ and $x - 2 <0)\}$ \nl
$= \{x \mid 9 > 3x-6$ and $x > 2)$ or $(9 < 3x - 6$ and $x < 2)\} $\nl
$= \{x \mid (5 > x$ and $x > 2)$ or $(5 < x$ and $x < 2)\}$ \nl
$=\{x \mid 2< x < 5\}$ \nl
$=(2,5)$ \nl

\textbf{Example 3.2.2.} Assuming \of  is an ordered field, express the truth set of $x^{2} -2x + 8 > 0$ in terms of intervals.\nl Solution  Suppose \of is an ordered field. \nl
$\{x \mid x^{2} - 2x - 8 > 0 \}$ \nl
$= \{x \mid (x - 4 > 0$ and $x + 2 > 0)$ or $(x - 4 < 0$ and $x + 2 < 0)\} $\nl
$= \{x \mid (x > 4)$ and $x > -2)$ or $(x < 4$ and $x < -2)\} $\nl
$= \{x \mid x > 4$ or $x < -2\} $\nl
$=\{x \mid x > 4\} \cup \{x \mid x < -2\} $\nl
$=(4,\infty) \cup (-\infty,-2) $\nl

\textbf{Exercises}\nl

 Assuming \of  is an ordered field, express each of the following in terms of intervals.

1. $\{x \mid -2< x < 3\}$ \nl

2. $\{x \mid x <1$ or $x > 2\}$ \nl

3. $\{x\mid x < 2$ and $x\leq 4\}$ \nl

4. $\{x \mid -5 < x$ and $x < 6\}$  \nl

5. $\{x \mid x < 2$ or $x < 6\}$ \nl


Assuming \of  is an ordered field, express each of the following as an interval. \nl

6. $[2,5) \cap [3,10)$\nl

7. $(0,6)\cap (1,4)$\nl

8. $(-\infty,2] \cup (0,\infty) $

\newpage
%page39

9. $(-\infty,1] \cup [0,3)$\nl

10. $(-\infty,2)\cap(0,2]$ \nl

11. $(-3,1)\cup[1,5]$ \nl


Assuming \of  is an ordered field, express the truth set of each of the following in terms of intervals. \nl

12. $[8\div(x-1)] < 4 $\nl

13. $x^{2} + x - 6 < 0 $\nl

14. $x^{2} + 3x - 10 > 0 $\nl

15. $x^{2} < 4 $\nl

16. $0 < [4\div (x-2)] < 1 $\nl

17. $x^{2} > 1 $\nl


\lhead{Section 3.3}
\textbf{3.3 Absolute Value}\nl

\textbf{Definition 3.3.1.} If \of  is an ordered field and $x\in F,$ then the \textit{absolute value} of $x$ is denoted by $\abs{x}$ and defined by \nl
\begin{tabular}{l l}
$\abs{x}$ &$=x$ if $x\geq 0$\\
&$=-x$ if $x < 0$\\
\end{tabular}\nl

\textbf{Theorem 3.3.1.} If \of  is an ordered field and $a,x\in F,$ then\nl
(1) $\abs{x} = a$ if and only if $a \geq 0$ and $(a = x$ or $a = -x)$ \nl
and\nl
(2) $\abs{x} \geq 0$ and $( \abs{x} = x$ or $\abs{x} = -x).$\nl
Proof. Suppose \of is an ordered field and $a,x\in F$.\nl
Proof of Assertion (1)
\begin{tabular}{l l}
$\abs{x}=a$& $\lra \abs{x}=a$ and $(x \geq 0$ or $x\leq 0)$\\
&$\lra (\abs{x} = a$ and $x\geq 0)$ or $(\abs{x} = a$ and $x\leq 0) $\\
&$\lra (x = a$ and $x\geq 0)$ or $(-x = a$ and $x\leq 0)$\\
&$\lra  (x\geq 0$ and $x = a)$ or $(-x\geq 0$ and $-x = a)$\\
&$\lra (a\geq 0$ and $a = x)$ or $(a \geq 0$ and $a = -x)$\\
&$\lra a\geq 0$ and $(a = x$ or $a = -x)$\\
\end{tabular}\nl
Proof of Assertion (2)
Since (1) holds for each number $a,$ it holds for $a=\abs{x}.$\nl
Thus $\abs{x} = \abs{x} \lra  [\abs{x} \geq 0$ and  $(\abs{x} = x$ or $\abs{x} = -x)].$ Therefore \nl

\newpage
%page40


$\abs{x} \geq 0$ and $(\abs{x} = x$ or $\abs{x} = -x)$\nl

\textbf{Example 3.3.1.} Suppose \of is an ordered field and $x\in F.$ Solve the equation $\abs{x + 2}=3x - 6.$\nl
Solution. Suppose \of is an ordered field and $x\in F.$ \nl
$\abs{x + 2} = 3x - 6$\nl
$\lra 3x-6 \geq 0$ and $(3x - 6 = x + 2$ or $3x - 6 = -(x+ 2))$\nl
$\lra x\geq 2$ and $(x=4$ or $x=1)$\nl
$\lra x=4$\nl

\textbf{Theorem 3.3.2.} If \of is an ordered field, $a,x\in F$, and $a > 0$, then
(1) $\abs{x} < a$ if and only if $-a < x < a$\nl
and\nl
(2) $\abs{x}\leq a$ if and only if$ -a \leq x \leq a.$\nl
Proof of (1). Suppose \of is an ordered field, $a,x\in F.$\nl
\begin{tabular}{l l}
$\abs{x} < a$ &$\lra \abs{x}<a$ and $(0 \leq x$ or $x\leq 0)$\\
&$\lra (\abs{x} < a$ and $0 \leq x)$ or $(\abs{x} < a$ and $x\leq 0)$\\
&$\lra (x < a$ and $-x\leq 0$ and $0\leq x)$ or $(-x < a$ and $0\leq -x$ and $x\leq 0)$\\
&$\lra (x<a$ and $-x < a$ and $0\leq x)$ or $(-x<a$ and $x<a$ and $x\leq 0)$\\
&$\lra (x < a$ and $-x < a)$ and $(0\leq x$ or $x\leq 0)$\\
&$\lra x < a$ and $x > -a$\\
&$\lra -a < x < a$\\
\end{tabular}\nl

\textbf{Example 3.3.2.} Assuming \of is an ordered field, find the truth set of $\abs{3x - 1}\leq 2x + 6$ in terms of intervals. \nl
Solution. Suppose \of is an ordered field.\nl
$\{x \mid \abs{3x - 1}\leq 2x+6\}$\nl
$=\{x \mid -(2x+6)\leq 3x-1 \leq 2x+6\}$\nl
$= \{x \mid -2x-6\leq 3x-1$ and $3x-1\leq 2x+6\}$\nl
$= \{x \mid -1\leq x$ and $x\leq 7\}$\nl
$=[-1,7] $\nl

\textbf{Theorem 3.3.3.} If \of is an ordered field and $a,x\in F,$ then \nl
(1) $\abs{x} > a$ if and only if $x < -a$ or $a < x,$ and \nl
(2) $\abs{x} \geq a$ if and only if  $x\leq -a$ or $a\leq x.$ \nl

\textbf{Example 3.3.3.} Assuming \of is an ordered field, find the set of $\abs{3x - 1} > 2x + 6$ in terms of intervals. \nl

\newpage
%page41


Solution.
~$\{x\mid \abs{3x - 1} > 2x + 6\} $\nl
$=\{x \mid -(3x-1) > 2x + 6$ or $3x - 1 > 2x + 6\}$\nl
$=\{x \mid 3x - 1  < -2x - 6$ or $2x + 6 < 3x - 1\}$ \nl
$=\{x\mid  x < - 1 $ or $7< x\}$\nl
$=\{x \mid x < -1\}\cup \{x \mid 7 < x\} $\nl
$=(-\infty,-1) \cup (7,\infty)$\nl

\textbf{Theorem 3.3.4.} If \of is an ordered field and $x\in F,$ then $x\leq \abs{x}.$\nl

\textbf{Theorem 3.3.5.} If \of is an ordered field and $x\in F,$ then $\abs{x}= \abs{-x}.$\nl

\textbf{Theorem 3.3.6.} If \of  is an ordered field and $x \in F,$ then $-x \leq \abs{x}$ and $-\abs{x} \leq x$ and $-\abs{x} \leq x \leq \abs{x}.$\nl

\textbf{Exercises}\nl

Assuming  \of is an ordered field, find the truth set of each of the following in terms of intervals. \nl

1. $\abs{2x - 3} < x $\nl
2. $\abs{x + 2} \leq 3x - 6 $\nl
3. $\abs{4x - 3} < 3x - 4 $\nl
4. $\abs{4x - 1} - 5 < 3x$\nl
5. $\abs{4x - 3} = 3x -4 $\nl
6. $\abs{4x - 1} = 3x + 5$\nl
7. $\abs{3x - 1} = 2x + 6$\nl
8. $\abs{x + 2} > 3x - 6$ \nl
9. $\abs{4x - 3} - 3x > -4$\nl
10. $\abs{4x - 1} \geq 3x + 5$\nl

Assuming \of is an ordered field and $x\in F,$ prove each of the following.\nl
11. $\abs{3x - 1} < 2 \lra 1 < x < 5 $\nl
12. $\abs{-2x + 5} <3  \lra  1 < x < 4$\nl
13. Theorem 3.3.3(1). (Hint: Negate both sides of part(2) of Theorem 3.3.2) \nl

14. Theorem 3.3.4. (Hint: Use Theorem 3.3.3) \nl

15. Theorem 3.3.5. (Hint: Use Theorem 3.3.1) \nl

16. Theorem 3.3.6. (Hint: Use Theorems 3.3.4 and 3.3.5) \nl

\newpage
%page42
\textbf{Theorem 3.3.7.} If \of  is an ordered field and $x,y \in F,$ then $\abs{x} = \abs{y}$ if and only if $x = y$ or $x = -y.$\nl
 Proof. Suppose \of  is an ordered field and $x,y \in F.$ \nl
 \begin{tabular}{r l}
$\abs{x} =\abs{y}$ &$\lra \abs{y} \geq 0$ and $(\abs{y} = x$ or $\abs{y}= -x)$\\
&$\lra \abs{y} = x$ or $\abs{y} = -x$\\
&$\lra [x \geq 0$ and $(x = y$ or $x = -y)]$ or $[-x\geq 0$ and $(-x = y$ or $-x = -y)]$\\
&$\lra [x \geq 0$ and $(x=y$ or $x=-y)]$ or $[x\leq 0$ and $(x=y$ or $x=-y)] $\\
&$\lra (x=y$ or $x=-y)$ and $(x \geq 0$ or $x \leq 0)$\\
&$\lra x = y$ or $x = -y$\\
\end{tabular}\nl

\textbf{Example 3.3.4.} Assuming \of is an ordered field, find the truth set of $\abs{2x - 2} = \abs{x - 4}.$\nl
Solution. Suppose \of is an ordered field. \nl
$ $ $\{x \mid \abs{2x - 2} = \abs{x - 4} \}$ \nl
$= \{x \mid 2x - 2 = x-4$ or $2x - 2= -x + 4 \}$\nl
$= \{x \mid x=-2$ or $x = 2\}$\nl
$=\{-2,2\}$\nl

\textbf{Theorem 3.3.8.} If \of is an ordered field and $x,y\in F$, then $\abs{x + y} \leq \abs{x} + \abs{y}.$\nl
Proof. Suppose \of is an ordered field and $x,y\in F.$ By Th 3.3.6, we have
\begin{tabular}{l }
$\,\,\,\,\,\,-\abs{x} \leq x \leq \abs{x}$ and \\
$-\abs{y}\leq y \leq \abs{y}.$\\
\end{tabular}\nl
By adding these two inequalities, we get $-(\abs{x} + \abs{y}) \leq x+y \leq \abs{x} + \abs{y}.$ By Th 3.3~2 we obtain $\abs{x + y} \leq \abs{x} + \abs{y}.$\nl

\textbf{Example 3.3.5.} Assuming \of is an ordered field and $x\in (2,4),$ find a number greater than $\abs{x^4 - 2x^{2} + 6x + 8}.$\nl
Solution. Suppose \of is an ordered field and $x\in (2,4).$ Thus $2<x<4$ and $\abs{x} < 4.$\nl
$\abs{x^4 - 2x^{2} + 6x + 8}.$\nl
$\leq \abs{x^4} + \abs{-2x^{2}}+ \abs{6x} + \abs{8}$\nl
$\leq x^4 + \abs{-2}x^{2} + 6\abs{x} +8 $\nl
$<(4^{4})+2(4^{2}) + 6(4) + 8$ \nl
=256 + 32 + 24 + 8 = 320 \nl

\newpage
%page43

\textbf{Theorem 3.3.9.} If \of is an ordered field and $x,y\in F,$ then
$\abs{x} - \abs{y} \leq \abs{x - y}$ and $\abs{\,\,\abs{x}- \abs{y}\,\, } \leq \abs{x - y}.$\nl

\textbf{Theorem 3.3.10.} If \of is an ordered field and $x,y\in F,$ then
$\abs{x}\cdot\abs{y} = \abs{x\cdot y}. $\nl
Proof. Suppose  \of is an ordered field and $x,y\in F.$  Let $a=\abs{x}$ and $b=\abs{y}.$ By Th 3.3.1, we have $[a \geq 0$ and $(a = x$ or $a = -x)]$ and $[b \geq 0$ and $(b = y$ or $b = -y)].$ Therefore $a\cdot b  \geq 0$  and $(a\cdot b = x\cdot y$ or $a\cdot b = -(x\cdot y).$ Applying Th 3.3.1 again, we get $\abs{x\cdot y} = a\cdot b =\abs{x}\cdot \abs{y}.$\nl

\textbf{Theorem 3.3.11.} If If \of is an ordered field and $x,y\in F,$ and $y\neq 0,$ then $\abs{x \div y} = \abs{x} \div \abs{y}.$\nl

\textbf{Theorem 3.3.12.} If  \of is an ordered field and $x\in F$ then ${\abs{x}}^{2}  = x^{2}. $\nl

\textbf{Exercises} \nl

Assuming \of  is an ordered field, find the truth set of each of the following. \nl
1. $ \abs{3x - 1} = \abs{-2x + 4} $\nl

2. $\abs{4x - 8} = -\abs{3x}$ \nl

3. $\abs{x} = \abs{y}  \lra x^{2}= y^{2}.$ \nl

4. $ \abs{x - y} \leq \abs{x} + \abs{y}$ \nl

5. $\abs{z - x} \leq \abs{z - y} + \abs{y - x}$\nl

6. $\abs{x}= 0 \lra x=0$ (Hint: Use Theorem 3.3.1)\nl

7. $x \in (a - c, a + c) \lra  \abs{x - a} < c $\nl

8. $a-c < x < a + c \lra \abs{x - a} < c $\nl

9. If $\abs{ x + y} < c,$ then $\abs{x} \leq \abs{y} + c.$\nl

10. If $\abs{x} < \dfrac{c}{2}$ and $\abs{y} < \dfrac{c}{2},$ then $\abs{x + y} < c.$\nl

11. If $\abs{x} < \dfrac{c}{\abs{a} + 1}$, then $\abs{ax} < c.$\nl

12. Theorem 3.3.9. (Hint: Apply Theorem 3.3.8 to $y$ and $x -y$)

13. Theorem 3.3.11.

14. Theorem 3.3.12. (Hint: Use Theorems 2.3.6 and 3.3.1)

\newpage

%page44
\lhead{Section 3.4}
\textbf{3.4 Mathematical Induction}\nl

\textbf{Definition 3.4.1.} If \of  is an ordered field, then a subset $S$ of $F$ is called an \textit{inductive set} if and only if (1) $1\in S$ and (2) $k\in S$ implies $k+1\in S.$\nl

\textbf{Definition 3.4.2.} If \of  is an ordered field, let $Z^{+}$ denote $\{n \mid n$ belongs to every inductive subset of $F \}$. The elements of  $Z^{+}$ are called \textit{positive integers.} \nl

\textbf{Theorem 3.4.1.} If \of  is an ordered field, then $Z^{+}$ is an inductive set.\nl
Proof. Suppose \of is an ordered field. \nl

(1) $1\in Z^{+}$ since 1 belongs to every inductive subset of $F.$\nl

(2)\begin{tabular}{r l} 	
$k\in Z^{+}$ &$\lra k$ belongs to every inductive subset of $F$\\
&$\lra k+1$ belongs to every inductive subset of $F$ \\
&$\lra k+1\in Z^{+} $\\
\end{tabular}\nl

\textbf{Theorem 3.4.2.} If \of is an ordered field and $S$ is an inductive subset of $Z^{+},$ then $S =Z^{+}.$\nl 
Proof. Suppose \of is an ordered field and $S$ is an inductive subset of $Z^{+}.$

\begin{tabular}{r l}
$n \in Z^{+}$& $\ra n$ belongs to every inductive subset of $F$.\\
&$\ra n\in S$\\
\end{tabular}\\
Thus $Z^{+} \subseteq S$ and $S \subseteq Z^{+}.$ Therefore $S = Z^{+}.$ \nl

\textbf{Theorem 3.4.3.} If \of is an ordered field and $n\in Z^{+},$ then $n \geq 1.$\nl

\textbf{Definition 3.4.3.} If \of is an ordered field and $a\in F,$ let $a^{0} = 1$ and for each $n\in Z^{+},$ let $a^{n} =a^{n-1}\cdot a.$\nl

Note: For the remainder of this book, we will let $ab$ denote $a\cdot b.$\nl

\textbf{Theorem 3.4.4.} If \of  is an ordered field, $a,b\in F, b \neq 0,$ and $n,m\in Z^{+}$, then \nl
a) $a^{1} = a$ \nl
b) $m + n \in Z^{+}$\nl
c) $mn\in Z^{+} $\nl
d) $m^{n}\in Z^{+} $\nl

\newpage

%page45

e) $a^{m+n} = a^{m}a^{n}$\nl
f) $(a^{m})^n = a^{mn} $\nl
g) $(a^{n})\div (b^{n})=(a\div b)^{n}$\nl
h) $a^{n} \div a^{m} = a^{n - m}$\nl
Proof of (b). Suppose \of  is an ordered field and $ m \in\zp , m + n \in \zp $.  \nl
Let $S = \{ n \mid n \in \zp$ and $ m + n \in \zp \}$.\nl
\qqu Proof of Assertion (1)\nl
$1 \in S$ since $1\in \zp$ and $m + 1\in \zp$\nl
$\quad \quad$ Proof of Assertion (2)\nl
\begin{tabular}{rl}
$k \in S$&$\ra  k\in \zp $ and $m + k \in \zp$\\
%&$\ra k \in \zp$ and $(m + k) + 1\in \zp$\\
&$\ra k + 1 \in \zp$\\ %and $m+(k+1)\in \zp$\\
&$\ra k+1\in S$\\
\end{tabular}\\
By Th 3.4.2, we have $S=\zp$ and for each $n\in \zp, m + n\in \zp .$\nl

Proof of (f). Suppose \of is an ordered field, $a\in F,$ and $m\in \zp$. Let $S= \{n \mid n \in \zp $ and $(a^{m})^{n} = a^{mn}\}.$ \nl
\qqu Proof of Assertion (1)\nl
$1\in S$ since $1 \in \zp$ and $(a^{m})^{1} = a^{m} = a^{m1} $\nl
\qqu Proof of Assertion (2)\nl
\begin{tabular}{r l}
$k \in S$&$\ra  k\in \zp $ and $(a^{m})^{k} = a^{mk}$\\
%&$\ra k + 1\in \zp$ and $(m + k) + 1\in \zp$\\
&$\ra k+1 \in \zp$ and $(a^{m})^{k+1} =  a^{m(k+1)}$\\
&$\ra k+1\in S$\\
\end{tabular}\\
By Th 3.4.2 we have $S = \zp$ and for each $n \in \zp, (a^{m})^{n} = a^{mn}.$\nl

\textbf{Theorem 3.4.5.} If \of is an ordered field and $n\in \zp$ then $n = 1$ or $n-1\in \zp$.\nl
Proof. Suppose \of is an ordered field. Let $S=\{ n \mid n\in \zp$  and $(n=1$ or $n-1\in \zp)\}$\nl
\qqu Proof of Assertion (1)\nl
$1 \in S$ since $1 \in \zp$ and $(k =1$ or $k -1 \in \zp ).$\nl
\qqu Proof of Assertion (2)\nl
\begin{tabular}{rl}
$k \in S$&$\ra k\in \zp$ and $k=1 $ or $k - 1 \in \zp$\\
&$\ra k \in \zp$ and $(k - 1) + 1\in \zp$\\
&$\ra k+1 \in \zp$ and $(k + 1) - 1 \in \zp$\\
&$\ra k+1\in S$\\
\end{tabular}\\

By Th 3.4.2, we have $S = \zp$ and for each $n\in \zp, n = 1$ or $n - 1\in \zp .$\nl

\textbf{Theorem 3.4.6.} If \of is an ordered field and $n \in \zp,$ then $\zp \cap (n, n+1) = \ns .$\nl
Proof. Suppose \of is an ordered field. Let $S=\{n \mid n\in \zp$ and $\zp \cap (n,n+1)= \ns \}.$ \nl
(1) Suppose there is an element $x$ in $\zp \cap (1,2).$ Thus $x\in \zp $ and $1 < x < 2.$\nl
By Th 3.4.5,we have $x-1\in \zp$ and $x - 1 < 1.$ This contradiction proves $1\in S.$\nl

\newpage

%page46


(2) Suppose $k\in S$ and $k+1\notin S$. Thus there is an element $x \in \zp \cap (k+1,k+2).$ Hence $x\in \zp$ and $k+1<x<k+2.$ Since 1 precedes every other element in $\zp , x\neq 1.$ By Th 3.4.5, we have $x-1\in \zp$ and $k < x-1 < k+1.$ There­fore $x-1\in \zp \cap (k,k+1).$ This contradiction proves $k\in S \ra k+1\in S.$ By Th 3.4.2, we have $S=\zp$ and for each $n\in \zp, \zp \cap (n,n+1) = \ns .$ \nl

\textbf{Definition 3.4.4.} If \of is an ordered field, let $Z= \{n \mid n\in \zp, n =0,$ or $-n\in \zp \}.$ The elements of $Z$ are called \textit{integers.}\nl

\textbf{Theorem 3.4.7}  If \of  is an ordered field and $x,y\in Z,$ then $xy\in Z.$\nl

\textbf{Theorem 3.4.8.} If \of  is an ordered field and $n\in Z,$ then $n-1\in Z.$\nl

\textbf{Theorem 3.4.9.} If \of  is an ordered field, $a\in Z,$ and $n\in \zp$, then $a - n\in Z.$\nl

\textbf{Theorem 3.4.10.} If \of is an ordered field and $x,y\in Z,$ then $x+y\in Z$ and $x - y \in Z$.\nl

\textbf{Theorem 3.4.11.}. If \of  is an ordered field, then every non-empty
subset of $\zp$ has a least element.\nl
 Proof. Suppose \of is an ordered field and $S$ is a non-empty subset of $\zp$ which does not have a least element. Since 1 precedes every positive integer, $S$ does not contain $1$. Let $M=\{n \mid n\in \zp $ and $n$ is less than every element in $S\}.$ Thus $1\in M.$ If $M$ is an inductive set, then $M = \zp$ and $S = \ns.$ Therefore there is an element $k$ in $M$ such that $k+1\in S.$ Since $S$ has no least element, there is an $h$ in $S$ such that $h < k+1$. Thus, $k < h < k + 1,$ since $k$ is less than every element of $S$. By Th 3.4.6, we have $h \notin \zp$ and $h \notin S$. This contradiction proves the theorem.\nl


\textbf{Definition 3.4.5.} If \of  is an ordered field, then for each $a$ in $Z$, let $Z_{a} = \{n \mid n\in Z$ and $n \geq a\}.$\nl

\textbf{Theorem 3.4.12.} If \of  is an ordered field, $a\in Z$ and $S$ is a subset of $Z_{a}$ such that (1) $a\in S$ and (2) $k\in S$ implies $k+1\in S,$ then $S =Z_{a}.$ \nl
Proof. Suppose \of  is an ordered field, $a\in Z, S\subseteq Z_{a}, a\in S,$ and $k\in S$ implies $k+1\in S.$ Let $T =\{n \mid n\in \zp$ and $ n-1+a\in S\}.$\nl
(1) $1\in T$ since $1\in \zp$ and $a\in S$\nl
(2) $k\in T \ra k\in \zp $ and $k - 1 + a\in S $\nl


\newpage


%page47
$\ra k\in\zp $ and $[k-1+ a] + 1 \in S$\nl
$\ra k + \in \zp$ and $(k+1) - 1 + a\in S$\nl
$\ra k+1\in T$\nl
Therefore $T = \zp$ and for each $n\in \zp, n - 1 + a\in S.$ Suppose $x\in Z_{a}$.  Thus $x\geq a$ and $x \in Z.$ Let $n = x + 1 - a.$ Hence $n \geq 1$ and $n\in Z.$ Therefore $n\in \zp, n - 1 + a\in S,$ and $x\in S.$ Thus $Z_{a} \subseteq S$ and $S \subseteq Z_{a}.$ Therefore $S = Z_{a}.$\nl

\textbf{Example 3.4.1.} If \of is an ordered field and $n\in Z_{2},$ then $2^{n}  > 2.$ \nl
 Solution. Suppose \of is an ordered field. Let $S=\{n \mid n\in Z_{2}$ and $2^{n} >2\}.$\nl
(1) $2\in S$ since $2\in Z_{2}$ and $2^{2} >2$\nl
Proof of Assertion (2)\nl
\begin{tabular}{r l}
$k\in S$ & $\ra k\in Z_{2}$ and $2^{k} >2$\\
&$\ra k+1\in Z_{2}$ and $2^{k+1} = 2^{k} \cdot 2 > 2\cdot 2 > 2 $\\
&$\ra k+1\in S$\\
\end{tabular}\\

By Th 3.4.2, we have $S=Z_{2}$ and for each $n$ in $Z_{2}, 2^{n} > 2.$\nl

\textbf{Theorem 3.4.13.} If \of  is an ordered field and $S$ is a subset of $\zp$ such that (1) $1\in S$ and (2) $Z^{n}_{1} \subseteq S$ implies $n+1\in S,$ then $S = \zp.$\nl

\textbf{Definition 3.4.6.} If \of is an ordered field and $x,y\in Z$, then $y$ is said to be \textit{divisible} by $x$ if and only if $y / x \in Z.$\nl

\textbf{Exercises}\nl
Suppose \of  is an ordered field, $A =\{x \mid x\in \zp$ and $x > 0\},$\nl $B= \{x \mid x\in F$  and $x^{2}\geq  2x\},$ and $C =\{x \mid x\in F$ and $x^{2} \geq 1\}.$\nl
1. State which of the sets $F, A, B,$ and $C$ are inductive. \nl

2. State which of the sets $F, A, B,$ and $C$ contain $\zp .$\nl

3. State which of the sets $F, A, B,$ and $C$ are equal to $\zp$. \nl
Assuming  \of is an ordered field, prove each of the following.

4. If $n\in Z_{3},$ then $2^{n} > 2n + 1. $\nl

5. If $n\in Z_{5},$ then $2^{n} > n^{2}.$\nl

6. If $n\in \zp,$ then $n^{2} + n$ is divisible by 2.\nl

7. If $n\in \zp,$ then $n^{3}+ 2n$ is divisible by 3.\nl

8. 	If $x,y\in F, n\in \zp,$ and $x > y,$ then $x^{n} - y^{n}$ is divisible by $x - y$\nl
 (Hint: $x^{n+1} - y^{n+1} = x(x^{n} -y^{n} ) + y^{n}(x-y).) $\nl

9. Theorem 3.4.3.

10. Theorem 3.4.4(a).

11. Theorem 3.4.4(c).

12. Theorem 3.4.4(d).

\newpage


%page48
13. Theorem 3.4.4(e).\nl

14. Theorem 3.4.4(g).\nl

15. Theorem 3.4.4(h).\nl

16. If $a\in F, n\in \zp,$ and $ m = 2n + 1,$ then $a^{m} = -(-a)^{m}.$ \nl

17. If $a\in F, n\in \zp$ ,and $m = 2n,$ then $a^{m} =(-a)^{m}.$ \nl


\textbf{Definition 3.4.7.} If \of is an ordered field, $h,k\in Z,$ and $~h < k,$  let $Z^{k}_{n}$~ denote $\{n \mid n\in Z$ and  $h\leq n\leq k\}.$\nl

\textbf{Definition 3.4.8.} A function $f$ is said to be \textit{one-to-one} if and only if no two ordered pairs of $f$
have the same second term, i.e., $x,y\in D_{f} $ and $f(x)= f(y) $ imply $x=y.$\nl

\textbf{Definition 3.4.9.} If \of  is an ordered field, then a set $A$ is said to have $n$ elements if and
only if $n\in \zp$ and there is a one-to-one function from $Z^{n}_{1}$ onto $A.$\nl

\textbf{Definition 3.4.10.} If \of is an ordered field, then a set $A$ is said to be \textit{finite} if and only if $A$
is empty or there is an element $n$ in $\zp $ such that $A$ has $n$ elements. \nl

\textbf{Definition 3.4.11.} If \of  is an ordered field, then a set $A$ is said to be \textit{infinite} if and only if
$A$ is not finite. \nl

\textbf{Definition 3.4.12.} If \of  is an ordered field, then $a$ is said to be an \textit{infinite sequence} of
elements of a set $A$ if and only if $a : \zp \ra A.$ If $a$ is an infinite sequence, then $a(i)$ is denoted
by $a_{i}$ for each $i$ in $\zp$ and $a$ is sometimes denoted by $a_{1}, a_{2}, a_{3},\, \dots \,\,\,.$\nl

\textbf{Definition 3.4.13.} If \of  is an ordered field and $a$ is an infinite sequence of elements of
$F,$ let  $\sum_{i=1}^{1}a_{i}$ denote $a_{1}$ and for each $n$ in $\zp$, let
$\sum_{i=1}^{n+1}a_{i}$ denote $\left( \sum_{i=1}^{n}a_{i}\right) + a_{n+1} .$\nl

\textbf{Definition 3.4.14.} If \of is an ordered field, $a : Z \ra F, k\in Z,$ and $n \in \zp$, let $\sum_{i=k+1}^{k+n}a_{i}$ denote $\sum_{i=1}^{n}a_{i+k}.$\nl

\textbf{Example 3.4.2.} If \of is an ordered field, then

\newpage
%page49

a) $\sum_{i=1}^{3}(i^{2} - 1) = (1^2 - 1) + (2^2 - 1) + (3^2 - 1) = 11$ \nl

b)  $\sum_{i=3}^{5}(2i-3) = (2\cdot 3-3)+(2\cdot 4-3)+ (2\cdot 5 - 3) = 15 $\nl

\textbf{Theorem 3.4.14.} If \of is an ordered field, $a$ and $b$ are infinite sequences of elements of $F, c\in F,$ and $n\in \zp,$ and $k \in Z_{1}^{n-1}$ then\nl

a)  $c\sum_{i=1}^{n}a_{i} = \sum_{i=1}^{n}(c\cdot a_{i}) \hfill c)\sum_{i=1}^{n}a_{i} = \sum_{i=1}^{k}a_{i} +\sum_{i=k+1}^{n}a_{i}.$\nl

b)  $\sum_{i=1}^{n}a_{i} + \sum_{i=1}^{n}b_{i} = \sum_{i=1}^{n}(a_{i} + b_{i})$\nl


Proof of (a). Suppose \of  is an ordered field, a is an infinite sequence of elements of $F,$
and, $c\in F.$ \nl
Let $S= \{n \mid n\in \zp$ and $c\sum_{i=1}^{n}a_{i} = \sum_{i=1}^{n}(c\cdot a_{i})\}$\nl

(1) $1\in S$ since $1\in \zp$ and $ca_{1} = ca_{1}.$\nl

\begin{tabular}{r l}
(2) $k\in s$&$\ra k\in \zp$ and $c\sum_{i=1}^{k}a_{i} = \sum_{i=1}^{k}(ca_{i})$\\
&$\ra k\in \zp$ and $c\left(\sum_{i=1}^{k}a_{i}\right) + ca_{k+1} = \left(\sum_{i=1}^{k}(ca_{i})\right) + ca_{k+1}$\\
&$\ra k\in \zp$ and $c\left(\sum_{i=1}^{k} a_{i}  \, + \, a_{k+1}\right) = \sum_{i=1}^{k+1}(ca_{i})$\\
&$\ra k+ 1\in \zp$ and $c\sum_{i=1}^{k+1}a_{i} = \sum_{i=1}^{k+1}(ca_{i})$\\
&$\ra k+1\in S$\\
\end{tabular}\\
By Th.3.4.2 we have $S = \zp$ and for each $n\in \zp ,$ $c\sum_{i=1}^{n}a_{i} = \sum_{i=1}^{n}(ca_{i}).$ \nl

\textbf{Example 3.4.3.} Show that if \of is an ordered field and $n\in \zp,$ then $\sum_{i=1}^{n}(2i) = n(n+1).$\nl

Solution. Suppose \of is an ordered field. \nl
Let $S=\{ n \mid n\in \zp$ and $\sum_{i=1}^{n}(2i) = n(n+1) \}$

\newpage
%page50

\begin{tabular}{r l}
(1) $1\in S$& since $1\in \zp$ and $2\cdot 1 =1(1+1)$\\
(2) $k\in S$&$\ra k\in \zp $ and $\sum_{i=1}^{k}(2i)=k(k+1) $\\
&$\ra k\in \zp $ and $\sum_{i=1}^{k}(2i) + 2(k+1)= k(k+1)+2(k+1) $\\
&$\ra k+1\in \zp$ and $ \sum_{i=1}^{k+1}(2i)=(k+1)(k+2)=(k+1)[(k+1)+1] $\\
&$\ra k+1\in S$\\
\end{tabular}


By Th 3.4.2, we have $S = \zp$ and for each $n\in \zp,$ $ \sum_{i=1}^{n}(2i)=n(n+1).$\nl

\textbf{Definition 3.4.15.} If \of is an ordered field. Let $Q$ denote $\{n \mid$ there exist elements $p$ and $q$ in $Z$ such that $q \neq 0$ and $n=p\div q\}$. The elements of $Q$ are called \textit{rational elements} of $F.$\nl

\textbf{Theorem 3.4.15.} If \of is an ordered field, then $(Q,+,\cdot,<)$ is an ordered field.\nl

\textbf{Theorem 3.4.16.} If \of  is an ordered field, $a$ is an infinite sequence of elements of $F,$ and $n\in \zp$, then the sum (product) of $a_{1} ,a_{2}, \,\, \dots , \,\, a_{n}$ is
independent of the way in which the parentheses are placed.\nl

\textbf{Theorem 3.4.17.} If \of is an ordered field, $a$ is an infinite sequence of elements of $F$, and $n\in \zp,$ then the sum (product) of $a_{1} ,a_{2}, \cdots ,a_{n} $ is independent of the order in which they are written. \nl

\textbf{Exercises}\nl

If \of is an ordered field, find the value of each of the following.\nl

1. $\sum_{i=1}^{4}(i^{2} - 2)$\nl

2. $\sum_{i=2}^{4}(i+1)\div i$\nl

Show that if \of is an ordered field, $a$ is an infinite sequence of elements of $F,$ and $n\in \zp ,$ then each of the following are true.

\newpage
%page51
\begin{large}
$3.  \sum_{i=1}^{n} i = [n((n+1))]\div 2$\nl

$4. \sum_{i=1}^{n} [i(i+1)]\div 2] = [n(n + 1)(n + 2)\div 6]$\nl

$5.  \sum_{i=1}^{n} 1\div [i(i + 1) = n\div (n+1)$\nl

$6.   \sum_{i=1}^{n} i^{2} =  [n(n+1)(2n+1)]\div 6 $\nl

$7.   \sum_{i=1}^{n} (3i - 2) = [n(3n - 1)]\div 2 $\nl

$8.   \abs{\sum_{i=1}^{n} a_{i}} \leq  \sum_{i=1}^{n}\abs{a_{i}}$\nl

$9.$   If $x,y \in F$ and $0 < x < y$, then $x^{n} < y^{n}.$\nl

$10.$ If $x,y \in F,$ then $\abs{x}^{n} = \abs{x^{n}}.$\nl

$11.$ Show that if $x,y\in P$ and $x < y,$ then $x^{n} < y^{n}.$\nl

$12.$ Prove 3.4.14(b).\nl

$13.$ Prove 3.4.14(c).\nl
\end{large}
\vfill
\textbf{3.5 Algebra of Functions}\nl

\lhead{Section 3.5}
\textbf{Definition 3.5.1.} If \of is an ordered field, $n\in \zp,$ and each of $f$ and $g$ is a function with range a subset of $F,$ then \nl

a) The \textit{sum} of $f$ and $g$ is denoted by $f + g$ and defined by $f +g = \{ (x,y) \mid x\in D_{f}\cap D_{g}$ and $y= f(x) + g(x)\}.$\nl

b) The \textit{negative} of $f$ is denoted by $-f$ and defined by $-f= \{(x,y) \mid x\in D_{f}$ and $y = -f(x)\}$.\nl

c) The \textit{difference} of $f$ and $g$ is denoted by $f - g$ and defined by $f - g = f + (-g).$ \nl

d) The \textit{product} of $f$ and $g$ is denoted by $f\cdot g$ and defined by $f\cdot g = \{(x,y) \mid x\in D_{f}\cap D_{g}$ and $y = f(x)\cdot g(x) \}. $\nl

e) The \textit{reciprocal} of $f$ is denoted by $1 \div f$ and defined by $1\div f =\{ (x,y) \mid x\in D_{f}, f(x) \neq 0,$ and $y = 1\div f(x) \}.$ \nl

f) The \textit{quotient} of $f$ by $g$ is denoted by $f\div g$ and defined by $f\div g =f\cdot (1 \div g).$ \nl

\newpage
%page52
g) The $n^{th} power$ of $f$ is denoted by $f^{n}$ and defined by $f^{n} = \{(x,y) \mid x\in D_{f}$ and $y=(f(x))^{n}\}.$ \nl

h) The \textit{composite} of $f$ by $g$ is denoted by $f\circ g$ and defined by $f\circ g = \{(x,y) \mid x\in D_{g}, g(x)\in D_{f},$ and $y = f(g(x)\}.$\nl

i) The \textit{absolute value} of $f$ is denoted by $\abs{f}$ and defined by $\abs{f} =\{(x,y) \mid x\in D_{f}$ and $y= \abs{f}\}.$\nl


\textbf{Example 3.5.1.} If \of is an ordered field, $f = \{(1,2), (2,-1), (3,4), (4,0) \}$ and
$g =\{ (-1,3), (0,2), (1,0), (2,1), (3,-1) \},$ then \nl

$ f + g = \{ (1,2), (2,0), (3,3) \}$\nl
$-f = \{ (1,-2), (2,1), (3,-4), (4,0) \}$\nl
$f - g = \{ (1,2), (2,-2), (3,5) \}$\nl
$f\cdot g = \{(1,0), (2,-1), (3,-4) \}$\nl
$1\div f =\{ (1,1/2), (2,-1), (3,1/4) \}$\nl
$f\div g = \{ (2,-1), (3,-4) \}$\nl
$f\circ g = \{(-1,4), (0,-1), (2,2)\}$\nl

\textbf{Example 3.5.2.} If \of is an ordered field, $f = \{(x,y) \mid x\in [0,2]$ and $y =x^{2} - 1$ and $g =\{(x,y) \mid x\in [1,3]$ and $y = x^{2} - x \}$, then \nl

$f + g = \{ (x,y) \mid x\in [1,2]$ and $y =2x^{2} - x - 1\}$\nl
$-f = \{(x,y) \mid x\in [0,2]$ and $y = -x^{2}+1 \}$\nl
$f - g = \{ (x,y) \mid x\in [1,2]$ and $y = x - 1\}$\nl
$f\cdot g= \{ (x,y) \mid x\in [1,2]$ and $y = x^{4} -x^{3} -x^{2} +x \} $\nl
$1\div f= \{ (x,y) \mid x\in ([ 0,2] -\{1\} )$ and $y= 1 \div (x^{2} - 1) \}$\nl
$f\div g = \{(x,y) \mid x\in (1,2$] and $y = (x+1)\div x\}.$\nl



\textbf{Theorem 3.5.1.} If \of is an ordered field, $n\in \zp,$ and each of $f$ and $g$ is a function with range a subset of $F$, then: \nl
a) $D_{f+g} = D{f}\cap D_{g}$ and for each $x \in D_{f+g} , (f+g)(x)=f(x)+g(x).$\nl
b) $D_{-f} = D_{f} $ and for each $x \in D_{f} , (-f)(x)=-f(x).$\nl
c) $D_{f-g} = D_{f}\cap D_{g} $ and for each $x \in  D_{f-g}, (f-g)(x)=f(x)-g(x).$\nl
d) $D_{f\cdot g} = D_{f}\cap D_{g}$ and for each $x$ in $D_{f\cdot g} , (f\cdot g)(x)=f(x)\cdot g(x).$\nl
e) $D_{1/f} = \{x \mid x\in D_{f}$ and $f(x) \neq 0\} $ and for each $x \in D_{1/f}, (1 \div f)(x)= 1
\div f(x)$\nl
f) $D{f/g} = \{x \mid x\in D_{f}\cap D_{g}$  and $g(x)\neq 0\}$ and for each $x$ in $D_{f/g},
(f/g)(x)=f(x)/g(x). $\nl
g) $D_{f^{n}} =D_{f}$ and for each $x$ in $D_{f^{n}}, f^{n}(x)=(f(x))^{n}.$\nl

\newpage
%Page53

h) $D_{f\circ g } = \{x \mid x\in  D_{g}$ and $g(x) \in D_{f}\}$ and for each $x$ in $D_{f\circ g} ,(f\circ g) (x) = f(g(x)).$ \nl

i) $D_{\abs{f}} = D_{f}$ and for each $x$ in $D_{\abs{f}} , \abs{f}(x)=\abs{f(x)} .$\nl

\textbf{Definition 3.5.2.} If $(F,+,\cdot)$ is a field and $c\in F,$ then $\{(x,y) \mid x\in F $ and $y = c\}$ is called a \textit{constant function} and denoted by $\underline{c}.$\nl

\textbf{Theorem 3.5.2.} If \fld is a field and each of $f, g,$ and $h$ is a function whose range is a subset of $F,$ then\nl
\begin{tabular}{l l l}
a) & $f+g=g+f$ & e) $f\cdot g=g\cdot f$ \\
b) &($f+g)+h=f+(g+h)$ &f) $(f\cdot g)\cdot h= f\cdot (g\cdot h$) \\
c) &$f + \underline{0} =f$ &g) $f\cdot \underline{1}  = f$ \\
d) &$f+(-f) = \underline{0}$ on $D_{f}$& h) $f\cdot (1\div f) =\underline{1}$ on $\{x \mid x\in D_{f}$and $f(x)\neq 0\}$ \\
i) &$f\cdot (g+h)=f\cdot g + f\cdot h$&\\
\end{tabular}\\


Proof of (b) and (c). \fld is a field and each of $f, g,$ and $h$ is a function whose range is a subset of $F.$\nl
    Proof of Assertion b)\nl
\begin{tabular}{l l }
$(f+g)+h$ &$= \{(x,y) \mid  x\in D_{f+g}\cap D_{h}$ and $y =(f+g)(x)+h(x) \}$\\
&$=\{(x,y) \mid x\in (D_{f}\cap Dg)\cap D_{h}$ and $y=(f(x)+g(x))+h(x)\}$\\
&$= \{(x,y) \mid x\in D_{f}\cap (D_{g}\cap D_{h})$ and $y = f(x)+(g(x)+h(x))\} $\\
&$=\{(x,y) \mid x\in D_{f}\cap D_{g+h}$ and $y=f(x)+(g+h)(x) \}$\\
&$=f+(g+h)$ \\
\end{tabular}\\
Proof of Assertion c) \nl
\begin{tabular}{l l }
$f + \underline{0}$ & $= \{(x,y) \mid x\in D_{f}\cap D_{\underline{0}}$ and $y=f(x) + \underline{0}(x) \} $\\
&$= \{(x,y) \mid x\in D_{f}\cap F$ and $y= f(x) + 0\} $\\
&$= \{(x,y) \mid x\in D_{f}$ and $y=f(x) \} $\\
&$=f$ \\
\end{tabular}\nl

\textbf{Exercises}\nl

1. 	If \fld is a field, $f =\{(-1,2), (0,3), (1,-4), (2,-2)\}$ and $g = \{ (0,2), (1,3), (2,-1), (3,1)\},$ find $f+g, -f, f-g, f\cdot g, 1\div f, f\div g,$ and $f\circ g.$\nl

2. 	If \of is an ordered field, $f = \{(x,y) \mid x\in  [0,2]$ and $y = x - 1\},$ and $g= \{(x,y) \mid  x\in  [-1,1]$ and $y = x^{2} - 1,$ find $f+g, -f, f-g, f\cdot g, 1\div f,$ and $f\div g.$ \nl

3. 	Prove Theorem 3.5.2. (a) and (d). \nl
4. 	Prove Theorem 3.5.2. (e) and (f). \nl
5. 	Prove Theorem 3.5.2. (g) and (h). \nl
6. 	Prove Theorem 3.5.2. (i). \nl


\textbf{Definition 3.5.3.} If \fld~is a field, $\{(x,y) \mid x\in F$ and $y = x\}$ is called the \textit{identity function} and denoted by $I.$\nl

\newpage
%Page54

\textbf{Theorem 3.5.3}  If \of is an ordered field, $c \in F, n \in \zp$,  and each of $f, h $ and $h$ is a function whose range is a subset of $F$, then\nl
a) $( f \pm g) \circ h = f\circ h \pm g\circ h$\nl
b) $(f\cdot g)\circ h = (f\circ h)\cdot (g\circ h)$\nl
c) $(f/g)\circ h = f\circ h/g\circ h$\nl
d) $(f\circ g)\circ h = f\circ (g \circ h)$\nl
e) $f^{n}\circ g = (f\circ g)^{n}$\nl
f) $I\circ f = f\circ I = f$\nl
g) $I^{n}\circ g = g^{n}$\nl
h) $\underline{c}\circ f = \underline{c}$ on $D_{f}$\nl
i) $f\circ \underline{c} = \underline{f(c)}$\nl

Proof of (a).   Suppose \fld is a field, $c \in F$, and each of $f, g,$ and $h$ is a function whose range is a subset of $F$.\nl
%\begin{small}

$(f + g)\circ h$\nl
$= \{ (x,y) \mid x \in D_{h}, h(x) \in D_{f+g},$ and $y = (f+g)(h(x))\}$\nl
$= \{ (x,y) \mid x \in D_{h},h(x)\in D_{f}\cap D_{g}$ and $y = f(h(x)) + g(h(x))  \}$\nl
$= \{ (x,y) \mid x\in D_{h}, h(x) \in D_{f}, h(x)\in D_{g},$ and $ y=(f\circ h)(x) + (g\circ h)(x)\}$\nl
$= \{ (x,y) \mid (x\in D_{h}$ and  $h(x) \in D_{f})$ , $(x \in D_{h})$ and $(h(x) \in D_{g}),$ 
 and $y = (f\circ h)(x) + (g\circ h)(x) \}$\nl
$= \{ (x,y) \mid x\in D_{f\circ h}, x\in D_{g\circ h},$ and $y = (f\circ h)(x) + (g\circ h)(x) \}$\nl
$= \{ (x,y) \mid x\in D_{f\circ h} \cap D_{g\circ h},$ and  $y = (f\circ h)(x) + (g\circ h)(x) \}$\nl
$= f\circ h + g\circ h$\nl
%\end{small}

\textbf{Example 3.5.3}  If \of is an ordered field, $f = \{ (x,y) \mid x\in [-2, 4]$ and $y = x^{2} - 1 \}$, and
$g = \{ (x,y) \mid x \in [0, 1) $ and $y = (x + 1)\div (x - 1)\}$, then \nl
%\begin{small}
$f\circ g$ \nl
$= \{ (x, y) \mid x\in D_{g}, g(x) \in D_{f},$~and~$y = f(g(x)) \}$\nl
$=\{ (x, y) \mid x\in [0,1) , (x+1)\div (x-1)  \in [-2, 4],\text{~and~}y = f((x+1)\div(x-1))\}$\nl
$= \{ (x, y) \mid x\in [0, 1), -2\leq [(x+1)\div (x-1)]\leq 4$ and \nl
$ $ \hfill $y = [(x+1)\div (x-1)]^{2}- 1\}$\nl
$= \{ (x, y) \mid x\in [0, 1),-2x + 2 \geq x + 1, x+ 1 \geq 4x - 4,\text{ ~and~}y = (4x)\div (x-1)^{2} \}$\nl
$= \{ (x, y) \mid 0 \leq x < 1, 1\div 3 \geq x, 5\div 3 \geq x\text{~and~} y = (4x)\div (x-1)^{2} \}$\nl
$= \{ (x, y) \mid 0 \leq x \leq 1\div 3\text{~and~}y = (4x)\div (x-1)^{2} \}$\nl
$= \{ (x, y) \mid x\in [0, 1/3]\text{~and~}y = (4x)\div (x-1)^{2} \}$\nl
%\end{small}





\newpage

\textbf{Exercises}\nl

Assuming \of is an ordered field, find $f\circ g$ for each of the following definitions of $f$ and $g$.\nl
%page55
\begin{tabular}{l l}
1.& $f = \{ (x,y) \mid x\in [-1, 1] $ and $y = x^{2} - 3 \}$\\
&$g= \{(x,y) \mid x\in [0,2]$ and $y = x -2\}$\\
2.& $f= \{(x,y) \mid x\in F$ and $y =1\div (x^{2} + 1)\}$ \\
&$g= \{(x,y) \mid x\in  [0,5]$ and $y = 2x - 1\}$\\
3. &	$f = \{(x,y)\mid x\in F$ and $y =x^{2}-2x-1 \}$\\
&$g =\{(x,y) \mid x\in [0,3]$ and $y = 2\}$\\

4. &$f=\{(x,y) \mid x\in [0,2]$ and $y = 1\div (x+1)\} $\\
&$g= \{(x,y) \mid x\in (1,2]$ and $y =1 \div (x-1)\}$\\

5. 	&$f= \{(x,y) \mid x\in (1,2)$ and $y=(x^2 + 1)\div (x-1)\}$\\
&$g= \{(x,y) \mid x\in (-\infty,3]$ and $y = x + 1\} $\\

6.& 	Prove Theorem 3.5.3. (b) and (c) \\

7.& 	Prove Theorem 3.5.3. (d) and (f) \\

8.& 	Prove Theorem 3.5.3. (e) and (g) \\
\end{tabular}\nl
\vfill
\textbf{3.6 Inverse Functions}\nl
\lhead{Section 3.6}

\textbf{Definition 3.6.1.} If $f$ is a one-to-one function, then $\{(x,y) \mid (y,x) \in f\}$ is called the \textit{inverse} of $f$ and denoted by $f^{-1}$ , i.e., $y= f^{-1}(x) \lra x =f(y).$\nl

\textbf{Theorem 3.6.1.} If $f$ is a one-to-one function, then\nl
a) $D_f  =R_{f^{-1}}$\nl
b) $R_{f} =D_{f^{-1}} $\nl
c) $f(f^{-1}(x)) = x$ for each $x$ in $D_{f^{-1}}$\nl
d) $f^{-1}(f(x)) = x$ for each $x$ in $D_{f}$\nl
Proof of (b) and (c).  Suppose $f$ is a one-to-one function. \nl
b) Suppose $f$ is a one-to-one function and $x\in D_{f^{-1}}$. There exists a $y$ such that $(x,y)\in f^{-1}.$ Hence $(y,x)\in f.$ Thus $x\in R_{f}.$\nl
c) Suppose $f$ is a one-to-one function and $x\in D_{f^{-1}}.$ There exists a $y$ such that $(x,y)\in f^{-1}.$  Hence $y = f^{-1}(x)$ and $x = f(y).$ Thus $f(f^{-1}(x))=x.$\nl

\newpage
%page56



\textbf{Section 3.6 }\nl

\textbf{Definition 3.6.2.} If \of  is an ordered field, $f$ is a function from a subset of $F$ into $F,$ and

$A\subseteq D_{f},$ then $f$ is said to be \textit{increasing}  \underline{\textit{( decreasing )}} on $A$ if and only if for each $x$ and $y$ in $A,~x < y$ implies $f(x) < f(y)$ \underline{$(f(x) >f(y))$}. A function is said to be increasing if and only if it is increasing on its domain. \nl

\textbf{Theorem 3.6.2.} If \of  is an ordered field and $f $ is an increasing
\underline{(decreasing)} function from a subset of $F$ into $F$,  then $f^{-1}$ is increasing \underline{(decreasing)}. \nl

Proof. Suppose \of is an ordered field and $f$ is an increasing function from a subset of $F$ into
$F.$  Suppose  $x,y \in D_{f}, x<y,$ and $f^{-1}(x) \geq f^{-1}(y).$
Since $f$ is increasing, we have $f( f^{-1}(x)) \geq f( f^{-1}(y) )$. Thus $x \geq y,$ which is
a contradiction. \nl


\textbf{Example 3.6.2.} Assuming \of  is an ordered field and $f$ is the function with domain $[1,2]$
such that $f(x) =(1\div x)- 1$ for each $x$ in $[1,2],$ find $D_{f^{-1}}$ and an expression for $f^{-1}(x).$\nl Solution.  Suppose \of  is an ordered field and $f$ is the function with domain $[1,2]$ such that $f(x) =(1\div x)- 1$ for each $x$ in $[1,2].$ Since $f$ is de­creasing,  $R_{f} = [f(2) ,f(1)] = [-1/2, 0].$ Suppose $x \in D_{f^{-1}}.$  By Definition 3.6.1, $x = f(y).$  Thus\nl

$x = 1/y -1$\nl
$x + 1= 1/y$\nl
$y = 1/(x + 1) $\nl
$f^{-1}(x) = 1/(x + 1)$\nl

\textbf{Exercises}\nl

Assuming \of is an ordered field, find $D_{f^{-1}}$ and an expression for $f^{-1}(x)$ for each of the following definitions of $f.$\nl

\begin{tabular}{l l}
1. $f(x)=2x-3$ & $ D_{f} = [-1,2]$\\
2. $f(x)=1 - 4x$& $ D_{f} = [0,2] $\\
3. $f(x)=1\div(x-1)$& $ D_{f}=[-1,0]$ \\
4. $f(x)=x\div (1-x)$ & $D_{f} =[2,4]$\\
5. $f(x)=(x+1)\div(x-1)$ &  $D_{f} = [2,3] $\\
%6. Which of the following are one-to-one functions? &  \\
\end{tabular}\\
~~6. Which of the following are one-to-one functions?
\newpage

%Page 57

a) $\{(1,2), (2,4), (5,3), (3,1)\}$\nl
b) $\{(1,3), (4,2), (1,5), (2,1)\}$\nl
c) $\{((1,1),2), ((1,2),3), ((2,1),3), ((2,2),4)\}$\nl
d) $\{((1,1),1), ((1,2),2), ((2,3),6)\}$\nl

7. Assuming \of  is an ordered field, determine which of the following are one-to-one functions.\nl

\begin{tabular}{l l}
a) $f(x)=x^{2} - 2$& $D_{f} = [-1,1]$\\
b) $f(x)=1\div(x+1)$ & $D_{f} = [1,3] $\\
c) $f(x)=x^{2} + 4$ & $D_{f} = [0,2] $\\
d) $f(x)=2x - 3$& $D_{f} = F$\\
\end{tabular}\\
$ $\nl

8. Prove Theorem 3.6.1(a). \nl

9. Prove Theorem 3.6.1(d). \nl

\vfill

\lhead{Section 3.7}
\textbf{3.7 Arithmetic} \nl

\textbf{Definition 3.7.1.} If \of  is an ordered field, $m,n\in Z, m\leq 0 \leq n,$  and $d : Z_{m}^{n} \ra Z^{9}_{0},  $ then $\sum_{i=m}^{n} d_{i}10^{i}$ will be denoted by setting the ordinates  of $d$ side by side in decreasing order of subscript with a period (called a decimal point) following $d_{0}.$\nl

\textbf{Example 3.7.1.} If $d_{-1} = 5, d_{0} = 1, d_{1} = 7,$ and $d_{2} = 3,$ then 371.5 denotes
\[
\sum_{i=-1}^{2} d_{i}10^{i} =3\cdot 10^{2} + 7\cdot 10^{1} + 1\cdot 10^{0}+ 5\cdot 10^{-1}.
\]

\textbf{Theorem 3.7.1.} If \of  is an ordered field, then

\newpage
%Page58
\begin{center}

\qqu  Table 1\nl

\begin{tabular}{|r| r| r| r| r| r| r| r| r| r| r |}\hline
+& 0&1& 2& 3& 4& 5& 6& 7& 8& 9  \\ \hline
0& 0&1& 2& 3& 4& 5& 6& 7& 8& 9 \\ \hline
1&1&2&3&4&5&6&7&8&9&10\\ \hline
2&2&3&4&5&6&7&8&9&10&11 \\ \hline
3&3&4&5&6&7&8&9&10&11&12 \\ \hline
4&4&5&6&7&8&9&10&11&12&13  \\ \hline
5&5&6&7&8&9&10&11&12&13&14 \\ \hline
6&6&7&8&9&10&11&12&13&14&15 \\ \hline
7&7&8&9&10&11&12&13&14&15&16 \\ \hline
8&8&9&10&11&12&13&14&15&16&17 \\ \hline
9&9&10&11&12&13&14&15&16&17&18 \\ \hline
\end{tabular}\nl

$ $\nl

\qqu  Table 2\nl

\begin{tabular}{|r| r| r| r| r| r| r| r| r| r| r |}\hline
$\cdot$& 0&1&2&3&4&5&6&7&8&9 \\ \hline
0&0&0&0&0&0&0&0&0&0&0 \\ \hline
1&0&1&2&3&4&5&6&7&8&9 \\ \hline
2&0&2&4&6&8&10&12&14&16&18 \\ \hline
3&0&3&6&9&12&15&18&21&24&27 \\ \hline
4&0&4&8&12&16&20&24&28&32&36 \\ \hline
5&0&5&10&15&20&25&30&35&40&45 \\ \hline
6&0&6&12&18&24&30&36&42&48&54 \\ \hline
7&0&7&14&21&28&35&42&49&56&63 \\ \hline
8&0&8&16&24&32&40&48&56&64&72 \\ \hline
9&0&9&18&27&36&45&54&63&72&81 \\ \hline
\end{tabular}\nl

\end{center}

Proof. Suppose \of is an ordered field. The procedure is to prove the entries in each table in row order. Suppose all the entries in Table 1 have been filled in up to but not including $4 + 5.$ Then
$4 + 5 = 4 +(4+1) = (4+4) + 1 = 8 + 1 = 9.$ The table can, be completed in the same manner. \nl
Suppose the entries in Table 2 have been filled in up to $3\cdot 6.$ Then\nl
$3\cdot 6 =3\cdot(5+1) =3\cdot 5 + 3\cdot 1 = 15 + 3 = (10 + 5) + 3 = 10 + (5+3)= 10 + 8 = 18.$ \nl
The table can be completed by continuing the process. \nl

\textbf{Example 3.7.2.} Find $3 - 8.$ \nl
Solution. Let $x = 8 -3.$ Thus $x + 3 = 8.$ By Table 1, we see that $x = 5.$ Hence $3 - 8 = -(8-3) = -5.$\nl

\vfill
\textbf{3.8 The Binomial Theorem}
\lhead{Section 3.8}\nl

\textbf{Definition 3.8.1.} If \of  is an ordered field and $a$ is an infinite
sequence of elements of $F$, let $\prod_{i=1}^{1}a_{i}$ denote $a_{1}$ and for each $n$ in $\zp$,   \nl

\newpage
%page59 continue sentence from prior page

let $\prod_{i=1}^{n+1} a_{i} $ denote $\left(\prod_{i=1}^{n}a_{i}\right)a_{n+1} .$ \nl

\textbf{Definition 3.8.2.} If \of is an ordered field, then for each $n$ in $\zp, \prod_{i=1}^{n} i $ is called $n$ \textit{factorial} and is denoted by $n!.$  $0!$ is called $0$ factorial and is defined to be $1$. \nl

\textbf{Definition 3.8.3.} If \of is an ordered field and $k \in Z_{0}^{n},$ let $\binom{n}{k} $ denote
$n!\div[(n-k)!(k!)].$\nl

\textbf{Theorem 3.8.1.}  If \of is an ordered field and $n\in \zp$, then\nl
a) $(n+1)! = n!(n+1) $\nl

b) $\binom{n}{0} = \binom{n}{n}$\nl

c) $\binom{n}{1} = \binom{n}{n-1} = n$\nl

\textbf{Theorem 3.8.2.}  If \of is an ordeed field , $n \in \zp,$ and $k\in Z_{1}^{n},$ then\nl
\[
\binom{n}{k} + \binom{n}{n-1} = \binom{n+1}{k}.
\]
Proof.\nl
\begin{tabular}{r l}
$\binom{n}{k}+\binom{n}{k-1}$&$=[n!\div((n-k)!k!)] + [n!\div((n-k+1)!(k-1)!)]$\\
&$=[(n-k+1)n! +k(n!)]\div[(n-k+1)!k!]$\\
&$=[n!(n+1)]\div[(n-k+1)!k!]$\\
&$=\binom{n+1}{k}$\\
\end{tabular}\\
$ $\nl
\textbf{Theorem 3.8.3.}  (The Binomial Theorem)  If \of is an ordered field, $a, b \in F,$ and $n\in \zp$, then
\[
\ (a + b)^{n} = \sum_{k=0}^{n}\binom{n}{k}a^{n-k}b^{k}.
\]
Proof.  Suppose \of is an ordered field and $a,b\in F.$\nl
Let $S = \{ n \mid n\in \zp $ and  $(a + b)^{n} = \sum_{k=0}^{n}\binom{n}{k}a^{n-k}b^{k} \}$\nl
(1) $1\in S$ since $1\in \zp$ and $(a+b)^{1} = a + b$\nl

\begin{tabular}{r l}
(2) $n\in S$&$\ra n\in \zp$ and $(a + b)^{n} = \sum_{k=0}^{n}\binom{n}{k}a^{n-k}b^{k} $\\
&\\
&$\ra n+1\in \zp$ and $(a+b)^{n+1} = (a+b)(a+b)^{n} $ \\
&\\
&$= a(a+b)^{n}  +b(a+b)^{n}$\\
&\\
&$= a\sum_{k=0}^{n}\binom{n}{k}a^{n-k}b^{k} +b\sum_{k=0}^{n}\binom{n}{k}a^{n-k}b^{k}$\\
&\\
&$=\sum_{k=0}^{n}\binom{n}{k}a^{n-k+1}b^{k} + \sum_{k=0}^{n}\binom{n}{k}a^{n-k}b^{k+1}$\\
&\\
&$= \sum_{k=0}^{n}\binom{n}{k}a^{n-k+1}b^{k} + \sum_{k=1}^{n+1}\binom{n}{k-1}a^{n-k+1}b^{k}$\\
\end{tabular}

\newpage
%page60
\begin{tabular}{r l}
&$= a^{n+1} + \sum_{k=1}^{n} \binom{n}{k} a^{n-k+1}b^{k} + \sum_{k=1}^{n} \binom{n}{k-1} a^{n-k+1}b^{k} + b^{n+1}$\\
&\\
\hspace{20pt}& $= a^{n+1} + b^{n+1}+\sum_{k=1}^{n}\binom{n+1}{k}a^{n-k+1}b^{k}$\\
&\\
\hspace{20pt}&$= \sum_{k=0}^{n+1}\binom{n+1}{k}a^{(n+1) - k}b^{k}$\\
&\\
\hspace{20pt}&$\ra n+1 \in S$ \\
&\\
\end{tabular}\\

By Th 3.4.2., we have $S = \zp$ and for each $n\in \zp,\, (a + b)^{n} = \sum_{k=0}^{n}\binom{n}{k}a^{n-k}b^{k}.$ \nl

\textbf{Theorem 3.8.4.} If \of is an ordered field, $a,b\in F,$ and $n\in \zp,$ then (1) $a^{n}b^{0}$ is the first term in the binomial expansion of $(a+b)^n$ and
(2) if $ca^{i}b^{j}$ is any term in the binomial expansion of $(a+b)^{n},$ then $((ci) \div (j+i))a^{i-1}b^{j+1}$ is the next term in the expansion.\nl
 Proof. Suppose \of is an ordered field, $a,b\in F,$ and $n\in \zp.$ Each term in the binomial expansion of $(a+b)^{n}$ may be expressed as $\binom{n}{k}a^{n-k}b^{k}$ for some $k$ in $Z_{0}^{n}.$ By letting $c=  \binom{n}{k}, i = n - k,$ and $j = k,$we see that\nl 
 $ca^{i}b^{j}= \binom{n}{k}a^{n-k}b^{k}.$  Therefore\nl
\begin{tabular}{r l}
&\\
$((ci) \div (j+1))a^{i-1}b^{j+1}$&$=\binom{n}{k} ((n-k) \div (k+1))a^{n-k-1}b^{k+1}$\\
&\\
& $= [(n!(n-k))\div ((n-k)!k!(k+1))]a^{n-(k+l)}b^{k+1}$\\
&\\
& $=[n! \div((n-(k+1))!(k+1)!)]a^{n-(k+1)}b^{k+1} \text{~if~} k\in Z_{0}^{n-1}$\\
&\\
or&$= 0$ if $k = n,$ \\
\end{tabular}\\
 which is the next term in the expansion. \nl

\textbf{Example 3.8.1.} Assuming \of is an ordered field and $x,y\in F,$ find the binomial expansion of
$(2x-y^{2})^{4}.$  \nl
 Solution. Suppose \of is an ordered field, and $x,y\in F.$ From Th 3.8.4, we see that \nl
\begin{tabular}{r l}
$(2x+(-y^{2}))^{4}$ & $=1\cdot(2x)^{4}(-y^{2})^{0} + [(1\cdot 4)\div 1](2x)^{3}(-y^{2})^{1}  + [(4.3)\div 2] (2x)^{2} (-y^{2})^{2}$\\
&$+[(6\cdot 2)\div 3](2x)^{1} (-y^{2})^{3} + [(4\cdot 1)\div 4] (2x)^{0}(-y^{2})^{4} $\\
&$=16x^{4} - 32x^{3}y^{2} + 24x^{2}y^{4} -8xy^{6} + y^{8}.$\\
\end{tabular}\nl

\textbf{Example 3.8.2.} Assuming \of is an ordered field and $x,y\in F,$ find the seventh term of $(2x - y^{-1})^{10}.$\nl
 Solution. Suppose \of is an ordered field and $x,y\in F.$ The seventh term of $(2x - y^{-1})^{10}$ is
 $\binom{10}{6}(2x)^{10-4}(-y^{-1})^{6}.$\nl
\begin{tabular}{r l}
$\binom{10}{6}(2x)^{10-4}(-y^{-1})^{6}$&$= (10!\div (4!6!))(2x)^{4}y^{-6}$\\
&$= [(1\cdot 2\cdot 3\cdot 4\cdot 5\cdot 6\cdot 7\cdot 8\cdot 9\cdot 10)$\\
&$ \div  (1\cdot 2\cdot 3\cdot 4\cdot 5\cdot 6\cdot 1\cdot 2\cdot 3\cdot 4)] 16x^{4}y^{-6}$\\
&$=210(16)x^{4}y^{-6}$\\
&$=3360x^{4}y^{-6}$\\
\end{tabular}\\

\newpage
%page61


\textbf{Exercises} \nl

Assuming \of is an ordered field and $a,b,x,y\in F,$ find the binomial expansion of
each of the following. \nl

1. $(a+b)^{6}$ \nl

2. $(2x-y)^{4}$ \nl

3. $(x+y^{2})^{5}$ \nl

4. $(x^{-1}+2{y})^{4}$ \nl


5. $(x^{1/2} +y^{1/2})^{4}$ \nl

6. $(x^{-2} - 2y)^{5}$\nl

7. $(x^{3/4} - 2y^{-1})^{4}$ \nl

Find the indicated term in the binomial expansion of each of the following. \nl

8. Sixth term of $(2x - y^{-2})^{12} $\nl

9. Ninth term of $(2 + (x\div 2))^{15}$\nl

10. Fifth term of $(x + (1\div x))^{10}$\nl

11. Seventh term of $(\sqrt{x} - (1\div 2\sqrt{x}))^{14}$\nl

12. Prove Theorem 3.8.1. \nl

%\include{Roach}

%\documentclass[12pt]{amsart}
%\documentclass{report}
%\usepackage{amssymb,latexsym}
%\usepackage{amsmath}
%\usepackage{amsfonts}
%\usepackage{graphics,graphicx}
\setlength{\topmargin}{-10pt}
\setlength{\voffset}{-40pt}
\setlength{\textheight}{700pt}
%\usepackage{amsscr}
%\usepackage[amsscr]{eucal}

%\newtheorem{theorem}{Theorem}
%\newtheorem{lemma}{Lemma}
%\newtheorem{proposition}{Proposition}
%\newtheorem{definition}{Definition}
%\newtheorem{corollary}{Corollary}
%\newtheorem{notation}{Notation}
%\newtheorem{remark}{Remark}
%\newtheorem{assertion}{Assertion}
%\newtheorem{claim}{Claim}

%\newcommand{\abs}[1]{\left|#1\right|}
%\newcommand{\norm}[1]{\left\|#1\right\|}
%\newcommand{\p}[1]{\left(#1\right)}
%\newcommand{\tail}[1]{#1^{\p{N}}}
%\newcommand{\trunk}[1]{#1^{\left[N\right]}}
%\newcommand{\dg}{^{\circ}}
%\newcommand{\nl}{\newline}
%\newcommand{\fl}{\flushleft}
%\newcommand{\tild}{\texttt{\symbol{126}}}
%\newcommand{\tild}{\sim}
%\newcommand{\up}{\wedge}
%\newcommand{\dn}{\lor}
%\newcommand{\lra}{\leftrightarrow}
%\newcommand{\ra}{\to}
%\newcommand{\qqu}{\qquad \qquad}
%\newcommand{\onto}{\dashv}
%\newcommand{\power}{\Pi}
%\newcommand{\univ}{\textbf{U}}
%\newcommand{\ns}{\emptyset}
%\newcommand{\of}{$(F, +, \cdot , <)\,$}
%\newcommand{\zp}{Z^{+}}
%\newcommand{\fld}{ $(F,+,\cdot)$}
%\newcommand{\po}[1]{\dfrac{\pi}{#1}}
%\newcommand{\lpo}[1]{\frac{\pi}{#1}}
%\newcommand{\q}{\indent}

%\pagestyle{plain}


%\begin{document}
\setcounter{page}{62}
\lhead{Section 4.1}
%\begin{large}
\textbf{Chapter 4}\nl

\textbf{The Real Number System} \nl

It is not possible to show that the basic set of an ordered field con­tains an element $x$
such that $x^{2} = 2.$ We, therefore, add the completeness axiom in order to completely
characterize the real number system. \nl

\textbf{4.1 The Completeness Axiom}\nl

\textbf{Definition 4.1.1.} If \of is an ordered field, $b\in F,$ and $\ns \neq S\subseteq F,$ then:\nl
a) $b$ is said to be an \textit{upper bound} of $S$ if and only if $x\leq b$ for each $x$ in $S.$\nl
b) b is said to be a \textit{lower bound} of $S$ if and only if $b\leq x$ for each $x$ in $S.$\nl
c) $S$ is said to be \textit{bounded} if and only if $S$ has an upper bound and a lower bound.\nl

\textbf{Definition 4.1.2.} If \of is an ordered field, $b\in F,$ and $\ns \neq S\subseteq F,$ then: \nl
a) $b$ is called the \textit{least upper bound} of $S,$ and denoted by \textit{lub} $S$, if and
only if
(1) $b$ is an upper bound of $S$ and (2) if $c$ is an upper bound of $S,$ then $b\leq c.$\nl
b) $b$ is called the \textit{greatest lower bound} of $S$, and denoted by \textit{glb} $S$, if and only if (1) $b$ is a lower
bound of $S$ and (2) if $c$ is a lower bound of $S,$ then $c\leq b.$\nl


\textbf{Definition 4.1.3.} $(R,+,\cdot ,<)$ is called a \textit{complete ordered field} if and only if $(R,+,\cdot ,<)$ is an ordered field and the following statement is true. \nl
C. Every nonempty subset of $R$ which has an upper bound also has a least upper bound. \nl

\parbox{5in}{\textbf{Definition 4.1.4.} If $(R,+,\cdot ,<)$ is a complete ordered field, then $(R,+,\cdot,<)$} is called the \textit{real number system} and the elements of $R$ are called \textit{numbers.} $R, \zp, Z,$ and $Q$ will be
assumed to belong to the real number system un­less otherwise stated.\nl

\textbf{Example 4.1.1.}\nl

a) glb $(0, 1) = 0.$ \nl

b) glb $[2,3] = 2$\nl

c) glb $\{1, \frac{1}{2}, \frac{1}{3}, \dots \} = 0$\nl


\newpage
%page63

$ $\nl
d) $(-1, \infty)$ has no least upper bound. \nl
e) lub $\{1, 2, 3 \} = 3$\nl
Proof of (a). 0 is a lower bound of $(0,1)$ since $0 < x < 1$ for each $x$ in $(0,1).$
Assume there is a lower bound $c$ of $(0,1)$ greater than $0.$\nl

Case 1. $c<1.$ Since $0 <\frac{c}{2} < c <1,$ then $\frac{c}{2}\in (0,1)$ and $\frac{c}{2}$ is less than the lower bound $c$ of $(0,1).$ This is a contradiction. \nl

Case 2. $c\geq 1.$ Since $0 <\frac{1}{2} <1 \leq c,$ then $\frac{1}{2}\in (0,1)$ and $\frac{1}{2}$ is less than the lower bound $c$ of $(0,1).$   This is a contradiction. \nl

\textbf{Theorem 4.1.1.} If $S$ is a number set which has a lower bound, then $S$ has a greatest
lower bound.\nl
Proof. Suppose $S$ is a nonempty number set which has a lower bound. Let
$T = \{x \mid x $ is a lower bound of $S\}$. Since $T$ is bounded above by each element in $S, T$ has a
least upper bound, say $b.$ Suppose $b$ is not a lower bound of $S.$ Thus there is a $x$ in $S$ such
that $x < b.$ This means $x$ is a upper bound of $T$ less than $b$, which is a contradiction.  Suppose $b$
is not the greatest lower bound of $S.$  Thus there is a lower bound $c$ of $S$ such that $c > b.$  Since
$c\in T, b$ is not an upper bound of $T,$ which is a contradiction.  Therefore $b =$ glb $S.$\nl

\textbf{Theorem 4.1.2.} $\zp$  does not have an upper bound.\nl
 Proof. Suppose $\zp$ has an upper bound. By $C$ of Definition 4.1.3, $\zp$ has a least upper bound, say $b.$ Thus $n\leq b$ for each
$n$ in $\zp$. If $n$ in $\zp,$ then $n+1\in \zp$ and $n+1 \leq b.$ Hence $n \leq  b-1,$ for each $n$ in $\zp,$ which contradicts the fact that $b$ is the least upper bound of $\zp.$ \nl


\textbf{Theorem 4.1.3.} If $a\in R$ and $a > 0,$ then there is a $n$ in $\zp$ such that $\frac{1}{n} < a.$\nl

\textbf{Theorem 4.1.4.} If $a,b\in R$ and $a < b,$ then there is a rational number $r$ such that $a < r < b.$\nl

Proof. Suppose $a$ and $b$ are numbers and $a < b.$  $\frac{1}{(b - a)}$ is less that some positive integer.
Let $k$ denote a positive integer such that $\frac{1}{(b - a)} < k.$  Thus, $\frac{1}{(b - a)} < k$ and $\frac{1}{k} < (b - a).$  It follows from Theorem 3.4.11 that there is a least integer $c$ greater than $k*a.$  Thus, $k*a < c \leq k*a + 1$ and $a < \frac{c}{k} \leq a + \frac{1}{k} < a + b - a = b.$  Thus $a < \frac{c}{k}
< b.$  $\frac{c}{k} $ is a rational number between $a$ and $b.$\nl

\newpage
%page64


Thus $m - 1 \leq un < m.$ Since $\frac{1}{v - u} < n$, we have $1 + un < vn.$ Hence $un < m < un + 1 < vn$ and
$u < \frac{m}{n} < v.$  Therefore, $a + h <\frac{m}{n} < b + h$ and $a < \frac{m - nh}{n} < b.$ By Th 3.4.7 and Th 3.4.10, $\frac{m +(-1)nh}{n} $ is rational.\nl

\textbf{Definition 4.1.5.} A number which is not rational is called \textit{irrational.}\nl

\textbf{Exercises }
State whether each of the following sets is bounded, bounded above, or bounded below.
In each case give the least upper bound and the greatest lower bound where possible.\nl

1. $[2,4)$\nl

2. $\{\frac{1}{2}, \frac{3}{4},\frac{7}{8},\frac{15}{16},\frac{31}{32},\dots \}$\nl

3. $\{n \mid n$ is an integer $\} $\nl

4. $(-\infty,1)$\nl

5. $\{0, \frac{1}{2}, -1,4 \}$ \nl

6. $\zp$\nl

7. Prove that glb $[0,1]=0.$ \nl

8. Prove that glb $\{ \frac{1}{n} \mid n\in \zp\} = 0.$ (Hint: Use Th 4.1.3.)\nl

9. Prove that if $a\in R, a > 0,$ and $b\in R,$ then there is an element $n$ in $\zp$ such that $na > b.$\nl

10. Prove that the sum of a rational number and an irrational number is an irrational number. (Hint: Use Th 3.4.15.) \nl

11. Prove that the product of a rational number and an irrational num­ber is an irrational number. \nl

12. Show that if $a,b,c\in R, a < b,$ and $c$ is irrational, then there is an irra­tional number $t$ such
that $a < t < b.$ (Hint: Apply Th 4.1.4 to $a - c < b - c.$) \nl

13. Prove Theorem 4.1.3. \nl

\newpage
%page65

\lhead{Section 4.2}
\textbf{4.2 Logarithmic Functions}\nl

\textbf{Definition 4.2.1.} Let $P = \{x \mid \in R$ and $x >0\}.$\nl

 The proof of the following theorem is given in Appendix I.\nl


\textbf{Theorem 4.2.1.} There is a function $L$ from $P$ into $R$ such that if $x,y\in P,$ then \nl
Ll. $L(xy)=L(x)+L(y)$\nl
L2. $L(x)\leq x - 1$\nl
$L$ is called the natural logarithm or logarithm to the base $e$ where $L(e) = 1.$\nl

\textbf{Theorem 4.2.2.} If $x,y\in P,$ then \nl
\begin{tabular}{l l}
a) $L(1) = 0$ & c) $L(\frac{1}{y}) = -L(y)$\\
b) $L(\frac{x}{y}) =L(x)-L(y)$& d) $ 1 - \frac{1}{x} \leq L(x)$\\
\end{tabular}\\
Proof of (b). Suppose $x,y\in P.$ Thus $\frac{x}{y}\in P$ and $L(x) = L(y\frac{x}{y}) = L(y)+L(\frac{x}{y}). $
Hence $L(\frac{x}{y}) =L(x) - L(y).$\nl

\textbf{Theorem 4.2.3.} $L$ is increasing.\nl

 Proof. Suppose $x,y\in P$ and $x < y.$ Thus $\frac{x}{y} < 1$ and $0 <1 - \frac{x}{y} = 1 - \dfrac{1}{(\frac{y}{x})} \leq L(\frac{y}{x})  = L(y)-L(x).$ Therefore $L(x) <L(y).$\nl

\textbf{Definition 4.2.2.} If $a\in (P-\{1\}),$ let $log_{a}$ denote the function with domain $P$ such that $log_{a}(x) = \dfrac{L(x)}{L(a)} $ for each $x$ in $P,$ i.e.,  $log_{a} = \dfrac{L}{L(a)}.$\nl

\textbf{Theorem 4.2.4.} If $a,b\in(P-\{1\})$ and $x,y\in P,$ then \nl
\begin{tabular}{l l}
&\\
a) $log_{a}a = 1$ & d) $log_{a}\frac{1}{y} = -log_{a}y$ \\
&\\
b) $log_{a}(xy) = log_{a}x + log_{a}y$  &\\
&\\
c) $log_{a}\left(\frac{x}{y}\right) = log_{a}x - log_{a}y$& e) $log_{a}x = \dfrac{log_{b}x}{log_{b}a}$\\
\end{tabular}\\

Proof of (e). Suppose $a,b\in(P-\{1\})$ and $x\in P.$ Thus
\[
\dfrac{log_{b}x}{log_{b}a} = \dfrac{L(x)/L(b)}{L(a)/L(b)} = \dfrac{L(x)}{L(a)} = log_{a}x.
\]

\newpage
%page66
%Page66
%Section 4.2
\textbf{Theorem 4.2.5.}\nl
%\begin{center}
%\begin{tabular}{|c| c| c| c| c| c|}\hline
%$n$&0&1&3&7&9\\\hline
%1.&&.0414&.1139&.2304&.2788\\
%2.&.3010&&.3617&&.4624\\
%3.&.4771&.4914&&.5682&\\
%4.&&.6128&.6335&.6721&\\
%5.&&& .7243&&.7709\\
%6.&&.7853&&.8261&\\
%7.&.8451&.8513&.8633&&.8976\\
%8.&&&.9191&&.9494\\
%9.&&&&.9868&\\\hline
%\end{tabular} \\
%\end{center}
\begin{tabular}{l l l}
$E(1) = e = 2.718282$&$L(e) = 1$&$L(10) = 2.302585$\\
$L(2) = .693147$&&$Log_{10}(2) = .301029$\\
$L(3) = 1.098612$&&$Log_{10}(3) = .477121$\\
$L(5) = 1.609438$&&$Log_{10}(5) = .698970$\\
$L(7) = 1.945910$&&$Log_{10}(7) = .845098$\\
\end{tabular}\nl

This theorem is proved in Appendix 2.\nl

%\textbf{Theorem 4.2.6.} If $x\in P$ and $n\in \zp.$ then $L(x^{n}) =n L(x)$ and $x^{n} = %E\left(n(L(x))\right).$   The proof is by induction. \nl

\textbf{Example 4.2.1.} Find $log_{10}2100$ and $log_{b} 1.5.$ \nl
Solution\nl

\begin{tabular}{r l}
$log_{10}2100$&$=log_{10}3\cdot 7\cdot 10^{2}$\\
&$=log_{10}3 + log_{10}7 + log_{10}10^{2}$\\
&$= .4771 + .8451 + 2log_{10}10 $\\
&$=1.3222+2$\\
&$=3.3222$\\
&\\
$log_{10}1.5$&$=log_{10}\frac{3}{2}$\\
&$=log_{10}3 - log_{10}2$\\
&$= .4771 - .3010$\\
&$= .1761$\\
\end{tabular}\\

\textbf{Theorem 4.3.7.} $L$ has range $R.$\nl

\textbf{Exercises}\nl
Prove each of the following.\nl

\begin{tabular}{l l}
&\\
1. $log_{10}4 =.6021$ &7. Theorem 4.2.2(c) \\
&\\
2. $log_{10}5 =.6990$  &8. Theorem 4.2.2(d) \\
&\\
3. $log_{10}6 =.7782$ &9. Theorem 4.2.4(a)\\
&\\
4. $log_{10}750 = 2.8751$ &10. Theorem 4.2.4(b)\\
&\\
5. $log_{10}72 = 1.8573$ &11. Theorem 4.2.4(c)\\
&\\
6. Theorem 4.2.2(a)& 12. Theorem 4.2.4(d)\\
\end{tabular}

\newpage
%Page67
\lhead{Section 4.3}
\textbf{4.3 Exponential Functions}\nl

\textbf{Definition 4.3.1.} Let $E$ denote the inverse of $L$ and let $e$ denote $E(1).$\nl

\textbf{Theorem 4.3.1.} If $x,y\in R,$ then\nl
a) $E(0) = 1$\nl
b) $E(x + y) = E(x)E(y)$\nl
c) $E(y - x) = \dfrac{E(y)}{E(x)}$\nl
d) $E(-x) = \dfrac{1}{E(x)}$\nl
Proof of (b). Suppose $x,y\in R.$ Let $a = E(x)$ and $b = E(y).$ Thus $L(a) = x$ and $L(b) = y.$ Hence $E(x + y)=E(L(a) + L(b)) = E(L(ab)) = ab = E(x)E(y).$\nl

\textbf{Definition 4.3.2.} If $a\in P$ and $x\in R,$ let $a^{x}$ denote $E(xL(a))$  and if $x\neq 0,$ let $0^{x}$ denote $0$.\nl

\textbf{Theorem 4.3.2.} If $x\in P$ and $n\in \zp.$ then $L(x^{n}) =n L(x)$ and $x^{n} = E\left(n(L(x))\right).$   The proof is by induction. \nl

\textbf{Theorem 4.3.3.} If $a,b\in P, u\in (P-\{1\}),$ and $x,y\in R,$ then\nl

\begin{tabular}{l l}
a) $L(e)=1$&\\
b) $E(0) = 1$&     k) $a^{-x} = \dfrac{1}{a^{x}}$ \\
c) $L(x) = log_{e}x$&\\
d) $E(x) =e^{x}$  &  l) $(ab)^{x} = a^{x}b^{x}$\\
e) $1^{x} = 1 = a^{0}$&\\
f) $a^{1} = a$&       m)  $ \dfrac{a^{x}}{b^{x}} = \left(\dfrac{a}{b}\right)^{x}$\\
g) $a^{x +y} = a^{x}a^{y}$&\\
h) $a^{x + 1} = a^{x}a$&\\
i) $(a^{x})^{y} = a^{xy}$&    n) $log_{u}a^{x} = xlog_{u}a$\\
j) $a^{x - y} = \dfrac{a^x}{a^y}$&    o) $a = u^{x} \lra x = log_{u}a$\\
\end{tabular}\\

Proof of (g). Suppose $a\in P$ and $x,y\in R.$ \nl

\begin{tabular}{r l}

$a^{x + y}$&$= E((x+y)L(a))$\\
&$= E(xL(a)+yL(a))$\\
&$=E(xL(a))E(yL(a))$\\
&$=a^{x}a^{y}$\\
\end{tabular}\\

\textbf{Example 4.3.1.} Find the truth set of $\{x \mid log_{10}(x^{2} - 4) = 1 - log_{10}2\}.$\nl

Solution.  $\{x \mid log_{10}(x^{2} - 4) = 1 - log_{10}2\}$

\newpage
%Page68

%\textbf{Section 4.3}\nl

\begin{tabular}{r l}
$\{x \mid log_{10}(x^{2} - 4) = 1 - log_{10}2\}$& $= \{x \mid log_{10}(x^{2} - 4) + log_{10}2 = 1\}$ \\
&$= \{x \mid log_{10}(2x^{2} -8) = 1\}$\\
&$= \{x \mid 2x^{2} - 8 = 10^{1}\}$\\
&$= \{x \mid x^{2} - 9 = 0\}$\\
&$= \{x \mid (x-3)(x+3) = 0\}$\\
&$= \{x \mid x-3 = 0$ or $x + 3 = 0\}$\\
&$= \{x \mid x = 3$ or $x = -3\}$\\
&$= \{3, -3\}$\\
\end{tabular}\\



\textbf{Example 4.3.2.} Solve $5\cdot 2^{x-1} =3^{2x}$ for $x.$\nl
Solution. \nl
$5\cdot 2^{x-1} =3^{2x}$ \nl
$log_{10}5\cdot 2^{x-1} = log_{10}3^{2x}$ \nl
$log_{10}5 + log_{10}2^{x-1} = log_{10}3^{2x}$\nl
$log_{10}5 + (x-1)log_{10}2 = 2x log_{10}3$ \nl
$.6990+(x-l).3010 = 2x(.4771)$ \nl
$.6990 + .3010x -.3010 = .9542x$ \nl
$.6532x = .3980$ \nl
$x=.6100$ \nl

\textbf{Definition 4.3.3.} If $x\in R, c\in Z, 1 \leq a < 10,$ and $x=a\cdot 10^{c},$ then $c$ is called the \textit{characteristic} and $log_{10}a$  is called the \textit{mantissa} of $x.$\nl

%\textbf{Theorem 4.3.4.} If $n\in \zp,$ then $\left( \dfrac{n + 1}{n}\right)^{n} < e < \left( \dfrac{n}{n - %1}\right)^{n}.$\nl

%A proof can be obtained by applying Theorem 3.4.2. to
%\[
%S = \left\lbrace n \mid n\in \zp \text{~and~} \left(1 + \frac{1}{k}\right)^{n} < E\left( \dfrac{n}{k}\right) < %\left(\dfrac{1}{1 - \frac{1}{k}}\right)^{n}\right\rbrace
%\]
%where $k\in Z_{2}.$ \nl
%\
%\textbf{Theorem 4.3.5.} $2.5 < \left( \dfrac{8}{7}\right)^{7} < e < \left(\dfrac{7}{6}\right)^{7}  < 3.$\nl

\textbf{Exercises} If $log_{10}2 =.3010, log_{10}3 = .4771, log_{10}4 = .6021,$ and $log_{10}5 = .6990,$ find the truth set of each of the following.\nl

\begin{tabular}{l l}
&\\
1. $log_{10}(x^{2} - 1) - log_{10}(x + 1) = 1$  &    5. $3^{4x - 3} = 27$\\
&\\
2. $log_{10}4x + log_{10}x = 2$&     6. $4^{x} = 2\cdot 2^{x}$\\
&\\
3. $4\cdot 3^{x} =2^{2x-1}$&              7. $0 < log_{10}x < 1$\\
&\\
4. $2^{x}\cdot 5^{x-1} = 3$        &          8. $1 <2^{log_{10}x} < 2$\\
\end{tabular}

\newpage

%Page69
\lhead{Section 4.4}
Simplify each of the following expressions.\nl

9. $\left(  \dfrac{x^{1/3}}{x^{1/4}}   \right)^{6}$\nl

10. $\left(\dfrac{x^{2/3}}{x^{-(1/6)}}                \right)^{2}$\nl


%$\left(\dfrac{a^{-1}b^{2}}{a^{2}b^{-2}}   \right)^{-1}$\nl

%$\left(\dfrac{a^{3}b^{-1}}{ab^{2}}   \right)^{3}$\nl


11. $\dfrac{\left(\dfrac{a^{-1}b^{2}}{a^{2}b^{-1}}   \right)^{-1}}{\left(\dfrac{a^{3}b^{-1}}{ab^{2}}   \right)^{3}}$\nl

12. Prove Theorem 4.3.1(a). \nl

13. Prove Theorem 4.3.1(c) and (d). \nl

14. Prove Theorem 4.3.2(a) and (b). \nl

15. Prove Theorem 4.3.2(c), (d), and (e). \nl

16. Prove Theorem 4.3.2(f) and (h). \nl

17. Prove Theorem 4.3.2(i) and (j). \nl

18. Prove Theorem 4.3.2(k) and (l). \nl

19. Show that $E$ is increasing. \nl

20. Show that if $a\in P$ and $f(x) = a^{x},$ for each $x$ in $R,$ then $f$ is increasing if $a > 1$ and $f$ is decreasing if $0 < a < 1.$\nl

\vfill

\textbf{4.4 Radicals}\nl

\textbf{Definition 4.4.1}   If $a \geq 0$ and $n \in \zp,$ let $\sqrt[n]{a}$ denote $a^{1/n}$ and $\sqrt{a}$ denote $a^{1/2}.$\nl

\textbf{Theorem 4.4.1.}  If $a,b\in P$ and $n,m\in \zp,$ then\nl

a) $(\sqrt[n]{a})^{m} = \sqrt[n]{a^{m}}$\nl

b) $\sqrt[n]{ab} = \sqrt[n]{a}\sqrt[n]{b}$\nl

c) $\sqrt[n]{\dfrac{a}{b}} = \dfrac{\sqrt[n]{a}}{\sqrt[n]{b}}$\nl

\newpage
%page70

d)  $\sqrt[n]{\sqrt[m]{a}} = \sqrt[nm]{a}$\nl
Proof of (a).  $(\sqrt[n]{a})^{m} = (a^{1/n})^{m}  = a^{m/n} = (a^{m})^{1/n} = \sqrt[n]{a^{m}}. $\nl

\textbf{Example 4.4.1.}  Assuming $x\in P,$ simplify $\sqrt[3]{x\sqrt{x}}$\nl
Solution.\nl
$\sqrt[3]{x\sqrt{x}} = (x\sqrt{x})^{1/3} = (xx^{1/2})^{1/3} = x^{1/3}x^{1/6} = x^{\frac{(2+1)}{6}} = x^{1/2}.$\nl

\textbf{Example 4.4.2.} Simplify $\dfrac{2}{\sqrt{3}} + \dfrac{\sqrt{6}}{\sqrt{2}} + \dfrac{\sqrt{12}}{6}.$\nl
Solution.\nl
\begin{tabular}{r l}
$\dfrac{2}{\sqrt{3}} + \dfrac{\sqrt{6}}{\sqrt{2}} + \dfrac{\sqrt{12}}{6}$&$=\dfrac{2\sqrt{3}}{\sqrt{3}\sqrt{3}} +
\dfrac{\sqrt{3}\sqrt{2}}{\sqrt{2}} + \dfrac{\sqrt{4}\sqrt{3}}{6}$ \\
&\\
&$=\dfrac{2\sqrt{3}}{3} + \sqrt{3} + \dfrac{\sqrt{3}}{3}$\\
&\\
&$=\sqrt{3}\left(\frac{2}{3} + 1 +\frac{1}{3}\right)$\\
&\\
&$=2\sqrt{3}$\\
\end{tabular}\\
\textbf{Theorem 4.4.2.} If $x \geq 0,$ then $\sqrt{x} \geq 0.$\nl
\textbf{Theorem 4.4.3.}  If $x \in R,$ then $\abs{x} = \sqrt{x^{2}}.$\nl
Proof.  Suppose $x \in R$.\nl

\begin{tabular}{l l}
$x^{2} = ((x^{2})^{1/2})^{2}$& Th 4.3.2(i)\\
&\\
$x^{2} = (\sqrt{x^{2}})^{2}$&Def 4.4.1.\\
&\\
$\sqrt{x^{2}} \geq 0$ and $(\sqrt{x^{2}} = x \text{~or~} \sqrt{x^{2}} = -x)$&Th 2.3.6, Th 4.4.2.\\
&\\
$\abs{x} = \sqrt{x^{2}}$&Th 3.3.1.\\
\end{tabular}\\

\textbf{Theorem 4.4.4.}  If $a > 0$ and $b > 0$, then $a = b$ if and only if $\sqrt{a} = \sqrt{b}.$\nl

\textbf{Example 4.4.3.}  Solve the equation $x^{2} + 4x - 5 > 0$ for $x$.\nl
Solution\nl

\begin{tabular}{l l}
$x^{2} + 4x - 5 > 0$&$\lra (x^{2} + 4x +4) - 4 - 5 > 0$\\
&\\
&$\lra (x + 2)^{2} > 9$\\
&\\
&$\lra \sqrt{(x + 2)^{2}} > 3$\\
&\\
&$\lra \abs{x + 2} > 3$\\
&\\
&$\lra x + 2 < -3$ or $3 < x + 2$\\
&\\
&$\lra x < -5$ or $1 < x$\\
\end{tabular}

\newpage
%Page71
\lhead{Section 4.5}
\textbf{Exercises}\nl

Assuming $a,b\in P,$ simplify each of the following.

1. $\sqrt[4]{a^{}\sqrt[3]{a^{-7/2}\sqrt{a}}}$\nl

2. $\sqrt{2a^{2}b^{3}}\sqrt{8b}$\nl

3. $6\sqrt{12} - 2\sqrt{75} -10 \sqrt{\frac{3}{25}}$\nl

4. $-\frac{1}{\sqrt{3}} + \frac{\sqrt{12}}{2} - \frac{2\sqrt{2}}{6}$\nl

5. If $f(x) = \sqrt{x} + 1$  and $g(x) = x^{2},$ and $D_{f\circ g} = D_{g} =R,$ find $f\circ g.$\nl

6. If $f(x) = x^{2} + 1$ and $g(x) = \sqrt{x},$ find $f\circ g.$\nl

7. Find the inverse of $f$ if $f(x) = 1 - x^{2}$ and $D_{f} = [0,2].$\nl

8. Find the inverse of $f$ if $f(x) = \sqrt{x} - 1.$\nl

9. Solve $x^{2} - 6x - 7 < 0$ for $x.$\nl

10. Prove Theorem 4.4.1(b). \nl

11. Prove Theorem 4.4.1(c). \nl

12. Prove Theorem 4.4.1(d). \nl

13. Prove Theorem 4.4.4.(Hint:Use Theorem 3.1.9) \nl

\vfill

\textbf{4.5 Factoring}\nl

\textbf{Theorem 4.5.1.} If $a,b,c,d,x,y\in R,$ then \nl
a) $ax + ay = a(x + y)$\nl
b) $x^{2} - y^{2} = (x + y)(x - y)$\nl
c) $x^{2} + 2xy + y^{2} = (x + y)^2$\nl
d) $x^{2} - 2xy + y^{2} = (x - y)^{2}$\nl
e) $x^{2} + (a + b)x + ab =(x + a)(x + b)$\nl
f) $acx^{2} + (ad + bc)x + bd = (ax + b)(cx + d) $\nl

\newpage
%page72

%Page72
\textbf{Example 4.5.1.}  Factor, i. e., write as a product, $(x - 2y)^{2} - 4y^{2}.$\nl
Solution\nl
\begin{tabular}{r l l}
$(x - 2y)^{2} - 4y^{2}$&$=[(x - 2y) + 2y][(x - 2y) - 2y]$&Th 4.5.1(b)\\
& $= x(x -4y)$&\\
\end{tabular}\\

\textbf{Example 4.5.2.} Factor $x^{4} + 2x^{2}y^{2}  + 9y^{4}.$ \nl
Solution.\nl
\begin{tabular}{r l l}
$x^{4} + 2x^{2}y^{2}  + 9y^{4}$&$=x^{4} + 6x^{2}y^{2} + 9y^{4} - 4x^{2}y^{2} $&\\
&$=(x^{2} + 3y^{2})^{2} -(2xy)^{2}$&Th 4.5.1(c)\\
&$=(x^{2} + 3y^{2} + 2xy)(x^{2} + 3y^{2} - 2xy)$&Th 4.5.1(b)\\
\end{tabular}\\

\textbf{Example 4.5.3.} Factor $xa - 2xb + ya - 2yb.$\nl
Solution.\nl
\begin{tabular}{r l l}
$xa - 2xb + ya - 2yb$&$=x(a - 2b) - y(a - 2b)$&Th 4.5.1\\
&$=(a - 2b)(x - y)$&Th 4.5.1(a)\\
\end{tabular}\\

\textbf{Example 4.5.4.} Factor $(x+3)^{2} + 2(x+3) - 15.$\nl
Solution.\nl
\begin{tabular}{r l l}
$(x+3)^{2} + 2(x+3) - 15$&$= [(x+3)-3] [(x+3)+5]$&Th 4.5.1(e)\\
&$=x(x + 8)$&\\
\end{tabular} \\

\textbf{Theorem 4.5.2.} If $x,y\in R,$ then \nl
a) $x^{3} + 3x^{2}y + 3xy^2 + y^{3} = (x + y)^{3}$\nl
b) $x^{3} - 3x^{2}y + 3xy^2 - y^{3} = (x - y)^{3}$\nl
c) $x^{3} + y^{3} =(x + y)(x^{2}  - xy + y^{2})$\nl
d) $x^{3} - y^{3} =(x - y)(x^{2}  + xy + y^{2})$\nl

\textbf{Example 4.5.5.} Factor $x^{3} - 6x^{2} + 12x - 8.$\nl Solution.\nl

\begin{tabular}{r l l}
$x^{3} - 6x^{2} + 12x - 8$&$=x^{3} - 3x^{2}(2) + 3x(2^{2})- 2^{3}$&\\
&$=(x - 2)^{3}$&Th 4.5.2(b)\\
\end{tabular}\\

\textbf{Example 4.5.6.} Factor $x^{3}+27.$\nl
Solution.\nl

\begin{tabular}{r l l}
$x^{3}+27$&$= x^{3} + y^{3}$&\\
&$=(x + 3)(x^{2} - 3x + 9)$&Th 4.5.2(c)\\
\end{tabular}

\newpage
%page73
\lhead{Section 4.6}
\textbf{Exercises}\nl
Factor each of the following. \nl

1. $(a+b)^{2} - (b-a)^{2}$\nl

2. $(x+3)^{2} + (x+3) - 12$\nl

3. $(x - y)^{2} + 3(x - y)y - 10y^{2}$\nl

4. $4x^{4} + 1$ \nl

5. $1 + 8x^{3}$\nl

6. $2y^{2} + 8y +8$\nl

7. $x^{2} - 6x + 2 -(y^{2} - 6x + 2)$\nl

8. $(x - y)^{2} - y^{2}$\nl

9. $x^{4} - 4x^{2}y^{2} + 4y^{4}$\nl

10. $x^{2} - 4x + 4 $\nl

11. $x^{2} + 4x + 4y - y^{2}$\nl

12. $x^{3} + 9x^{2}y + 27xy^{2} + 27y^{3}$\nl
Simplify each of the following.

13. $[(a+b)(a - b)^{-1} -(a-b)(a+b)^{-1}](a^{2} - b^{2})$\nl

14. $\dfrac{-(1 - x^{2})+ (1 + x)^{2}}{1 + x}$\nl

\vfill
\textbf{ 4.6 Polynomials}\nl

\textbf{Definition 4.6.1.}  A function $f$ is said to be a \textit{polynomial} if and only if there is an $n\in Z_{0}$ (called the \textit{degree} of $f$) and there is a function $a : Z_{0}^{n} \ra R$ such that $a_{0} \neq 0$ and \nl
\[
f = \sum_{i = 0}^{n}a_{i}I^{n - i}, \text{~i. e.~} f(x) = \sum\limits_{i=0}^{n}a_{i}x^{n - i}, \text{~for~each~} x \text{~in~} R.
\]

\textbf{Theorem 4.6.1.} If $r \in R, n\in Z_{0}, a: Z_{0}^{n} \ra R, b_{0} = a_{0} \neq 0, b_{i} = b_{i-1}r + a_{i},$
for each $i$ in $Z_{1}^{n},$ and $f = \sum_{i=0}^{n-1}a_{i}I^{n-i},$ then\nl
a) $f=(I - r)\left(\sum_{i=0}^{n-1}b_{i}I^{n-1-i}\right) + b_{n}$\nl
b) $f(r) = b_{n}$\nl
c) If $f(r) = 0,$ then $f = (I-r)\sum_{1=1}^{n-1}b_{i}I^{n-1-i}$\nl

\newpage
%page74

Proof of (a). Suppose $x,r \in R, n\in Z_{0}, a:Z_{0}^{n} \ra R, b_{0} = a_{0} \neq 0, b_{i} = b_{i-1}r + a_{i},$
for each $i$ in $Z_{1}^{n},$ and $f = \sum_{i=0}^{n}a_{i}I^{n-i}.$\nl

\fl $(x - r)\left(\sum\limits_{i=0}^{n-1}b_{i}x^{n-1-i} \right) + b_{n}$\nl

$= \left(\sum\limits_{i=0}^{n-1}b_{i}x^{n-i} \right) - \left(\sum\limits_{i=0}^{n-1}b_{i}rx^{n-1-i} \right) + b_{n}$\nl

$= b_{0}x^{n} + \left(\sum\limits_{i=1}^{n-1}b_{i}x^{n-i} \right) -  \left(\sum\limits_{i=0}^{n-1}b_{i}rx^{n-1-i} \right) + b_{n}$\nl

$= b_{0}x^{n} + \left(\sum\limits_{i=1}^{n-1}(b_{i-1}r + a_{i})x^{n-i} \right) -  \left(\sum\limits_{i=0}^{n-1}b_{i}rx^{n-1-i} \right) + b_{n}$\nl

$= a_{0}x^{n} + \overbrace{\left(\sum\limits_{i=1}^{n-1}b_{i-1}rx^{n-i} \right)}^{u} + \left(\sum\limits_{i=1}^{n-1}a_{i}x^{n-i} \right)  \overbrace{-~~\left(\sum\limits_{i=0}^{n-2}b_{i}rx^{n-1-i} \right)}^{-u} -b_{n-1}r$\nl
 $ + (b_{n-1}r +a_{n}) $\nl

We pause to note that the two overbraced sums differ only in sign and that both $b_{n-1}r$ and $-b_{n-1}r$ appear in the sum.  Thus, the expression in the line above reduces to\nl

$\sum\limits_{i=0}^{n} a_{i}x^{n-i} = f(x).$  Thus, $f = (I - r)\left(\sum\limits_{i=0}^{n-1}b_{i}I^{n-1-i}\right) + b_{n}.$\nl

\textbf{Example 4.6.1.}  If $f(x) = 2x^{3} - 7x^{2} + 9$ for each $x$ in $R,$ find $f(3).$  If $f(3) = 0,$ find a polynomial $q$ such that $f(x) = (x - 3)q(x),$ for each $x$ in $R.$ \nl
Solution.  Suppose $f(x) = 2x^{3} - 7x^{2} + 9.$
The computation may be done in the following manner.\nl
\begin{tabular}{lllll|l}
$a_{0}$&$a_{1}$&$\cdots$&$a_{n-1}$&$a_{n}$&$r$\\
&$b_{0}r$&$\cdots$&$b_{n-2}r$&$b_{n-1}r$&\\ \cline{1-5}
$b_{0}=a_{0}$&$b_{1}=b_{0}r + a_{1}$&$\cdots$&$b_{n-1}=b_{n-2}r+a_{n-1}$&$b_{n}=b_{n-1}r + a_{n}$&\\
\end{tabular}\nl
%\begin{tabular}{lllll}\hline
%$b_{0}=a_{0}$&$b_{1}=b_{0}r + a_{1}$&$\cdots$&$b_{n-1}=b_{n-2}r+a_{n-1}$&$b_{n}=b_{n-1}r + a_{n}$\\
%\end{tabular}\nl

For this problem the computation would be done in the following manner.\nl

\begin{tabular}{rrrr|l}
2&-7&0&9&3\\
&6&-3&-9&\\
\end{tabular}\nl
\begin{tabular}{rrrl}\hline
2&-1&-3&~0
\end{tabular}

Thus $f(3) = 0$ and $f(x) = (x - 3)(2x^{2} - x - 3),$ for each $x$ in $R.$  This process is called synthetic division.\nl

\newpage
%page 75

\textbf{Theorem 4.6.2.}  If $x,a\in R$ and $n\in \zp,$ then $x^{n} - a^{n} = (x - a)\sum\limits_{i=1}^{n}a^{i-1}x^{n-i}.$\nl

\textbf{Definition 4.6.2.} A number $c$ is called a \textit{root of a polynomial} $f$ if
and only if $f(c) = 0.$\nl

\textbf{Theorem 4.6.3.} A number $c$ is a root of a polynomial $f$ if and only if $(x - c)$ is a factor of $f.$\nl

\textbf{Theorem 4.6.4.} A polynomial of degree $n$ has at most $n$ roots. \nl

\textbf{Theorem 4.6.5.} If $f= \sum\limits_{i=0}^{n}a_{i}I^{n-i}$ and $g= \sum\limits_{i=0}^{n}b_{i}I^{n-i}$ are polynomials such that $f(x) = g(x)$ for $n + 1$ values of $x,$ then $a_{i} =b_{i},$ for each $i,$ and $f = g.$\nl

\textbf{Example 4.6.2.} Find values of $A, B, C,$ and $D$ such that


(1) $\dfrac{x^{2} + 8}{x(x - 2)^{3}} = \dfrac{A}{x} + \dfrac{B}{x - 2} + \dfrac{C}{(x - 2)^{2}} + \dfrac{D}{(x - 2)^{3}},$\nl
for each $x$ in $(R-\{ 0, 2\}).$\nl

Solution. By multiplying both sides of equation (1) by $x(x-2)^{3},$ we get \nl
(2) $x^{2} + 8 = A(x - 2)^{3} + Bx(x - 2)^{2} + Cx(x - 2) +Dx.$\nl

Suppose equation (2) is true for each $x$ in $R.$ Letting $x = 0,$ equation (2) becomes $8 = -8A.$ Thus $A = -1.$ By letting $x = 2,$ equation (2) yields $12 =2D$ and $D = 6.$ Therefore, equation (2) may be written as \nl

(3) $x^{2} + 8= (B-1)x^{3} + (6 - 4B + C)x^{2} + (-6 + 4B -2C)x + 8.$\nl

Since $\zp \subseteq R$ and $\zp$ is infinite, then $R$ is infinite. By Th 4.6.5, the coef­ficients of
like terms of equation (3) must be equal. Equating coefficients, we get $0 = B - 1$ and $1=6 -4B + C.$ Hence $B = 1$ and $C = -1.$ Equation (2) may now be written as \nl

(4) $x^{2} + 8= -1(x-2)^{3} + x(x-2)^{2} - x(x - 2) + 6x.$\nl

The fact that equation (4) is true for each $x$ in $R$ may be checked by mul­tiplication. By dividing both sides of equation (4) by $x(x-2)^{3},$ we have \nl
$\dfrac{x^{2} + 8}{x(x - 2)^{3}}= \dfrac{-1}{x} + \dfrac{1}{x - 2} + \dfrac{-1}{(x - 2)^{2}}  + \dfrac{6}{(x - 2)^{3}},$\nl
for each $x$ in $(R-\{0, 2 \}).$  The terms on the right-hand side of equation (1) are called partial fractions. \nl

\textbf{Definition 4.6.3.} If the number $c$ is a root of a polynomial $f,$ then the greatest integer $n$ such that $(x - c)^{n}$ is a factor of $f$ is called the \textit{multiplicity of the root} $c.$\nl

\newpage
%page76

\textbf{Theorem 4.6.6.} If the number $c$ is a root of a polynomial $f,$ then $n$ is the multiplicity of the
root $c$ if and only if there is a polynomial $q$ such that $f(x)=( x - c)^{n}q(x),$ for each $x$ in $R,$ and
$q(c) \neq 0$

\textbf{Example 4.6.3.} Find the multiplicity of the root $-2$ of $f(x) =x^{4} + 5x^{3} + 6x^{2} -4x -8.$\nl Solution. \nl

\begin{tabular}{rrrrr|l}
1&-5&6&-4&-8&-2\\
&-2&-6&0&8&\\
\end{tabular}\nl
\begin{tabular}{rrrr|l}\hline
1&$~3$&$~$0&-4&-2\\
&-2&-2&4&\\
\end{tabular}\nl
\begin{tabular}{rrr|l}\hline
1&1&-2&-2\\
&-2&2&\\
\end{tabular}\nl
\begin{tabular}{rrl}\hline
$1$&$-1$&$~~~~$\\
\end{tabular}

Therefore $f(x)=(x+2)^{3}(x-1)$ and $(-2,-1) \neq 0$ By Th 4.6.6, we see that $-2$ is a root of $f$ of
multiplicity 3. \nl

\textbf{Theorem 4.6.7.} If $a,b,c\in R$ and $f(x)=ax^{2} + bx + c,$ for each $x$ in $R,$ then

a) $\dfrac{-b + \sqrt{b^{2} - 4ac}}{2a}$ and $\dfrac{-b - \sqrt{b^{2}- 4ac}}{2a}$ are the only roots of $f$
if $b^{2} - 4ac \geq 0.$\nl

b) $f$ has no roots if $b^{2} - 4ac < 0.$\nl

The formula in a) for finding the roots of a quadratic equation is called the \textit{Quadratic Formula.}\nl

\textbf{Exercises}\nl
For each of the following definitions of $f$ and $r,$ find a polynomial $q$ and a
number $s$ such that $f(x)=(x - r)q(x) + s,$ for each $x$ in $R.$\nl

\begin{tabular}{l l}
1. $r=-1$& $f(x)= 2x^{3} + 3x^{2} - 4x -1$ \\

2. $r=3$ &$f(x)= 2x^{4}- 5x^{3} - 4x^{2} + 8$ \\

3. $r=2$ &$f(x)= x^{4} - 3x^{2} -x - 3$ \\
\end{tabular}\nl


For each of the following definitions of $f$ and $r,$ find the multiplicity of the root $r$ of $f.$ \nl

\begin{tabular}{l l}

4. $r=1$ &$f(x)=x^{4}+ x^{3}- 5x^{2} + 3x$ \\

5. $r=-2$ &$f(x)=x^{3} -x^{2} - 2x + 8$ \\

6. $r=-1$& $f(x)=x^{5} + 3x^{4} + 2x^{3} - 2x^{2} -3x -1$ \\

\end{tabular}

7. Finds the roots of $f(x)= 2x^{2} - 3x - 1.$ \nl

\newpage
%page77

8. Find the roots of $f(x) = x^{2} + 3x - 2.$\nl

9. Find values of $A, B, C,$ and $D$ such that $\dfrac{x^{3} -3x^{2} + 2x - 3}{(x^{2} + 1)^{2}} = \dfrac{Ax + B}{x^{2} + 1} + \dfrac{Cx + D}{(x^{2} + 1)^{2}}$ for each $x$ in $R.$\nl

10. Prove Theorem 4.6.2. \nl­

11. Prove Theorem 4.6.3. (Hint: Apply Th 4.6.2 to $f(x)-f(c)$.) \nl

12. Prove Theorem 4.6.4. (Hint: Apply Th 3.4.2 to $S= \{n \mid n\in \zp $ and every polynomial of degree $n$ has at most $n$ roots$\}.$\nl

13. Prove Theorem 4.6.5. (Hint: Apply Th 4.6.4 to $f - g$.) \nl

14. Prove Theorem 4.6.6.\nl

15. Prove Theorem 4.6.7.\nl

\newpage
%page78
\lhead{Section 5.1}
\textbf{Chapter 5}\nl

\textbf{Trigonometry}\nl

\textbf{5.1 Six Fundamental Identities}\nl

\textbf{Definition 5.1.1.} Let $R^{2}$ denote $R\times R$ and let $0$ denote $(0, 0). $\nl

\textbf{Definition 5.1.2.}  If $(a, b), (x,y)\in R^{2},$ let $d((a,b),(x, y))$ denote
\[
\sqrt{(x-a)^{2}+ (y-b)^{2}}.
\]

\textbf{Definition 5.1.3.}  Let $c_{1}$ denote $\{ u \mid u\in R^{2}$ and $d(u, 0) = 1\}.$\nl

The proof of the following theorem will not be given, as it is beyond the scope of this course.\nl

\textbf{Theorem 5.1.1.}   There is a positive number $\pi$ and a function $\alpha$ (called the \textit{wrapping function})
with domain $R$ and range $c_{1}$ such that\nl

T1.  $\alpha(0) = (1,0), \alpha\left(\dfrac{\pi}{2}\right) = (0, 1),$ and $\alpha(\pi) = (-1, 0).$\nl

T2.  If $x,y\in R,$ then $d(\alpha(x-y), \alpha (0)) = d(\alpha (x), \alpha(y)).$\nl

T3.  If $0 < x < \dfrac{\pi}{2},$ and $\alpha (x) = (a, b),$ then $a \neq 0$ and $0 < b < x < \dfrac{b}{a}.$\nl

\textbf{Definition 5.1.4.} Let $\sin$ and $\cos$ denote the functions with domain $R$ such that for each $x$ in $R, \alpha (x) = (\cos(x), \sin(x)).$  The trigonometric functions are generally written without parentheses, i. e., $\sin(x)$ is written $\sin~ x.$\nl

\textbf{Definition 5.1.5.}  Let $\tan, \cot, \sec, $ and $\csc$ denote the functions defined by
\[
\tan = \dfrac{\sin}{\cos},~~ \cot = \dfrac{\cos}{\sin},~~ \sec = \dfrac{1}{\cos}, \text{~and~} ~~\csc = \dfrac{1}{\sin}.
\]
$ $\nl

\textbf{Definition 5.1.6.}  If each of $f$ and $g$ is a function, then $f(x) = g(x)$ is said to be an \textit{identity} if and only if $f(x) = g(x)$ for each $x$ in $D_{f}\cap D_{g}.$\nl

\textbf{Theorem 5.1.2.}  The following are identities.\nl

ID1.  $\sin~x ~\text{~csc~}x = 1$\nl

ID2.  $\cos~x~\sec~x = 1$\nl


\newpage
%page79
\begin{tabular}{l l}
ID3. (a)$~ \tan ~x ~\cot~ x ~=~1,$ &(b) $ \tan~ x~= \dfrac{1}{\cot~x}$\\

ID4. (a)$~\sin^{2~}x + \cos^{2}x~=~1,~$ &(b) $ \sin^{2}x~=~1 - \cos^{2}x$\\

ID5. (a)$~\sec^{2}x = 1 + \tan^{2}x,$& (b) $ 1 = \sec^{2}x - \tan^{2}x$\\

ID6. (a)$~\csc^{2}x = 1 + \cot^{2}x,$& (b) $1 = \csc^{2}x - \cot{2}x$\\
\end{tabular}\\
Proof of ID4(a). Suppose $x\in R$.\nl
$\alpha(x) \in c_{1}$\nl
$(\cos~ x, \sin ~x)\in c_{1}$\nl
$\sqrt{(\cos~x - 0)^{2} + (\sin~x - 0)^{2}} = 1 $ \nl
$\cos^{2}x + \sin^{2}x = 1. $\nl
Proof of ID5(a). Suppose $x\in R$. \nl

$\sin^{2}x + \cos^{2}x = 1$\nl

$\dfrac{\sin^{2}x}{\cos^{2}x} + \dfrac{\cos^{2}x}{\cos^{2}x} = \dfrac{1}{\cos^{2}x}$\nl

$\tan^{2}x + 1 = \sec{2}x$\nl

\textbf{Example 5.1.1.} Show that $\tan x + \cot x = \sec x \csc x$ is an identity. \nl
Solution. \nl
\begin{tabular}{r l l}
$\tan ~x + \cot~ x$ & $= \dfrac{\sin~ x}{\cos ~x} + \dfrac{\cos ~x}{\sin~ x}$ & Def 5.1.5 \\
&&\\
& $= \dfrac{\sin^{2}x + \cos^{2}x}{\sin~x ~\cos~x}$& \\
&&\\
&$=\dfrac{1}{\cos~x ~\sin~x}$&ID4\\
&&\\
&$= \sec~x ~\csc~x$&Def 5.1.5

\end{tabular}


\textbf{Example 5.1.2.} Show that  $\dfrac{\sin~x}{1 - \cos~x} = \dfrac{1 + \cos~x}{\sin~x}$is an identity. \nl
Solution.  \nl
\begin{tabular}{r l}
&\\
$\dfrac{\sin ~ x}{1 -\cos~x}$&$= \dfrac{\sin~x}{1 - \cos~x}\cdot \dfrac{1 + \cos~x}{1 + \cos~x}$\\
&\\
&$=\dfrac{\sin~x(1 + \cos~x)}{1 - \cos^{2}x}$\\
&\\
&$=\dfrac{\sin~x(1 + \cos~)}{\sin^{2}~x}$\\
&\\
&$=\dfrac{1 + \cos~x}{\sin ~x}$\\
\end{tabular}\\
\newpage
%page80
\textbf{Exercises}\nl

Prove the following identities.\nl

1. $\dfrac{\cos x}{\sec x} + \dfrac{\sin x}{\csc x} = 1$\nl

2. $\dfrac{1}{1 - \sin x} + \dfrac{1}{1 + \sin x} = 2\sec ^{2} x$\nl

3. $\dfrac{1 + \cos x}{1 - \cos x} - \dfrac{1 - \cos x}{1 + \cos x} = 4\cot x \csc x$\nl

4. $\dfrac{1 + \cos x}{\sin x} + \dfrac{\sin x}{1 + \cos x} = 2 \csc x$\\

5. $\dfrac{1}{\tan x + \cot x } = \dfrac{\sin x}{\sec x}$\nl

6. $\cos x + \tan x \sin x = \sec x$\nl

7. $\dfrac{\cot x + 1}{\sin x + \cos x}= \text{~csc~} x$\nl

8. $\dfrac{\sec x + 1}{\tan x} = \dfrac{\tan x}{\sec x - 1}$\nl

9. $\dfrac{1 - \cos x}{1 + \cos x} = (\csc x - \cot x)^{2}$\nl

10. $(\sin x + \cos x)(\cot x + \tan x) = \sec x + \csc x$\nl

11. $\dfrac{\tan x + \tan y}{\cot x + \cot y} = \tan x \tan y$\nl

12. $\dfrac{\sin x + 2\sin x \cos x}{2 + \cos x -2\sin^{2} x} = \tan x$\nl

13. $\dfrac{\sec x}{\csc x} + \dfrac{\sin x}{\cos x}  = 2\tan x$\nl

14. $\dfrac{\tan^{2}x}{\sec x - 1} = 1 + \sec x$\nl

15. $\dfrac{1}{\tan x + \sec x} = \sec x- \tan x$\nl

16. $\dfrac{\sin x + \cos x \tan x}{\tan x} = 2\cos x$\nl

17. $\dfrac{\tan x - \sin x \cos x}{\sin^{2}x} = \tan x$\nl

18. Prove that ID6(a) is an identity.

\newpage
%page81
\lhead{Section 5.2}
\textbf{5.2 The Sine and Cosine Addition Formulas}\nl

\textbf{Theorem 5.2.1.} The following are identities. \nl

ID7. $\cos(x - y) = \cos x\cos y + \sin x \sin y $\nl
ID8. $\cos(\dfrac{\pi}{2} - x) = \sin x$\nl
ID9. $\sin(\dfrac{\pi}{2} - x) = \cos x $\nl
ID10. $\cos(-x) = \cos x $\nl
ID11. $\sin(-x) = -\sin x$\nl
ID12. $\tan(-x) = -\tan x$\nl
ID13a. $\cos(x+y) = \cos x \cos y - \sin x \sin y $\nl
ID13b. $\cos(x-y) = \cos x \cos y + \sin x \sin y $\nl
ID14. $\sin(x+y) = \sin x \cos y + \cos x \sin y$\nl
ID15. $\sin(x-y) = \sin x \cos y - \cos x \sin y$\nl
ID16. $\tan(x + y) = \dfrac{\tan x + \tan y}{1 - \tan x \tan y}$\nl
ID17. $\tan(x - y) = \dfrac{\tan x - \tan y}{1 + \tan x \tan y} $\nl


\textbf{Proof of ID7.} Suppose $x,y\in R.$ \nl
\begin{tabular}{r l l}
$d(\alpha(x-y),\alpha(0))$&$= d(\alpha(x), \alpha(y)) $& T2 \\
$d((\cos (x-y), \sin (x-y)), (1,0))$&$=d(((\cos x, \sin x),( \cos y, \sin y))$ &Def 5.1.4\\
$\sqrt{(\cos(x-y)- 1)^2 + \sin^{2}(x-y)}$&$=\sqrt{(\cos x - \cos y)^2 + (\sin x - \sin y)^2}$ &  Def 5.1.2 \\
\end{tabular}\\
By squaring both sides and collecting terms, we obtain \nl
$2 - 2\cos(x-y) = 2-2\cos x\cos y- 2\sin x\sin y,$ which simplifies to
$\cos(x-y)$ $= \cos x\cos y + \sin x \sin y.$\nl
\textbf{Proof of ID8.} Suppose $x \in R.$ \nl
\begin{tabular}{r l l}
$\cos\left(\dfrac{\pi}{2}-x\right)$ &$= \cos\left(\dfrac{\pi}{2}\right)\cos x + \sin\left(\dfrac{\pi}{2}\right) \sin x$ &ID7\\
&$= 0\cos x + 1\sin x$ &T1\\
&$ =\sin x$ &\\
\end{tabular}\\
\textbf{Proof of ID9.} Suppose $x\in R.$\nl
\begin{tabular}{r l l}
$\cos x$ &$= \cos\left(\left(\dfrac{\pi}{2}\right) - \left(\dfrac{\pi}{2} -x\right)\right)$ & \\
&$=\sin\left(\po{2} - x\right)$ &ID8\\
\end{tabular}\\
\textbf{Proof of ID10.} Suppose $x\in R.$ \nl
\begin{tabular}{r l l}
$\cos(-x)$ &$= \cos(0 - x)$& \\
&$ = \cos 0 \cos x + \sin 0 \sin x$& by ID7. \\
%\newpage
%page82
&$= 1\cos x + 0\sin x$&T1\\
&$= \cos x$&\\
\end{tabular}\\
\newpage
%page82

%page82
Proof of ID11. Suppose $x\in R.$ \nl
\begin{tabular}{r l l}
$\sin (-x)$&$=\cos\left(\dfrac{\pi}{2} - (-x)\right)$&ID8 \\
&&\\
&$=\cos\left(\pi - \left(\dfrac{\pi}{2}- x\right)\right)$&\\
&&\\
&$= \cos \pi \cos \left(\dfrac{\pi}{2} - x\right) + \sin \pi \sin\left(\dfrac{\pi}{2} - x\right)$&ID7\\
&&\\
&$=-1\sin x + 0\sin x$&T1\\
&&\\
&$=-\sin x$& \\
\end{tabular}\\
Proof of ID13. Suppose $x,y\in R.$ \nl
\begin{tabular}{r l l}
$\cos(x+y)$&$= \cos(x-(-y))$&\\
&$=\cos x\cos(-y) + \sin x\sin(-y)$&ID7\\
&$=\cos x\cos y - \sin x\sin y$&ID10,ID11\\
\end{tabular}\\


Proof of ID14. Suppose $x,y\in R.$\nl
\begin{tabular}{r l l}
$\sin (x + y)$&$=\cos \left(\po{2} - (x + y)\right)$&ID8\\
&&\\
&$=\cos \left(\left(\po{2} - x\right) - y\right) $&\\
&&\\
&$=\cos \left( \po{2} - x \right)\cos y + \sin\left(\po{2} - x\right)\sin y $&ID7\\
&&\\
&$=\sin x \cos y + \cos x \sin y$&ID8,ID9\\
\end{tabular} \\

Proof of ID15. Suppose $x,y\in R.$ \nl
\begin{tabular}{r l l}
$\sin(x-y)$&$= \sin(x+(-y))$& \\
&$= \sin x \cos(-y) + \cos x \sin(-y)$&ID14\\
&$= \sin x \cos y - \cos x \sin y$&ID10,ID11\\
\end{tabular}\\

\textbf{Theorem 5.2.2.} The following are identities. \nl
ID18. $\sin x = \sin( \pi-x)$\nl
ID19. $\cos x = -\cos (\pi -x)$\nl
ID20. $\tan x = -\tan(\pi -x )$\nl
ID21. $\sin x = -\sin(x - \pi)$\nl
ID22. $\cos x = -\cos(x - \pi)$\nl
ID23. $\tan x = \tan(x - \pi)$\nl

\textbf{Theorem 5.2.3.} $\alpha (2\pi)=(1,0)$ and $\alpha\left(\dfrac{3\pi}{2}\right) = (0,-1). $\nl

\textbf{Theorem 5.2.4.} The following are identities. \nl
ID24. $\sin x = -\sin(2\pi - x)$\nl
1D25. $\cos x = \cos( 2\pi - x)$\nl
ID26. $\tan x = -\tan( 2\pi - x)$\nl
ID27. $\sin x = \sin(x \pm 2\pi) $\nl
ID28. $\cos x = \cos(x \pm 2\pi)$\nl
ID29. $\tan x = \tan(x \pm 2\pi) $\nl
\newpage

%page83
\lhead{Section 5.3}
\textbf{Exercises}\nl
1. Prove 1016. \nl
2. Prove ID17. \nl
3. Prove ID18. \nl
4. Prove ID19. \nl
5. Prove ID20. \nl
6. Prove ID21. \nl
7. Prove ID22. \nl
8. Prove ID23. \nl
9. Prove ID24. \nl
10. Prove ID25. \nl
11. Prove ID26. \nl
12. Prove ID27. \nl
13. Prove ID28. \nl
14. Prove ID29. \nl
15. Prove Theorem 5.2.3. \nl

16. Show that $\cot (x + y) = \dfrac{\cot x \cot y - 1}{\cot x + \cot y}$ is an identity.\nl

17. Show that $\cot (x - y) = \dfrac{\cot x \cot y + 1}{\cot y - \cot x}$ is an identity.\nl

%\vfill
\textbf{5.3 Application of the Addition Formulas}\nl

Here is a memory crutch to help you remember the following theorem:\nl
All Students Take Calculus.\nl
\textbf{A} in All stands for ``all three functions $\sin, \cos,$ and $\tan$ are positive in the first quadrant.''\nl
\textbf{S} in Students stands for ``only $\sin$ is positive in the first quadrant.''\nl
\textbf{T} in Take stands for ``only $\tan$ is positive in the third quadrant.''\nl
\textbf{C} in Calculus stands for ``only $\cos$ is positive in the fourth quadrant.''\nl

\textbf{Theorem 5.3.1.}\nl

a) If $x\in \left(0, \po{2}\right),$ then $\sin x >0, \cos x > 0,$ and $\tan x > 0.$

%\newpage
b) 	If $x\in \left(\po{2},\pi\right),$ then $\sin x >0, \cos x < 0,$ and $\tan x < 0$. \nl
c) 	If $x\in\left(\pi,\dfrac{3\pi}{2}\right),$ then $\sin x < 0, \cos x < 0,$ and $\tan x > 0.$ \nl
d) 	If $x\in \left(\dfrac{3\pi}{2},2\pi\right),$ then $\sin x < 0, \cos x > 0,$ and $\tan x < 0.$ \nl

Proof of (a). Suppose $x\in \left(0, \po{2}\right).$ By Def 5.1.4 and T3, we have $\alpha(x) = (\cos x, \sin x), \cos x \neq 0, $ and $ 0 < \sin x < x < \dfrac{\sin x}{\cos x} = \tan x.$\nl
If $\cos x < 0,$ then $0 < \dfrac{\sin x}{\cos x}$ implies $0 > \sin x,$ which is a contradiction.  Thus, $\cos x > 0.$\nl
\newpage
%page84
Proof of (b). Suppose $x\in \left(\po{2}, \pi \right).$ Thus $\po{2} < x < \pi, -\po{2} > -x> -\pi,$ and
$\po{2} > \pi - x > 0.$ Hence\nl
\begin{tabular}{r l l}
$\sin x$&$= \sin(\pi - x) > 0$ &ID18, Th 5.3.1(a) \\
$\cos x$&$=-\cos(\pi - x) < 0$& ID19, Th 5.3.1(a)\\
$\tan x$&$= -\tan(\pi - x) < 0$& ID20, Th 5.3.1(a) \\
\end{tabular}\\

\textbf{Theorem 5.3.2.}\nl

a) If $x\in \left[0,\po{2}\right],$ then $\sin x \geq 0$ and $\cos x \geq 0.$\nl

b) 	If $x\in  \left[\po{2},\pi\right],$ then $\sin x \geq 0$ and $\cos x \leq 0.$\nl

c) 	If $x\in  \left[\pi,\dfrac{3\pi}{2}\right],$ then $\sin x \leq 0$ and $\cos x \leq 0.$\nl

d)  If $x\in  \left[\dfrac{3\pi}{2},2\pi\right],$ then $\sin x \leq 0$ and $\cos x \geq 0.$\nl

The folowing diagram tells which identies to use to convert from any trigonometric function to any other trigonometric function.\nl
\begin{center}
\begin{tabular}{c c c c c c c}
$\cos$&$-$&$\sec$&$-$&ID5&$-$&$\tan$\\
$\mid $&&&&&&\\
ID4&&&&&&$\mid$\\
$\mid$&&&&&&\\
$\sin$&$-$&$\text{csc}$&$-$&ID6&$-$&$\text{cot}$\\
\end{tabular}
\end{center}










\textbf{Examp1e 5.3.1.} If $x\in \left[\pi,\dfrac{3\pi}{2}\right] ,\cos x = -\dfrac{3}{5},
y\in  \left[\po{2},\pi\right], \sin y = \dfrac{4}{5},$ find the value of $\sin (x + y)$\nl

Solution. Suppose $x\in \left[\pi,\dfrac{3\pi}{2}\right] ,\cos x = -\dfrac{3}{5},
y\in  \left[\po{2},\pi\right], \sin y = \dfrac{4}{5},$ By ID4,
$\sin^{2}x + \cos^{2}x = 1.$ By Th 5.3.2, $\sin x \leq 0.$   Thus $\sin^{2}x + \dfrac{9}{25} = 1,$ and $\sin x = -\dfrac{4}{5}.$\nl
Similarly, $\cos y = -\dfrac{3}{5}.$ Thus,

\begin{tabular}{r l}
$\sin (x + y)$&$=\sin x \cos y + \cos x \sin y$\\
&$= \left(-\dfrac{4}{5}\right)\left(-\dfrac{3}{5}\right) + \left(-\dfrac{3}{5}\right)\left(\dfrac{4}{5}\right)$\\
&$= 0$\\
\end{tabular}\\

\textbf{Exercises}\nl
1. Simplify $\sin 3x \cos 2x - \cos 3x \sin 2x.$\nl

2· Simplify $\cos \dfrac{x}{3} \cos \dfrac{2x}{3} - \sin \dfrac{x}{3} \sin \dfrac{2x}{3}.$\nl

\newpage
%page85
\lhead{Section 5.4}
3.  If $x\in \left[ \dfrac{3\pi}{2}, 2\pi \right], y\in \left[\pi, \dfrac{3\pi}{2}\right] , \cos x = +\dfrac{5}{13},$ and $\sin y = - \dfrac{4}{5},$ find $\sin (x + y).$\nl

4.  If $x\in \left[ \dfrac{\pi}{2}, \pi \right], y\in \left[\pi, \dfrac{3\pi}{2}\right] , \cos x = -\dfrac{5}{13},$ and $\cos y = - \dfrac{12}{13},$ find $\cos (x + y).$\nl

5.  If $x\in \left[ \pi, \dfrac{3\pi}{2} \right], y\in \left[\dfrac{3\pi}{2}, 2\pi \right] , \sin x = -\dfrac{3}{5},$ and $\cos y = \dfrac{4}{5},$ find $\tan (x + y).$\nl
(Hint : Use ID5 and ID6.)\nl

6.  If $x\in \left[ \dfrac{\pi}{2}, \pi \right], y\in \left[\pi, \dfrac{3\pi}{2}\right] , \csc x = \dfrac{5}{3},$ and $\csc y = - \dfrac{5}{4},$ find $\cos (x - y).$\nl

7.  If $x\in \left[ \dfrac{3\pi}{2}, 2\pi \right], y\in \left[\dfrac{\pi}{2}, \pi \right] , \sin x = -\dfrac{3}{5},$ and $\cos y = - \dfrac{5}{13},$ find $\cot (x - y).$\nl
(Hint : Use Problem 17 of Section 5.2.)\nl

8. Express $\sin 3x + \sin x$ as a product.\nl

9. Prove Theorem 5.3.1(c).\nl

10. Prove Theorem 5.3.1(d).\nl

11. Prove Theorem 5.3.2(a).\nl
\vfill

\textbf{5.4 The Double-and Half-Angle Formulas} \nl

\textbf{Definition 5.4.1.} A function $f$ is said to have \textit{period} $p$ if and only if $p$ is the smallest positive number such that for each $x$ in $D_{f}, (x + p)$ is in $D_{f}$ and $f(x + p) =f(x).$ If $f$ is a function with period $p$ and $a$ is in $R,$ then $\{(x,f(x)) \mid a \leq x \leq a + p\}$ is called a \textit{cycle} of $f.$\nl

\textbf{Theorem 5.4.1.} $\sin, \cos,$ and $\alpha$ have period $2\pi$ and $\tan$ has period $\pi.$\nl
Proof that $\sin$ has period $2\pi.$ By ID27, the period of $\sin$ is less than or equal to $2\pi.$ Suppose $\sin$ has period $p$ and $0< p <2\pi.$ Since $0 <\pi - \dfrac{p}{2} < \pi$ and $\pi < \pi + \dfrac{p}{2} <2\pi$ we have, by assumption and Th 5.3.1,
\[
0< \sin\left(\pi - \dfrac{p}{2}\right) = \sin \left[\left(\pi - \dfrac{p}{2}\right) + p \right] = \sin\left( \pi + \dfrac{p}{2}\right) < 0.
\]
This contradiction proves the theorem.
\newpage

%Page86

\textbf{Theorem 5.4.2.} If $n\in Z,$ then the following are identities.\nl
ID30. $\sin x = \sin(x + 2\pi n)$\nl
ID31. $\cos x = \cos(x + 2\pi n)$\nl
ID32. $\tan x = \tan(x + \pi n)$\nl
Proof of ID30. Let $S = \{n \mid n\in\zp$ and $\sin x = \sin(x + 2\pi n) \}.$\nl
(1) $1\in S$ since $1\in\zp$ and $\sin x = \sin(x +2\pi) $\nl
\begin{tabular}{r l l}
(2) $k\in S$& $\ra k\in\zp $and $\sin x = \sin(x + 2\pi k)$&\\
  & $\ra k+1\in \zp $ and $\sin x = \sin (x+2\pi k) = \sin((x+2\pi k)+2\pi)$&\\
  &\qquad \qquad$ = \sin(x+2\pi(k+1))$& \\
  & $\ra k+1\in S $&\\
\end{tabular}\\
By Th 3.4.2, $S = \zp$ and for each $n$ in $\zp,  \sin x = \sin(x+2\pi n).$   Similarly for each $n$ in $\zp$, $\sin x = \sin(x+2\pi(-n)).$ $  \sin x = \sin(x +2\pi\cdot 0).$ Therefore for each $n \in$ $Z,~ \sin x = \sin(x +2\pi n).$ \nl

\textbf{Theorem 5.4.3.} The following are identities. \nl
ID33. $\sin 2x = 2 \sin x \cos x$\nl
ID34. (a) $\cos 2x = \cos^{2}x - \sin^{2}x$\nl
ID34. (b) $\cos 2x =1 -2\sin^{2}x$\nl
ID34. (c) $\cos 2x =2\cos^{2}x -1 $\nl
ID35. $\cos 3x = 4\cos^{3}x -3\cos x $\nl

ID36. $\tan 2x = \dfrac{2\tan x}{1 - \tan^{2}x}.$\nl

ID37 $\sin \dfrac{x}{2} = \pm\sqrt{\dfrac{1 - \cos x}{2}}$\nl

ID38 $\cos \dfrac{x}{2} = \pm\sqrt{\dfrac{1 + \cos x}{2}}$\nl

ID39  $\tan \dfrac{x}{2} = \dfrac{1 - \cos x}{\sin x} = \dfrac{\sin x}{1 + \cos x}$\nl

Proof of ID35\nl
\begin{tabular}{r l l}
$\cos 3x$&$= \cos(2x + x)$&\\
& $=\cos 2x\cos x - \sin 2x\sin x$ & ID13 \\
& $=(2\cos^{2}x -1)\cos x -2\sin^{2}x\cos x$& ID34(c),ID3\\
& $=2\cos^{3}x - \cos x-2(1 - \cos^{2}x)\cos x$& ID4(b) \\
& $=2\cos^{3}x - \cos x - 2\cos x + 2\cos^{3}x $&\\
& $=4\cos^{3}x - 3\cos x$&\\
\end{tabular}\\

Proof of ID39\nl

$\tan \dfrac{x}{2} = \dfrac{\sin \dfrac{x}{2}}{\cos \dfrac{x}{2}}$\nl

\newpage

%page87

\qquad \qquad $= \dfrac{2\sin \dfrac{x}{2} \cos \dfrac{x}{2}}{2\cos^{2} \dfrac{x}{2}}$\nl

\qquad \qquad $= \dfrac{\sin x}{1 + \cos x}$\nl

\textbf{Exercises}\nl
 1. Prove ID31.\nl
 2. Prove ID32.\nl
 3. Prove ID33.\nl
 4. Prove ID34(a). \nl
 5. Prove ID34(b). \nl
 6. Prove ID34(c). \nl
 7. Prove ID36.    \nl
 8. Prove ID37.    \nl
 9. 	Prove ID38. \nl
 Prove the following are identities. \nl

 10. $\dfrac{1 - \cos 2x}{\sin 2x} = \tan x $\nl

 11. $\dfrac{1 + \sin 2x + \cos 2x}{1 + \sin 2x - \cos 2x} = \cot x$\nl

 12. $(\sin x + \cos x)^{2} = 1 + \sin 2x$\nl
 13. $\tan x + \cot x = 2\csc 2x$\nl
 14. $\sin 3x = 3\sin x -4\sin^{3}x$\nl
 15. 	Prove $\cos$ has period $2\pi.$\nl
 16. 	Prove $\tan$ has period $\pi$.

 \newpage
 %page88
 %Page88
\lhead{Section 5.5}
\textbf{5.5 Finding the Value of a Trigonometric Function of a Number}\nl
\textbf{Theorem 5.5.1.}
\begin{center}
\begin{tabular}{|c |c |c |c |c |c |c |c |c|}\hline
&0&$\po{6}$ &$\po{4}$&$\po{3}$&$\po{2}$&$\pi$&$\dfrac{3\pi}{2}$&$2\pi$\\ \hline
$\sin$&0&$\dfrac{1}{2}$&$\dfrac{1}{\sqrt{2}}$&$\dfrac{\sqrt{3}}{2}$&1&0&-1&0\\ \hline
$\cos$&1&$\dfrac{\sqrt{3}}{2}$&$\dfrac{1}{\sqrt{2}}$&$\dfrac{1}{2}$&0&-1&0&1\\ \hline
$\tan$&0&$\dfrac {1}{\sqrt{3}}$&1&$\sqrt{3}$&$-$&0&$-$&0\\ \hline
\end{tabular}\\
\end{center}
Proof\nl

\begin{tabular}{r l l}
$\cos \po{2}$&$= 4\cos^{3}\po{6} - 3\cos\po{6}$&ID35\\
&&\\
0&$=\cos\po{6}\left(4\cos^{2}\po{6} - 3\right)$&T1\\
&&\\
0&$= 4\cos^{2}\po{6} - 3$ since $\cos \po{6}\neq 0$ &Th 5.3.1\\
&&\\
$\cos \po{6}$&$=\dfrac{\sqrt{3}}{2}$&Th5.3.1\\
&&\\
$\sin^{2}\po{6} + \cos^{2}\po{6}$& $= 1$&ID4\\
&&\\
$\sin^{2}\po{6} + \dfrac{3}{4}$&$ = 1$&\\
&&\\
$\sin\po{6}$&$= \dfrac{1}{2}$& Th 5.3.1\\
&&\\
$\tan \po{6} = \dfrac{\sin \po{6}}{\cos \po{6}}$&$ = \dfrac{\dfrac{1}{2}}{\dfrac{\sqrt{3}}{2}} = \dfrac{1}{\sqrt{3}}$& Def 5.1.5\\
&&\\
%$\sin \po{4}$&$ = \sqrt{\dfrac{1 - \cos \po{2}}{2}}$&ID37\\
%&&\\
%&$= \frac{1}{2}$&T1\\

%$\sin \po{4}$&$= \dfrac{1}{\sqrt{2}}$&Th 5.3.1\\
%&&\\
$\cos\po{4}$&$= \cos\left(\po{2} - \po{4}\right)$&\\

&$= \sin\po{4}$&ID8\\

&$= \dfrac{1}{\sqrt{2}}$&\\

&&\\
$\sin \po{4}$& $= \sqrt{\dfrac{1 - \cos(\pi / 2)}{2}}$&ID37\\
&&\\

&$=\dfrac{1}{\sqrt{2}}$&T1\\

\end{tabular}\\


\newpage
%page89
%\begin{small}

\begin{tabular}{r l r l}
$\tan \po{4}$&$= \dfrac{\sin \po{4}}{\cos \po{4} } = \dfrac{\dfrac{1}{\sqrt{2}}}{\dfrac{1}{\sqrt{2}}} = 1$&$\tan \po{3}$&$= \dfrac{\sin \po{3}}{\cos \po{3} } = \dfrac{\dfrac{\sqrt{3}}{2}}{\dfrac{1}{2}} = \sqrt{3}$\\
&&&\\
$\sin \po{3}$&$= \sin\left(\po{2} - \po{6}\right)$&$\cos \po{3}$&$=\cos\left(\po{2} - \po{6}\right)$\\
&&&\\
&$= \cos \po{6}$&&$= \sin \po{6}$\\
&&&\\
&$= \dfrac{\sqrt{3}}{2}$&&$= \frac{1}{2}$\\

%&&&\\
%$\cos \po{3}$&$=\cos\left(\po{2} - \po{6}\right)$&&\\
%&&&\\
%&$= \sin \po{6}$&&\\
%&&&\\
%&$= \frac{1}{2}$&&\\
\end{tabular}\\
%\end{small}
The remainder of the table follows from T1 and Th 5.2.3. This completes the proof. \nl

%\textbf{Theorem 5.5.2.}   $3 < \pi < 2\sqrt{3}$\nl

To find the value of a trigonometric function of a number not in $\left[0, \po{2}\right],$ use ID30$-$ID32 to reduce the problem to that of finding the value of that trigonometric function of a number in  $\left[0, 2\pi \right].$  Then use ID18$-$ID26 to reduce the problem to that of finding the value of that trigonometric function of a number in $\left[0, \po{2} \right].$\nl

Identities ID18$-$ID26 can be summarized as follows:  The absolute value of a trigonometric function of a number $x$ in $[0, \pi]$ is equal to the value of that trigonometric function evaluated at the number which denotes the distance from  $\alpha (x)$ along the unit circle to the $x$-axis.  The sign of the trigonometric function is given by Theorem 5.3.1.

\textbf{Example 5.5.1.} Find $\sin \dfrac{7\pi}{6}$\nl

Solution. By ID21 and Th 5.5.1, $\sin \dfrac{7\pi}{6} = -\sin\left(\dfrac{7\pi}{6}  -\pi\right)= -\sin\po{6} = -\dfrac{1}{2}$\nl

\textbf{Example 5.5.2.} Find $\cos\left(-\dfrac{2\pi}{3}\right).$\nl

Solution.  By ID10, ID19, and Th 5.5.1,\nl


\[
\cos \left( - \dfrac{2\pi}{3}\right) = \cos \left(\dfrac{2\pi}{3} \right) = -\cos \left(\pi - \dfrac{2\pi}{3}    \right)  = -\cos \left(\po{3}\right) = -\dfrac{1}{2}.
\]

\newpage
%page90
\textbf{Example 5.5.3.} Find $\tan \dfrac{7\pi}{4}.$\nl
Solution. By ID26 and Th 5.5.1,\nl
$\tan \left(\dfrac{7\pi}{4}\right) = -\tan \left(2\pi -\dfrac{7\pi}{4}\right) = - \tan \left(\po{4}\right) = -1.$\nl

\textbf{Definition 5.5.1.} For each $x$ in $R, x\dg,$ read ``$x$ degrees", will denote $\dfrac{x\pi}{180}.$\nl

\textbf{Theorem 5.5.2.} If $y\in R,$ then $y = \left(\dfrac{y~180}{\pi}\right)\dg. $\nl
Proof. Suppose $y\in R.$ By Def 5.5.1,\nl

$\left(\dfrac{y~180}{\pi}\right)\dg   =  \left(\dfrac{y~180}{\pi} \right)\cdot \left(\dfrac{\pi}{180} \right) = y.$\nl

 \textbf{Theorem 5.5.3.} \nl

\begin{tabular}{l l}
a) $0\dg = 0$& e) $90\dg = \po{2}$\\
&\\
b) $30\dg = \po{6}$& f) $180\dg = \pi$\\
&\\
c) $45\dg = \po{4}$& g) $270\dg = \dfrac{3\pi}{2}$\\
&\\
d) $60\dg = \po{3}$& h) $360\dg = 2\pi$\\
\end{tabular}\nl



Proof of (b). $30\dg = \dfrac{30\pi}{180} = \po{6}.$\nl
Proof of (c). $45\dg = \dfrac{45\pi}{180} = \po{4}.$\nl

\textbf{Example 5.5.4.} Convert $\dfrac{2\pi}{3}$ to degrees, i.e., find a number $y$ such that $y\dg = \dfrac{2\pi}{3}.$\nl
Solution. By Th 5.5.3, $\dfrac{2\pi}{3} = \left( \dfrac{2\pi}{3}\cdot\dfrac{180}{\pi}\right)\dg = 120\dg $\nl

\textbf{Example 5.5.5.} Find $\tan 240\dg .$\nl
Solution. By Def 5.5.1 and ID23, $\tan 240\dg = \tan \dfrac{240\pi}{180} =\tan \dfrac{4\pi}{3} = \tan\left( \dfrac{4\pi}{3} - \pi\right) = \tan \po{3} = \sqrt{3}.$\nl

\textbf{Exercises}\nl
1. 	Convert each of the following to degrees.\nl
\begin{tabular}{l l}
&\\
a) $\dfrac{3\pi}{4}$  &  c) $\dfrac{7\pi}{6}$\\
&\\
b) $\dfrac{5\pi}{6}$   &  d) $\dfrac{5\pi}{4}$\\
\end{tabular}

\newpage
%page91
\lhead{Section 5.6}
%Page91

Find the value of each of the following. \nl

2. $\sin 135\dg$ \nl

3. $\cos 240\dg$ \nl

4. $\tan -300\dg$ \nl

5. $\sin -45\dg$ \nl

6. $\cos 300\dg$ \nl

7. $\tan 120\dg$ \nl

Prove each ofthe following. \nl

8. $\sin 15\dg = \dfrac{\sqrt{3} - 1}{2\sqrt{2}}$ (Hint: $15\dg = 45\dg - 30\dg.$) \nl

9. $\cos 15\dg = \dfrac{\sqrt{3} + 1}{2\sqrt{2}}$ \nl

10. $\tan 15\dg = \dfrac{\sqrt{3} - 1}{\sqrt{3} + 1}    $\nl

11 $\sin 75\dg = \dfrac{\sqrt{3} + 1}{2\sqrt{2}}$\nl

12. $\cos 75 = \dfrac{\sqrt{3} - 1}{2\sqrt{2}}$ \nl

13. $\tan 75\dg = \dfrac{\sqrt{3} + 1}{\sqrt{3} - 1}$\nl


14. Prove the rest of Theorem 5.5.4. \nl


15. Prove Theorem 5.5.2. (Hint: Use T3 and Th 5.5.1.) \nl

\vfill

\textbf{5.6 The Sum, Difference, and Product of Sines and Cosines}\nl
\textbf{Theorem 5.6.1.} The following are identities. \nl

ID40. $\sin x \cos y = \frac{1}{2}[\sin (x + y) + \sin(x - y)]$ \nl

ID41. $\cos x \cos y = \frac{1}{2}[\cos (x + y) + \cos(x - y)] $\nl

\newpage
%page92
ID42. $\sin x \sin y = \frac{1}{2}[\cos (x - y) - \cos (x + y)]$\nl

ID43. $\sin x + \sin y = 2\sin \frac{x + y}{2} \cos\frac{x - y}{2}$\nl

ID44. $\sin x - \sin y = 2\cos \frac{x + y}{2} \sin\frac{x - y}{2}      $\nl

ID45. $\cos x + \cos y = 2\cos \frac{x + y}{2}\cos\frac{x - y}{2}$\nl

ID46. $\cos x - \cos y = -2\sin \frac{x + y}{2}\sin\frac{x-y}{2}$\nl

Proof of ID43.\nl

$2\sin\frac{x+y}{2}\cos\frac{x-y}{2}$\nl

$=2\left( \frac{1}{2}\right)[\sin \left(\frac{x+y}{2} + \frac{x-y}{2}\right) + \sin \left(\frac{x+y}{2} - \frac{x-y}{2} \right)]$         ID40\nl

$=\sin x + \sin y$\nl

\textbf{Theorem 5.6.2.}\nl

a) $\sin$ is increasing on $[-\lpo{2}, \lpo{2}]$ and is decreasing on $[\lpo{2}, \frac{3\pi}{2}]$\nl

b) $\cos$ is decreasing on $[0, \pi ]$ and increasing on $[\pi, 2\pi]$\nl

c) $\tan$ is increasing on $\left(-\lpo{2}, \lpo{2}\right)$\nl

Proof of (a).   Suppose $x,y \in [0, \lpo{2}]$ and $y < x.$  Since $0 \leq y < x \leq \lpo{2},$ we have $0 < \frac{x+y}{2} < \lpo{2}$ and $0 < \frac{x-y}{2} \leq \lpo{4} < \lpo{2}.$  By Th 5.3.1, $0 < \cos \frac{x+y}{2}$ and $0 < \sin \frac{x-y}{2}.$\nl
Therefore $\sin x - \sin y = 2\cos \frac{x + y}{2}\sin\frac{x - y}{2}> 0,$ by ID44.  Thus $\sin y < \sin x.$ Hence $\sin$ is increasing on $[0, \lpo{2}].$\nl

Suppose $x,y \in [-\lpo{2}, 0]$ and $x < y.$  Thus $-\lpo{2} \leq x < y \leq 0$ and $\lpo{2} \geq -x > -y \geq 0.$  Since $\sin$ is increasing on $[0, \lpo{2}],$ we have $\sin (-x) > \sin (-y).$ By ID11, we see that $-\sin x > -\sin y$ and that $\sin x < \sin y.$  Therefore $\sin$ is increasing on $[-\lpo{2}, \lpo{2}].$\nl

Suppose $x,y \in [\lpo{2}, \frac{3\pi}{2}]$ and  $x < y.$  Thus $\lpo{2} \leq x < y \leq \frac{3\pi}{2}$ and
$-\lpo{2} \leq x - \pi  < y - \pi  \leq \lpo{2}.$  Since $\sin$ is increasing on $[-\lpo{2}, \lpo{2}],$ we have $\sin (x - \pi) < \sin (y - \pi).$  Thus by ID21, $-\sin x < -\sin y.$  Hence $\sin x > \sin y.$  Therefore the theorem is proven.\nl

From Th 5.5.1, Th 5.6.2, ID18, ID21, ID24, and ID30, we can obtain the following graph of a portion of $\sin .$\nl

%Similarly, we obtain the following graphs.\nl


\newpage
%page93
%From Th 5.5.1, Th 5.6.2, ID18, ID21, ID24, and ID30, we can obtain the following graph of a portion if $\sin .$\nl

Similarly, we obtain the following graphs.\nl

\begin{figure}[h!]
	\centerline{\includegraphics[scale=1.0]{Page93.eps}}
%	\caption{Desert Phoo.}
%	\caption{Top-Left}
%	\label{Fi: contractibility}
\end{figure}




%\textbf{The graphs on page 93 go here.}\nl
\newpage
%page94
\textbf{Example 5.6.1.} Graph one cycle of $f(x) = \frac{1}{2}\sin\left(2x - \po{3}\right).$\nl

 Solution. $f$ obtains a maximum height of $\frac{1}{2}$ and a minimum height of $- \frac{1}{2}.$   Since $\sin$ has a period of $2\pi, f$ starts a cycle when $2x - \po{3} = 0.$ i.e., when
$x = \po{6}$ and $f$ completes a cycle when $2x - \po{3} = 2\pi,$ i.e., when $x = \dfrac{7\pi}{6}.$   Each cycle of $f$ has the same shape as a corresponding cycle of $\sin.$

%\textbf{The graph on page 94 goes here.}\nl

\begin{figure}[h!]
	\centerline{\includegraphics[scale=0.7]{Page94.eps}}
%	\caption{The Set-Up for the Section on Contractibility.}
%	\label{Fi: contractibility}
\end{figure}

\textbf{Exercises}\nl
 Graph one cycle of each of the following. \nl
1. $f(x) = 2\sin x$ \nl

2. $f(x) = 3\cos 2x$ \nl

3. $f(x) = -2\sin \pi x$ (Hint: Draw $g(x)=2\sin \pi x$ upside down.)\nl

4. 	$f( x) = 2\cos \frac{x}{2}$­\nl

5. 	$f(x) = \sin(-x)$ \nl

6. $f(x) = 3\sin(-2x)$\nl

7. Prove ID40. \nl

8. 	Prove ID41.\nl

9. 	Prove ID42. \nl

10. Prove ID44.\nl

11. Prove ID45. \nl

12. Prove ID46. \nl

13. Prove $\cos$ is decreasing on $[0,\frac{\pi}{2}].$ \nl

14. Prove $\tan$ is increasing on $[0, \frac{\pi}{2}].$ \nl

\newpage
\lhead{Section 5.7}
%page95
\textbf{5.7 Inverse Trigonometric Functions}\nl

\textbf{Example 5.7.1.} Find the truth set of $\sin x = \frac{1}{2}.$\nl

Solution. From the graph of $\sin,$ we see that there is one element of
the truth set in $[0,\frac{\pi}{2}],$ one in $[\frac{\pi}{2}, \pi],$ and none in $[\pi, 2\pi].$ By Th 5.5.1,
we have $\sin \frac{\pi}{6} = \frac{1}{2}$ and by ID18, $\frac{1}{2} = \sin\frac{\pi}{6} = \sin(\pi - \frac{\pi}{6}) = \sin\frac{5\pi}{6}. $  Since $\sin$ has period $2\pi,$ the truth set of $\sin x = \frac{1}{2}$
is $\{x \mid x= \frac{\pi}{6} + 2\pi n$ or $x = \frac{5\pi}{6} + 2\pi n$ for
some integer $n\}.$\nl

There are many instances in which it is desirable to find the inverse of a trigonometric function. Since none of the trigonometric functions are one-to-one, we must restrict their domain to a set on which they are one-to-one before taking their inverse. \nl

\textbf{Definition 5.7.1.} \nl
a) 	Let $\arcsin$ or $\sin^{-1}$ denote the inverse of $\{(x,y) \mid x\in [-\frac{\pi}{2}, \frac{\pi}{2}]$ and $y = \sin x\}$,
i.e., $y = \arcsin x \lra y \in [-\frac{\pi}{2}, \frac{\pi}{2}] $ and $x = \sin y.$ \nl

b) 	Let $\arccos$ or $\cos^{-1}$ denote the inverse of $\{(x,y) \mid x\in [0, \pi]$ and $y = \cos x\}$,
i.e., $y = \arccos x \lra y \in [0, \pi] $ and $x = \cos y.$ \nl

c) 	Let $\arctan$ or $\tan^{-1}$ denote the inverse of $\{(x,y) \mid x\in (-\frac{\pi}{2}, \frac{\pi}{2})$ and $y = \tan x\}$,
i.e., $y = \arctan~x \lra y \in (-\frac{\pi}{2}, \frac{\pi}{2}) $ and $x = \tan y.$ \nl

d) Let arccsc or $\csc^{-1}$ denote the inverse of $\{(x,y) \mid x\in  [-\lpo{2},0)\cup (0,\lpo{2}]$ and $y = \csc x\},$ i.e., $y = \text{arccsc~}x \lra y\in [-\lpo{2},0)\cup (0,\lpo{2} ]$ and $x = \csc~y.$ \nl

e) Let arcsec or $\sec^{-1}$ denote the inverse of $\{(x,y) \mid x\in [0,\lpo{2})\cup (\lpo{2}, \pi]$ and $y =\sec x \},$ i.e., $y = \text{arcsec}~x \lra  y \in [0,\lpo{2})\cup (\lpo{2}, \pi]$ and $x = \sec~y.$ \nl

f) Let arccot or $\cot^{-1}$ denote the inverse of $\{(x,y) \mid x \in (0,\pi)$ and $y = \text{cot} ~x \},$ i.e., $y = $ arccot $x \lra y \in (0,\pi)$ and $x = \cot~y.$ \nl

\textbf{Example 5.7.2.} Find $\arcsin (-\frac{1}{2}).$\nl
Solution. Let $x = \arcsin(-\frac{1}{2}).$ By Def 5.7.1, $x = \arcsin -\frac{1}{2} \lra$  $\sin x = -\frac{1}{2}$ and $x\in [-\lpo{2}, \lpo{2}].$   By Th 5.5.1 and ID11, we have $- \frac{1}{2} = -\sin\lpo{6}=
\sin-\lpo{6}$ and $-\lpo{6} \in [-\lpo{2}, \lpo{2}].$   Thus $x = -\lpo{6}.$\nl

\textbf{Example 5.7.3.}  Find the value of $\sin\left(\arccos (-\frac{3}{5}) + \arcsin \frac{4}{5}\right).$\nl

\newpage
%page96
Solution. Let $x = \arccos (-\frac{3}{5}) $ and $ y = \arcsin \frac{4}{5}.$  By Def 5.7.1, $x = \arccos \left(-\frac{3}{5}\right) \lra
 \cos x = \left(-\frac{3}{5}\right)$ and $x \in  [0,\pi],$
 and $y = \arcsin \frac{4}{5} \lra \sin y = \frac{4}{5}$ and $y\in  [-\lpo{2}, \lpo{2}].$\nl
 By ID4 and Th 5.3.2, we have $\sin x = \sqrt{1 - \cos^{2}x} = \sqrt{1 - \frac{9}{25}} = \frac{4}{5}$ and $\cos y =
\sqrt{1 - \sin^{2}y} = \sqrt{1 - \frac{16}{25}} = \frac{3}{5}.$ Thus $\sin\left(\arccos (-\frac{3}{5}) + \arcsin \frac{4}{5}\right) = \sin(x + y)$ $= \sin x \cos y + \cos x \sin y = \frac{4}{5}\frac{3}{5} + (\frac{-3}{5})(\frac{4}{5}) = 0.$\nl

\textbf{Example 5.7.4.} Find the value of $\sin\left(\arctan\left(-\frac{5}{12}\right)\right) $\nl

Solution. Let $x = \arctan\left(-\frac{5}{12}\right).$ By Def 5.7.1, $x = \arctan\left(-\frac{5}{12}\right) \lra \tan x = -\frac{5}{12}$ and $x \in (-\lpo{2}, \lpo{2}).$   By ID6 and Th 5.3.2, we have $\sin\left(\arctan\left(-\frac{5}{12}\right)\right) = \sin x$\nl
 $= \dfrac{1}{\csc x} = \dfrac{1}{-\sqrt{1 + \cot^{2}}} = \dfrac{1}{-\sqrt{1 + \frac{144}{25}}} = -\dfrac{5}{13}.$

\textbf{Exercises}\nl
Find the truth set of each of the following. \nl

1. $\sin x = \frac{1}{\sqrt{2}}$ \nl

2. $\tan x= -\sqrt{3}$\nl

3. $\sec x =\frac{2}{\sqrt{3}}$ \nl		

4. $\cot x= -1$ \nl

5. $\csc x = 2$\nl
 Find the value of each of the following. \nl

6. arctan$(-\sqrt{3})$ \nl

7. $\sin^{-1} (\cos(-60\dg))$ \nl

8. $\tan(\sin^{-1}(-\frac{1}{\sqrt{2}}))$ \nl

9. $\sin(2\arccos \frac{3}{5})$ \nl

10. $\tan\left( \frac{1}{2}\arccos (\frac{-5}{13})\right)$

\newpage

%page97

11. $\sin\left(\sin^{-1} \frac{1}{2} + \cos^{-1}(-\frac{4}{5})\right) $\nl

12. $\cos(\arcsin x + \arccos y)$\nl

13. $\tan(\tan^{-1} x + \tan^{-1} y)$\nl

14. $\sin^{2} (\frac{1}{2}\cos^{-1} x)$\nl


\newpage
%page98
\lhead{Section 6.1}
\textbf{Chapter 6}\nl

\textbf{Complex Numbers and Vector Spaces}\nl

\textbf{6.1 Complex Numbers }\nl

\textbf{Definition 6.1.1.} $(C,+,\cdot)$ is called the \textit{complex number system} and the elements of $C$ are called \textit{complex numbers} if and only if \nl
a) $R\subseteq C$\nl
b) $i\in C$ and $i^{2}=-1$\nl
c) $(C,+,\cdot)$ is a field \nl
d) For each $x$ in $C,$ there exist a real number $a$ (called the \textit{real
part} of $x$ and denoted by $Re~x$) and a real number $b$ (called the \textit{imagin­ary part} of $x$ and denoted by $Im~ x$) such that $x = a + ib.$ \nl

\textbf{Definition 6.1.2.} If $a,b\in R,$ and $x = a + ib,$ then $a - ib$ is called the \textit{complex conjugate} of $x$ and denoted by $\overline x.$\nl

\textbf{Theorem 6.1.1.} If $x,y\in C,$ then \nl
a) $\overline{i}= -i $\nl
b) $\overline{\overline{x}} = x $\nl
c) $\overline{x} = x$ if and only if $x =Re~x $\nl
d) $\overline{x}= -x$ if and only if $x =iIm~x $\nl
e) $x + \overline{x} = 2Re~x$\nl
f) $x - \overline{x}= 2iIm~x$\nl
g) $x\overline{x}=(Re~x)^{2} + (Im~x)^{2}$\nl
h) $\overline{x + y} = \overline{x} + \overline{y}$\nl
i) $\overline{xy} = \overline{x}~\overline{y} $\nl
j) $\overline{x - y} = \overline{x} - \overline{y}$\nl

k) $\overline{\left(\dfrac{x}{y}\right)}  = \dfrac{\overline{x}}{\overline{y}}$ if $y \neq 0$\nl

l) $x = y$ if and only if $Re~x=Re~y$ and $Im~x =Im~y.$\nl
m) $x\in R$ if and only if $Im~x = 0$\nl

\textbf{Defniition 6.1.3.} If $x\in C,$ then $\sqrt{x\overline{x}}$ is called the absolute value of $x$ and denoted by $\abs{x}.$ \nl

\textbf{Theorem 6.1.2.} If $x,y\in C,$ then\nl
a) $\abs{x} = \sqrt{(Re~x)^{2} + (Im~x)^{2}} = \abs{\overline{x}}$\nl
b) $\abs{x} \geq Re~x$\nl

\newpage
%page99

c) $\abs{xy} = \abs{x}\abs{y}$\nl
d) $\abs{x + y} \leq \abs{x} + \abs{y}$\nl

e) $\abs{\dfrac{x}{y}} = \dfrac{\abs{x}}{\abs{y}}$ if $y \neq 0$\nl
Proof of (d).  Suppose $x,y\in C.$  Thus\nl
\begin{tabular}{r l}
$\abs{x + y}^{2}$&$= (x + y)(~\overline{x + y}~)$\\
&$=(x + y)(\overline{x} + \overline{y})$\\
&$=x\overline{x} + x\overline{y} + y\overline{x} + y\overline{y}$\\
&$=x\overline{x} + x\overline{y} +\overline{x\overline{y}} +y\overline{y}$\\
&$=\abs{x}^{2} + 2Re~(x\overline{y}) + \abs{y}^{2}$\\
&$\leq \abs{x}^{2} + 2\abs{x\overline{y}} + \abs{y}^{2}$\\
&$=\abs{x}^{2} + 2\abs{x}\abs{y} + \abs{y}^{2}$\\
&$=(\abs{x} + \abs{y})^{2}$\\
\end{tabular}\\

Therefore $\abs{x + y} \leq \abs{x} + \abs{y}.$\nl

\textbf{Definition 6.1.4.}  If $a,b\in R$ and $x = a + ib,$ let $(a, b)$ denote $x$.\nl

\textbf{Theorem 6.1.3.}  If $t,a,b,c,d\in R, x = (a, b),$ and $y = (c, d),$ then\nl
a) $x + y = (a+c, b+d)$\nl
b) $xy = (ac - bd, ad + bc)$\nl
c) $tx = (ta, tb)$\nl

\textbf{Theorem 6.1.4.}  If $x\in C,$ and $x \neq 0,$ then there is a positive number $r$ and a number $\theta$ in $[0, 2\pi)$ such that $x = r\alpha (\theta).$\nl
Proof.  Suppose $x$ is in $C$ and $x\neq 0.$  Let $r = \abs{x}.$  Since\nl
\[
d\left(\dfrac{x}{\abs{x}}, 0 \right) = \sqrt{\dfrac{(Re~x)^{2}}{\abs{x}^{2}} + \dfrac{(Im~x)^{2}}{\abs{x}^{2}}}
\]
\[
=\sqrt{\dfrac{(Re~x)^{2} + (Im~x)^{2}}{\abs{x}^{2}}} = \sqrt{\dfrac{\abs{x}^{2}}{\abs{x}^{2}}} = 1,
\]\nl
we see that $\dfrac{x}{\abs{x}}\in c_{1}.$  By th 5.1.1, there is a number $s$ such that $\alpha(s) = \dfrac{x}{\abs{x}} = \dfrac{x}{r}.$\nl
Suppose $s \geq 2\pi.$  Let $M = \{n \mid n\in \zp \text{~and~} \dfrac{s}{2\pi} - 1 < n \}.$  Let $m$ denote the least element of $M$.  Thus $m - 1 \leq \dfrac{s}{2\pi} - 1 < m$ and $0 \leq s - 2\pi m < 2\pi.$
By Th 5.4.2, $\alpha(s) = \alpha (s - 2\pi m).$   If $s < 0,$ then by a similar argument there is a number $t$ in $[0, 2\pi)$ such that $\alpha(s) = \alpha(t).$  Therefore there is a number $\theta$ in $[0, 2\pi)$ such that $\alpha (\theta) = \alpha(s) = \dfrac{x}{r}.$  Thus $x = r\alpha(\theta)$ and $\theta \in [0, 2\pi).$

\newpage
%page100

\textbf{Example 6.1.1.}  Express $(2 - 5i) - (3 - 2i)$ as an ordered pair.\nl
Solution. $(2-5i) -(3-2i)= 2 - 5i - 3 + 2i = -1 - 3i =(-1,-3) $\nl

\textbf{Example 6.1.2.} Express $(2 - 3i)(2 + 3i)$ as an ordered pair\nl
 Solution. $(2 - 3i)(2 + 3i)=4 + 6i - 6i -9i^{2} = 4 + 9 =(13,0)$\nl



\textbf{Example 6.1.3.} Express $\dfrac{3 + 2i}{3 - 2i}$ as an ordered pair.\nl
\[
\text{Solution.~} \dfrac{3 + 2i}{3 - 2i} = \dfrac{3 + 2i}{3 - 2i}\cdot \dfrac{3 + 2i}{3 + 2i} = \dfrac{9 + 6i +6i + 4i^{2}}{9 + 6i -6i - 4i^{2}}
\]
\[
=\dfrac{9 + 12i - 4}{9 + 4} = \dfrac{5 + 12i}{13} = \left(\dfrac{5}{13}, \dfrac{12}{13}\right).
\]

\textbf{Exercises}\nl

Express each of the following as an ordered pair.

1. 	$(3 - 2i) + (4 + 3i)$\nl
2. 	$(3i - 5) + (4 - 2i)$\nl
3. 	$(6 + 2i) - (2 + 4i)$\nl
4. 	$(2 + i)(2 - i)$\nl
5. 	$(3 - 2i)(2 + i)$\nl

6. 	$\dfrac{3 + 2i}{5 - 3i}$\nl

7.  $\dfrac{3}{-i}$\nl

8. 	Prove Theorem 6.1.1 (a), (b), and (c). \nl

9. 	Prove Theorem 6.1.1 (d) and (e).

10. Prove Theorem 6.1.1 (f) and (g).

11. Prove Theorem 6.1.1 (h) and (i).

12. Prove Theorem 6.1.1 (j) and (k).

13. Prove Theorem 6.1.1 (l) and (m).

14. Prove Theorem 6.1.2 (a) and (b).

15. Prove Theorem 6.1.2 (c) and (d).


\newpage
%page101
\lhead{Section 6.2}
\textbf{6.2 Generalized Exponents}\nl

\textbf{Definition 6.2.1.} Let $L$ denote the function with domain $(C-\{0\})$ such that if $r\in P, \theta\in[0, 2\pi),$ and $x =r\alpha (\theta) = r(\cos \theta, i~\sin \theta),$ then $L(x) = (L(r),\theta).$ \nl

\textbf{Definition 6.2.2.} Let $E$ denote the function defined on $C$ such that if $a,b\in R$ and $x = a + ib,$ then $E(x) = E(a)\alpha(b).$\nl

The same symbol was used for the $L$ defined in Section 4.2 as for the $L$ defined above. This should cause no confusion since the former $L$ is a subset of the latter one. A similar statement is true for $E.$\nl

\textbf{Theorem 6.2.1.} If $r\in P, \theta \in  [0,2\pi),$ and $x =r\alpha(\theta),$ then $E(L(x))=x.$\nl
 Proof. Suppose $r\in P, \theta \in  [0,2\pi)$ and $x =r\alpha(\theta).$ Thus
 \[
 E(L(x)) = E(L(r\alpha(\theta)))=E(L(r),\theta) = E(L(r))\alpha(\theta) = r\alpha(\theta) = x
 \]

\textbf{Theorem 6.2.2.} If $a,b\in R$ and $x = a + ib,$ then $L(E(x))=x.$\nl
Proof. Suppose $a,b\in R$ and $x = a + ib.$
\[
L(E(x)) = L(E(a)\alpha (b)) = (L(E(a)), b) = (a,b) = x
\]

\textbf{Definition 6.2.3.} If $a\in (C-\{ 0, 1 \}),$ let $log_{a}$ denote the function with domain $(C-\{0\})$ such that $log_{a}x=L(x)/L(a)$ for each $x$ in $(C-\{0\}).$\nl

\textbf{Definition 6.2.4.} If $a\in (C-\{0\})$ and $x\in C,$ let $a^{x}$ denote $E(xL(a)).$\nl

\textbf{Theorem 6.2.3.} If $a,b\in R$ and $x = a+i·b,$ then $e^{x} =e^{a}(\cos b +i\sin b).$\nl

\textbf{Theorem 6.2.4.} If $r\in P, \theta \in  [0,2\pi),$ and $x = r(\cos \theta + i\sin \theta),$ then
$log_{e}x = (log_{e}r, \theta).$\nl

\textbf{Theorem 6.2.5.} All of the theorems of Sections 4.2 and 4.3 are true with $R$ replaced by $C$ and $P$ replaced by $(C - \{0\})$ with the excep­tion of Theorems L2, 4.2.2(d) and 4.2.3.\nl

\textbf{Theorem 6.2.6.} If $r,s\in P, \theta,\phi \in [0,2\pi), x = r(\cos \theta + i\sin \theta),$ and $y = s(\cos \phi + i\sin \phi ),$ then $xy = rs[\cos( \theta + \phi) + i\sin( \theta + \phi )].$\nl

Proof. Suppose  $r,s\in P, \theta,\phi \in [0,2\pi), x = r(\cos \theta + i\sin \theta),$ and $y = s(\cos \phi + i\sin \phi ).$\nl
\begin{tabular}{l l}
$xy = r\left(\cos\theta + i\sin\theta\right)\cdot s\left(\cos\phi + i\sin\phi\right)$&\\
$=rs[\cos\theta\cos\phi - \sin\theta\sin\phi + i\left(\sin\theta\cos\phi + \cos\theta\sin\phi\right)]$&Theorem 6.1.3.\\
$=rs[\cos(\theta + \phi) + i\sin(\theta + \phi)]$&Theorem5.2.1,ID13a and ID14\\
\end{tabular}\\

\newpage
%page102



\indent $=E(L(rs), \theta + \phi)$\nl
\indent $=E(L(rs))\alpha ( \theta + \phi)$\nl
\indent $=rs(\cos( \theta + \phi) + i\sin( \theta + \phi))$\nl

\textbf{Theorem 6.2.7.} If $r,s\in P, \theta,\phi  \in [0,2\pi), x=r(\cos \theta  + i\sin \theta),$ and
$y = s(\cos \phi + i\sin \phi),$ then $\dfrac{x}{y} = \dfrac{r}{s}(\cos( \theta - \phi) + i\sin( \theta - \phi )).$\nl

\textbf{Theorem 6.2.8.} If $n\in R, r\in P, \theta \in [0,2\pi),$ and $x = r(\cos \theta + i\sin \theta),$ then $x^{n} = r^{n}(\cos n\theta + i\sin n\theta).$\nl
Proof. Suppose  $n\in R, r\in P, \theta \in [0,2\pi),$ and $x = r(\cos \theta + i\sin \theta).$\nl
\begin{tabular}{r l}
$x^{n}$&$=(r\alpha(\theta))^{n}$\\
&$=E(nL(r\alpha (\theta)))$\\
&$=E(n(L(r), \theta))$\\
&$=E(nL(r),n\theta ) $\\
&$=E(L(r^{n}),n\theta )$\\
&$=r^{n}\alpha(n\theta)$\\
&$=r^{n}(\cos n\theta + i\sin n\theta)$\\
\end{tabular}\\

\textbf{Definition 6.2.5.} If $n\in \zp$ and $a,x\in C,$ then $a$ is called an $n$th root of $x$ if and only if $a^{n} = x.$\nl

\textbf{Theorem 6.2.9.} If $n\in \zp, k\in Z_{0}^{n-1}, r\in P, \theta \in [0,2\pi),$ and $x = r(\cos \theta + i\sin \theta),$ then\nl
\[
r^{1/n} \left[\cos\left( \dfrac{\theta}{n}+ \dfrac{k2\pi}{n}\right) + i\sin\left( \dfrac{\theta}{n} + \dfrac{k2\pi}{n}\right)\right] \text{~is~ an~} n\text{th~root~of~} x.
\]
Proof. Suppose $n\in \zp, k\in Z_{0}^{n-1}, r\in P, \theta \in [0,2\pi),$ and $x = r(\cos \theta + i\sin \theta).$\nl
Thus, $\left(\dfrac{\theta}{n} + \dfrac{k2\pi}{n}\right)\in [0, 2\pi)$ and
\[
 \left[r^{1/n} \left[\cos\left( \dfrac{\theta}{n}+ \dfrac{k2\pi}{n}\right) + i\sin\left( \dfrac{\theta}{n} + \dfrac{k2\pi}{n}\right)\right]\right]^{n}
\]
$= \left(r^{1/n}\right)^{n}\left[\cos(\theta + k2\pi) + i\sin(\theta + k2\pi)\right]$\nl
%$=(r^{1/n})^{n}[\cos(\theta + k2\pi) + i\sin(\theta + k2\pi)]$\nl
$=r(\cos\theta + i\sin \theta) = x$\\

\textbf{Example 6.2.1.} Find the value of $i^{i}.$ Solution.
$i^{i}=E(iL(i)) = E\left(iL\left(1\cdot\alpha\left(\po{2}\right)\right)\right) = E\left(i\left(L(1),\po{2}\right)\right)$ \nl
$=E\left(i\left(0,\po{2}\right)\right) = E\left(ii\po{2}\right) = E\left(-\po{2}\right) = e^{-\lpo{2}}$\nl


\textbf{Example 6.2.2.} Express $\dfrac{10(\cos 120\dg + i\sin 120\dg)}{2(\cos 30\dg + i\sin 30\dg)}$ as an ordered pair.

\newpage
%page103

Solution.\nl
\[
\dfrac{10(\cos 120\dg + i\sin 120\dg)}{2(\cos 30\dg + i\sin 30\dg)} = 5\left(\cos 90\dg + i\sin 90\dg\right) = 5i = (0, 5)
\]

\textbf{Exercises} \nl

1. 	If $x = 3(\cos 15\dg + i\sin 15\dg)$ and $y = 2(\cos 45\dg + i\sin 45\dg),$ find each
of the following. \nl

a) $xy$ \quad b) $x^{4}$ \quad c) $\dfrac{y}{x}$ \quad d) $y^{-2}$\nl

Express each of the following as an ordered pair. \nl

2. 	$(2+2i)^{2}$\nl
3. 	$(\sqrt{3} - i)^{2}$\nl
4. 	$\left[3\left(\cos \po{3} + i\sin \po{3}\right)\right]^{5}$ \nl

5. 	$\dfrac{6(\cos(\frac{3\pi}{4} +i\sin \frac{3\pi}{4}))}{3(\cos(\frac{\pi}{2} +i\sin \frac{\pi}{2}))} $\nl

6. 	$[4(\cos 135\dg + i\sin 135\dg)] [3(\cos 45\dg + i \sin 45\dg)]$\nl

7. 	$[5(\cos \frac{7\pi}{6} + i\sin \frac{7\pi}{6}][2(\cos \lpo{3} + i\sin \lpo{3}] $\nl

\begin{large}
8.  $\left(\dfrac{\sqrt{3}}{2}, \dfrac{1}{2}\right)^{(3,6)}$\nl

9.  $\left(\dfrac{1}{\sqrt{2}}, -\dfrac{1}{\sqrt{2}}\right)^{\left(\frac{8}{7},0\right)}$\nl

10.  $(1 + i)^{-i}$\nl

\end{large}


11. Find three cube roots of $27^{i}.$\nl

12. Find three cube roots of $-i.$\nl

13. Find four fourth roots of $-1 + \sqrt{3}i.$\nl

14. Find three cube roots of $-1 + i$\nl

15. Prove Theorem 6.2.3. \nl

16. Prove Theorem 6.2.4. \nl

17. Prove Theorem 6.2.7. \nl

%18. Prove Theorem 6.2.7. \nl

\newpage
%page104
\lhead{Section 6.3}
\textbf{6.3 Vector Spaces}\nl

\textbf{Definition 6.3.1.} $(V,F,+,\cdot)$ is said to be a \textit{vector space} if $(F,+,\cdot)$ is a field, $V$ is a non-empty set whose elements are called vectors, $+$ is a binary operation over $V,~ \cdot$ is a function whose domain contains $F\times V,$ and the following statements are true. \nl

\qquad V1. If $x,y\in V,$ then $x + y\in V.$ \nl

\qquad V2. If $x,y\in V,$ then $x + y = y + x.$ \nl

\qquad V3. If $x,y,z\in V,$ then $(x + y) + z = x + (y + z).$ \nl

\qquad V4. There is an element $0$ in $V$ (called the zero vector) such that
if $x\in V,$ then $x + 0 = x.$ \nl

\qquad V5. If $x\in V,$ then there is an element $-x$ in $V$ (called the inverse
of $x$) such that $x + (-x) = 0.$ \nl

\qquad V6. If $a\in F$ and $x\in V,$ then $a\cdot x\in V.$ \nl

\qquad V7. If $a,b\in F$ and $x\in V,$ then $(ab)\cdot x = a(b\cdot x).$ \nl

\qquad V8. If $a\in F$ and $x,y\in V,$ then $a\cdot(x + y) =a\cdot x + a\cdot y.$ \nl

\qquad V9. If $a,b\in F$ and $x\in V,$ then $(a + b)\cdot x = a\cdot x + b\cdot x.$ \nl

\qquad V10. If $x\in V,$ then $1\cdot x = x,$ where $1$ is in $F$ and $1$ is the identity for $\cdot.$ \nl

\textbf{Definition 6.3.2.} If $(V,F,+,\cdot)$ is a vector space, $a,b\in F$ and $x,y\in V,$ we will use $ab,~ ax,$ and $x - y$ to denote $a\cdot b, a\cdot x,$ and $x + (-y)$ respectively. \nl

\textbf{Theorem 6.3.1.} If $(V,F,+,\cdot )$ is a vector space, then the following statements are theorems. \nl
a) If $x,y,z\in V$ and either $x + z = y + z$ or $z + x = z + y,$ then $x = y.$\nl
b) If $x,y\in V$ and $x + y = x,$ then $y = 0.$\nl
c) If $x,y\in V$ and $x + y =0,$ then $y = -x.$\nl
d) If $x\in V,$ then $0x = 0.$\nl
e) If $a\in F,$ then $a0 = 0.$\nl
f) If $a\in F, x\in V,$ and $ax = 0,$ then $a = 0$ or $x = 0.$\nl
g) If $x\in V,$ then $-x = (-1)x.$\nl
h) If $a\in F$ and $x\in V,$ then $(-a)x = a(-x) = -(ax).$\nl

\textbf{Definition 6.3.3.} $(V,F,+,\cdot)$ is called an \textit{inner product} space if and only if $(V,F,+,\cdot)$ is a vector space, the domain of $\cdot$ contains $V\times V,$ and if $a\in F$ and $x,y,z\in V,$ then $x\cdot y\in R$ and \nl
a) $x\cdot x \geq 0$\nl
%\newpage
%page105
b) $x\cdot x = 0$ if and only if $x = 0.$ \nl
c) $x\cdot y = y\cdot x$\nl
d) $a(x\cdot y) =(ax)\cdot y $\nl
e) $x\cdot (y + z) = x\cdot y + x\cdot z$\nl

Note: $(V,+,\cdot)$ is a vector space means $(V,R,+,\cdot)$ is a vector space, where $(R,+,\cdot)$ is the real
number system. \nl

\textbf{Definition 6.3.4} If $n\in \zp,$ then $x$ is called an $n$\textit{-tuple of elements} of $R$ if and only if $x : Z_{1}^{n} \ra R.~$  If $n\in\zp$ and $x$ is an $n$-tuple, then for each $i$ in $Z_{1}^{n}, x(i)$ will be denoted by $x_{i}$ and $x$ will be denoted by $(x_{1} , x_{2},~ \dots ~,x_{n}).$ If $n\in\zp,$ let $R^{n}$ denote the set of all $n$-tuples. \nl

\textbf{Theorem 6.3.2.} If $n\in\zp,$ then $(R^{n},+,\cdot)$ is an inner product space
with $~\cdot~$ and $~+~$ defined such that for each $x,y\in R^{n}$ and $t\in R,$ we have\nl

$x + y=(x_{1} + y_{1}, x_{2} + y_{2} , ~\dots ~ ,x_{n} + y_{n}),$\nl

$t\cdot x = (tx_{1}, tx_{2}, ~\dots ~ ,tx_{n}),$ and
\[
x\cdot y = \sum\limits_{i=1}^{n}x_{i}y_{i}.
\]

\textbf{Exercises} \nl

1. Prove Theorem 6.3.1(a) and (b). \nl

2. Prove Theorem 6.3.1(c) and (d).\nl

3. Prove Theorem 6.3.1(e) and (f).\nl

4. Prove Theorem 6.3.1(g) and (h).\nl

5. Prove Theorem 6.3.2. \nl

\newpage
%page106
\lhead{Section 6.4}
\textbf{6.4 Measure of an Angle}\nl

\textbf{Definition 6.4.1.} $(V,F,+,\cdot)$ is said to be a \textit{normed linear} space if and only if $(V,F,+,\cdot)$ is a vector space and for each $x$ in $V$ there is a real num­ber $\abs{x}$ (called the \textit{norm} of $x$) such that \nl

a) $\abs{x} \geq 0$\nl
b) $\abs{x} = 0$ if and only if $x = 0$\nl
c) $\abs{a\cdot x} = \abs{a}\abs{x}$ for every $a$ in $F$ \nl
d) $\abs{x + y} \leq \abs{x} + \abs{y}$ for each $y$ in $V$.\nl


\textbf{Theorem 6.4.1.} If $(V,F,+,\cdot)$ is an inner product space, then $(V,F,+,\cdot )$ is a normed linear
space, with $\abs{x}= \sqrt{x\cdot x}$ for each $x$ in $V.$ \nl

\textbf{Definition 6.4.2.} $(S,d)$ is called a \textit{metric space} if and only if $S$ is a non-empty set and $d : S\times S\ra R$ such that if $x,y,z\in S,$ then \nl
a) $d(x,y) \geq 0$\nl
b) $d(x,y) = 0$ if and only if $x = y$\nl
c) $d(x,y) = d(y,x)$\nl
d) $d(x,z) \leq d(x,y) + d(y,z)$\nl

\textbf{Theorem 6.4.2.} If $(V,F,+,\cdot)$ is a normed linear space and $d(x,y) = \abs{x - y}$ for each $x$ and $y$ in $V,$ then $(V,d)$ is a metric space. \nl

\textbf{Definition 6.4.3.} Suppose $(V,+,\cdot)$ is a vector space, $x,y\in V$ and $x \neq y.$\nl

a) The \textit{line} $xy$ is denoted by $\overleftrightarrow{xy}$ and defined by $\overleftrightarrow{xy} = \{p \mid p = ay + (1-a)x\text{~for~some~} a \text{~in~} R\}.$ \nl

b) The \textit{ray} $xy$ is denoted by $\overrightarrow{xy}$ and defined by $\overrightarrow{xy} = \{p \mid p =ay +( 1 - a)x$ for some $a\geq 0\}.$ $x$ is called the \textit{end point} of $\overrightarrow{xy}.$ \nl

c) 	The \textit{segment} $xy$ is denoted by $\overline{xy}$ and defined by $\overline{xy}  =\{p \mid p = ay + (1 - a)x$ for some $a$ in $[0,1] \}.$ $x$ and $y$ are called \textit{end points} of $\overline{xy}.$ \nl

\textbf{Definition 6.4.4.} If $(V,+,\cdot)$ is a vector space and $x,y,z\in V,$ then $x,y,$ and $z$ are said to be \textit{collinear} if and only if some line contains all three points.\nl

\textbf{Definition 6.4.5.} Suppose $(V, +, \cdot)$ is a vector space and $x,y,$ and $z$
are non-collinear (i.e., no line contains all three points) elements of V.\nl
a) $\overrightarrow{yx} \cup \overrightarrow{yz}$ is called the \textit{angle} $xyz$ and denoted $\angle xyz.$\nl
b) $\overline{xy}\cup \overline{yz} \cup \overline{zx}$ is called the \textit{triangle} $xyz$ and denoted by
$\triangle xyz$.

\newpage
%page107

\textbf{Theorem 6.4.3.} If $(V,+,\cdot)$ is an inner product space, $x,y,$ and $z$ are non-collinear elements of $V,~ u\in(\overrightarrow{xy} - \{y\}),$ and $v\in(\overrightarrow{yz} - \{y\}),$ then\nl
\[
\dfrac{(x - y)\cdot(z - y)}{\abs{x - y}\abs{z - y}} =  \dfrac{(u - y)\cdot(v - y)}{\abs{u - y}\abs{v - y}}
\]

\textbf{Definition 6.4.6.} Suppose $(V,+,\cdot)$ is an inner product space and $x,y,$ and $z$ are non-collinear elements of $V.$ The \textit{measure} of $\angle xyz$  is denoted by $m\angle xyz$ and defined by
\[
m\angle xyz = \arccos \left(  \dfrac{(x - y)\cdot(z - y)}{\abs{x - y}\abs{z - y}}  \right).
\]

Let $\sin \angle xyz, \cos \angle xyz,$ and $\tan\angle xyz$ denote $\sin(m\angle xyz), \cos(m\angle xyz),$ and $\tan (m\angle xyz),$ respectively. \nl

\textbf{Definition 6.4.7} If $(V,+,\cdot)$ is an inner product space, then an angle $xyz$ is said to be a \textit{right angle} if and only if $(x - y)\cdot(z - y) = 0.$\nl
\textbf{Exercises}\nl
Assuming $(V,R,+,\cdot)$ is an inner product space and $x,y,z\in V,$ prove each of the following. \nl
1. $\abs{x} - \abs{y} \leq \abs{x - y}$\nl

2. $\abs{x + y}^{2} = \abs{x}^{2} + 2x\cdot y + \abs{y}^{2}     $ \nl

3. $\abs{x - y}^{2} = \abs{x}^{2} - 2x\cdot y + \abs{y}^{2}$\nl

4. $\abs{x + y}^{2} - \abs{x - y}^{2} = 4x\cdot y$\nl

5. If $\angle xyz$ is a right angle, then $0 =(x - y)\cdot(z - y) = (y - x)\cdot(z - y) =(x - y)\cdot(y - z) = (y - x)\cdot(y - z).$\nl

6. 	If $\angle xyz$ is a right angle, then\nl
\begin{tabular}{r l}
&a) $\abs{x - y}^{2} = (x - y)\cdot (x - z)$ and \\
&b) $\abs{z - y}^{2} = (z - y)\cdot(z - x)$\\
\end{tabular}\nl

7. $\abs{ax + (1 - a)y} \leq \max \{ \abs{x},\abs{y}\}$ for each $a$ in $[0,1].$\nl

8. If $x,y,$ and $z$ are non-collinear, $u\in (\overrightarrow{yx} -\{y\}),$ and $v\in(\overrightarrow{yz} -\{y\}),$ then  $\angle xyz = \angle uyv.$ \nl

9. 	If $\angle xyz$ is a right angle, $r$ and $s$ are distinct elements of $\overrightarrow{xy}, u$ and $v$ are distinct elements of $\overrightarrow{yz},$ then $(r - s)\cdot (u - v) = 0.$\nl

10. Prove Theorem 6.4.1. (Hint: Show $0 \leq \abs{(x\cdot y)x - (x\cdot x)y}^{2} \ra     $\nl
$\abs{x\cdot y} \leq \abs{x}\abs{y}.$  Then show $\abs{x + y}^{2} \leq \left(\abs{x} + \abs{y}\right)^{2}.$)\nl

11. Prove Theorem 6.4.2.\nl
12. Prove Theorem 6.4.3.\nl

\newpage
%page108
\lhead{Section 6.5}
\textbf{6.5 Triangles} \nl

\textbf{Definition 6.5.1.} If $(V,+,\cdot)$ is an inner product space and $A, B,$ and
$C$ are non-collinear elements of $V,$ let\nl
\begin{tabular}{r r l}
&$a = d(B, C)$&$A = \angle CAB$\\
&$b = d(A, C)$&$B = \angle ABC$\\
&$c = d(A, B)$&$C = \angle BCA$\\
\end{tabular}\nl


\textbf{Theorem 6.5.1.} If $(V,+,\cdot)$ is an inner product space and $\angle ACB$ is a right angle, then\nl
\begin{tabular}{r l}
&a) $a^{2} + b^{2} = c^{2}$\\
&\\
&b) $\sin\angle A = \dfrac{a}{c}$\\
&\\
&c) $\cos\angle A = \dfrac{b}{c}$\\
&\\
&d) $\tan\angle A = \dfrac{a}{b}$\\
&\\
\end{tabular}\nl

\textbf{Theorem 6.5.2.} If $(V,+,\cdot)$ is an inner product space and $A, B,$ and $C$ are non-collinear elements of $V,$ then\nl

\begin{tabular}{r l}
a) &$a^{2} = b^{2} + c^{2} - 2bc\cos\angle A$\\
&$b^{2} = c^{2} + a^{2} -2ca\cos\angle B$\\
&$c^{2} = a^{2} + b^{2} -2ab\cos\angle C$\\
\end{tabular}(Law of Cosines)\\
\begin{tabular}{r l}
&\\
b) &$\dfrac{\sin\angle A}{a} = \dfrac{\sin\angle B}{b} = \dfrac{\sin\angle C}{c}$\\
&\\
\end{tabular}(Law of Sines)\\

\begin{tabular}{r l}
c) &$m\angle PAC + m\angle CAB = \pi$ for each $P$ in $(\overleftrightarrow{AB} - \overrightarrow{AB})$\\
%d) &$m\angle BAC + m\angle ACB \leq \pi$\\
d) &$m\angle A + m\angle B + m\angle C  = \pi$\\
\end{tabular}\nl

\textbf{Example 6.5.1.} If $(V,R,+,\cdot)$ is an inner product space, $A, B,$ and $C$ are non-collinear elements of $V, a = 2, m\angle A = 45\dg,$ and $m\angle B =60\dg,$ find $c.$\nl
Solution. Suppose $(V,R,+,\cdot)$ is an inner product space, $A, B,$ and $C$ are non-collinear elements of $V, a = 2, m\angle A = 45\dg,$ and $m\angle B =60\dg.$
\[m\angle C = 180\dg - m\angle A - m\angle B = 180\dg -45\dg -60\dg =75\dg.\]

By problem 11, Sect. 5.5 on page 91,

\newpage
%page109
\[
\sin\angle C = \sin 75\dg = \dfrac{\sqrt{3} + 1}{2\sqrt{2}}
\]
From the Law of Cosines, we have
\[
c = \dfrac{a\sin\angle C}{\sin\angle A} = \dfrac{2\left(\dfrac{\sqrt{3} + 1}{2\sqrt{2}}\right)}{\dfrac{1}{\sqrt{2}}} = \sqrt{3} + 1.
\]

\textbf{Example 6.5.2.} If $(V,R,+,\cdot)$ is an inner product space, $A, B,$ and $C$ are non-collinear
elements of $V, a=3, b=4,$ and $m\angle C =60\dg,$ find $c.$\nl
Solution. Suppose $(V,R,+,\cdot)$ is an inner product space, $A,B,$ and $C$ are non-collinear elements of $V, a = 3, b = 4,$ and $m\angle C =60\dg.$ By the Law of Cosines, we have \nl

\begin{tabular}{r r l}
&$c^{2}$&$= a^{2} + b^{2} -2ab\cos\angle C$ \\
&&$=3^{2} + 4^{2} -2(3)(4)(1/2) = 13$ \\
&c&$= \sqrt{13}$ \\
\end{tabular}\nl

\textbf{Exercises}\nl

Assuming $(V,R,+,\cdot)$ is an inner product space and $A, B,$ and $C$ are non-collinear elements
of $V,$ prove or solve each of the following. \nl

1. If $a = 2, m\angle C = 30\dg,$ and $ m\angle B =105\dg,$ find $b, c,$ and $m\angle A.$\nl

2. If $b = 3, c = 5,$ and $m\angle A = 120\dg,$ find $a.$\nl

3. If $a =\sqrt{2}, b = 2,$ and $m\angle A = 45\dg,$ find $m\angle B.$ \nl

4. If $a = 24, b = 25,$ and $c = 7,$ find $m\angle B.$\nl

5. If $m\angle A = 30\dg, c = 2,$ and $a = \sqrt{2},$ find all possible values for $m\angle C.$\nl

6. If $\angle ABC$ is a right angle, then $c < b$ and $a < b.$\nl

7. 	If $\angle ABC$ is a right angle, then \nl
\begin{tabular}{r l}
&a) $m\angle BAC < m\angle ABC$ and\\
&b) $m\angle BCA <  m\angle ABC.$\\
\end{tabular}\nl

8. If $P\in(\overleftrightarrow{AB} - \overrightarrow{AB}),$ then $m\angle BCA < m\angle CAP.$ \nl

9. Prove Theorem 6.5.1. \nl

10. Prove Theorem 6.5.2. \nl

\newpage
%page110
\lhead{Section 6.6}
\textbf{6.6 Lines}\nl

Suppose line means line in $(R^{2},+,\cdot)$ as defined in Definition 6.4.3 and $+$ and $\cdot$ are defined as in Theorem 6.3.2.\nl

\textbf{Theorem 6.6.1.} If $L$ is a line and $A,B\in L,$ then $L = \overleftrightarrow{AB}.$\nl

\textbf{Theorem 6.6.2.} The intersection of two distinct lines contains exactly one point.\nl

\textbf{Theorem 6.6.3.} $L$ is a line if and only if there exist numbers $A, B,$ and $C$ such that $L= \{(x,y) \mid x,y\in R \text{~and~} Ax + By + C= 0\}.$\nl

\textbf{Definition 6.6.1.} If $L$ is a line, then $L$ is said to have \textit{slope} if and only if there is a number $m$ (called the slope of $L$) such that for each two points $(u,v)$ and $(a,b)$ of $L,$
\[
m = \dfrac{v - b}{u - a}.
\]

%\textbf{Definition 6.6.2.} If $E(x,y)$ is a statement for each $x$ and $y$ in $R,$ then $E(x,y)$ is said to be %an \textit{equation} for a set $S$ if and only if \nl
%$S = \{ (x,y) \mid x,y\in R \text{~and~} E(x,y)\}.$ \nl

\textbf{Theorem 6.6.4.} If $A,B,C\in R, B \neq 0,$ and $L$ is a line with the
equation $Ax + By + C = 0,$ then $L$ has slope $-\dfrac{A}{B}.$\nl

\textbf{Theorem 6.6.5.} If $L$ is a line containing two points $(x_{1}, y_{1})$ and $(x_{2}, y_{2}),~x_{1} \neq x_{2}$ and $(x, y) \in L,$ then $y =\left( \dfrac{y_{2} - y_{1}}{x_{2} - x_{1}}\right)(x - x_{1}) + y_{1}.$ This is called the two-point form of a line.\nl

\textbf{Example 6.6.1.} Find the slope of the line which contains $(-1,2)$ and $(2,4).$\nl
Solution. Suppose $L$ is a line which contains $(-1,2)$ and $(2,4).$ By Th6.6.5 and Def 6.6.1, $L$ has slope
$\dfrac{4 - 2}{2 - (-1)} = \dfrac{2}{3}.$\nl

\textbf{Theorem 6.6.6.} $L$ is a line with slope $m$ and $(a,b)\in L$ if and only if $y = m(x-a) + b$ is an equation for $L.$ This equation is called the point-slope form of $L.$\nl

\textbf{Example 6.6.2.} If $L$ is a line with slope $- 1/2$ and $(-3,4)\in L,$ find an equation for $L.$\nl

\newpage
%page111

Solution. Suppose $L$ is a line with slope $- 1/2$ and $(-3,4)\in L.$ By Th 6.6.6, $y =-1/2(x-(-3)) + 4$ is an equation for $L.$ This equation simplifies to $2y + x - 5 =0. $\nl

\textbf{Definition 6.6.3.} $\{(x,y) \mid x\in R \text{~and~} y = 0\}$ is called the $x$-axis. $\{ (x,y) \mid y\in R \text{~and~} x = 0\}$ is called the $y$-axis. \nl

\textbf{Definition 6.6.4.} If $L$ is a line which intersects the $y$-axis ($x$-axis),
then each point of the intersection is called a $y-intercept$ ($x-intercept$) of $L.$\nl

\textbf{Theorem 6.6.7.} $L$ is a line with slope $m$ and $y$-intercept $(0, c)$ if and only if $y = mx + c$ is an equation of $L.$ This equation is, called the slope-intercept form of $L. $\nl

\textbf{Example 6.6.3.} Find the slope and $y$-intercept of the line with equation $4x - 3y + 1 =0.$\nl
Solution. Suppose $L$ is a line with equation $4x - 3y + 1 = 0.$
\begin{tabular}{r r l}
&&$4x - 3y + 1 = 0$\\
&$\lra$&$3y = 4x + 1$\\
&$\lra$&$y = \dfrac{4}{3}x + \dfrac{1}{3}$\\
\end{tabular}\nl

By Th 6.6.7, $L$ has slope $\dfrac{4}{3}$ and $y$-intercept $\left(0,\dfrac{1}{3}\right).$\nl

\textbf{Definition 6.6.5.} Two intersecting lines $\overleftrightarrow{AB}$ and $\overleftrightarrow{CD}$ are said to be \textit{perpendicular} if and only if $(A-B)\cdot(C-D) = 0.$\nl

\textbf{Theorem 6.6.8.} Two lines with slopes $m$ and $n$ are perpendicular if and only if $m = -\dfrac{1}{n}.$\nl

\textbf{Definition 6.6.6.} Two lines are said to be \textit{parallel} if and only if they do not intersect.\nl

\textbf{Theorem 6.6.9.} Two lines with slope are parallel if and only if they have the same slope.\nl

\textbf{Exercises}\nl

1. Find an equation of the line which contains the given point and
has the given slope in each of the following problems. \nl

\newpage
%page112
\begin{tabular}{r l}
&a) $(-1,3), m = 2$\\
&\\
&b) $(2, -1), m = -\frac{1}{3} $\\
&\\
&c) $(0, -4), m = -\frac{3}{4}$\\
\end{tabular}\nl

2. 	Find the slope and an equation of the line which contains the given
points in each of the following problems. \nl
\begin{tabular}{r l}
&a) $(1,2), (3,-2)$\\
&b) $(2,3), (-1,5)$ \\
&c) $(0,2), (5,-6)$ \\
\end{tabular}\nl

3. 	In each of the following problems, find the slope and $y$-intercept
of the line with the given equation. \nl
\begin{tabular}{r l}
&a) $5x - 2y + 2 = 0$\\
&b) $2x - 3y - 1 = 0$\\
&c) $3x + 2y - 3 = 0$\\
\end{tabular}\nl


4. If $L$ is a line parallel to the line with equation $2x - 6x - 1 = 0$ and $(2,5)\in L,$ find an equation for $L.$\nl

5. If $L$ is a line parallel to the line containing the two points $(3,1)$ and $(-1,2)$ and $(-1,3)\in L,$ find an equation for $L.$\nl

6. Find the point of intersection of the line with equation $3y - x + 1 = 0$ and the line perpendicular to it which contains $(2,3).$\nl

7. 	If $h,k\in R,$ find the intercepts of the line with equation $\dfrac{x}{h} + \dfrac{y}{k} = 1.$\nl

8. Prove Theorem 6.6.1. \nl

9. Prove Theorem 6.6.2. \nl

10. Prove Theorem 6.6.3.\nl

11. Prove Theorem 6.6.4. \nl

12. Prove Theorem 6.6.5. \nl

13. Prove Theorem 6.6.6. \nl

14. Prove Theorem 6.6.7. \nl

15. Prove Theorem 6.6.8. \nl

16. Prove Theorem 6.6.9. \nl

\newpage
%page113
\lhead{Section 6.7}
\textbf{6.7 Matrices}\nl

A matrix is a rectangular array of numbers.  Matrices have been found to be useful in dealing with a large number of problems people encounter.  In this section we will develop the algebra of matrices and in the next section we will show how they can be used to solve a system of linear equations.\nl

We will use $A_{ij}$ or $a_{ij}$ to denote the element in the $i^{\text{th}}$ row and $j^{\text{th}}$ column.  A matrix with $m$ rows and $n$ columns is called a \textit{m-\text{by}-n matrix}.  The word \textit{dimension}
will refer to the number of rows and columns in a matrix.  A matrix is said to be \textit{square} if it has the same number of rows and columns.  \nl

The first of the arithmetical operations that we will consider is addition. \nl

\textbf{Definition 6.7.1.}  If $A$ and $B$ are matrices with the same dimensions, then there is a sum $A + B$ whose $ij^{\text{th}}$ element is defined by $(A + B)_{ij} = a_{ij} + b_{ij}.$  \nl

\textbf{Example 6.7.1}\nl


\begin{equation*}
\left[
\begin{array}{rrr}
2 & -1 & 3\\
-5&4&6\\
7&0&-3\\
\end{array}
\right] +
\left[
\begin{array}{rrr}
5 & 3 & -2\\
3&-5&1\\
-4&5&4\\
\end{array}
\right]=
\left[
\begin{array}{rrr}
7 & 2 & 1\\
-2&-1&7\\
3&5&1\\
\end{array}
\right]
\end{equation*}\nl

\textbf{Theorem 6.7.1}  If $A,B,$ and $C$ are matrices with the same dimensions, then\nl
\begin{tabular}{l l}
$A + B = B + A$&(Commutative Law)\\
$A + (B+C) = (A + B) + C$&(Associative Law)\\
\end{tabular}\nl

\textbf{Definition 6.7.2.}  There is a zero matrix $\underline{0}$ called the identity for addition whose $ij^{\text{th}}$ element is defined by $\underline{0}_{~ij}$ = 0.\nl
Note : The dimension of $\underline{0}$ will be assumed to be the dimension necessary for the operation to take place.\nl
\textbf{Example 6.7.2}\nl
\begin{equation*}
\left[
\begin{array}{rrr}
4 & -7 & 3\\
-5&2&6\\
7&3&2\\
\end{array}
\right] +
\left[
\begin{array}{rrr}
0 & 0 &0 \\
0&0&0\\
0&0&0\\
\end{array}
\right]=
\left[
\begin{array}{rrr}
4 & -7 & 3\\
-5&2&6\\
7&3&2\\
\end{array}
\right]
\end{equation*}\nl

\textbf{Theorem 6.7.2.}  If $A$ is a matrix, then $\underline{0} + A = \underline{0} = A + \underline{0}.$  (Identity for Addition)\nl

\textbf{Definition 6.7.3.}  If $A$ is a matrix, then there is a matrix $-A$ called the additive inverse of $A$ whose $ij^{\text{th}}$ element is defined by $(A)_{ij} = -a_{ij}.$\nl
\newpage
\textbf{Example 6.7.3.}\nl
\begin{equation*}
\left[
\begin{array}{rrr}
-5 & 2 & -4\\
3&-6&2\\
1&7&-3\\
\end{array}
\right] +
\left[
\begin{array}{rrr}
5 & -2 & 4\\
-3&6&-2\\
-1&-7&3\\
\end{array}
\right]=
\left[
\begin{array}{rrr}
0 & 0 & 0\\
0&0&0\\
0&0&0\\
\end{array}
\right]
\end{equation*}\nl

%\newpage
%Page 114
\textbf{Theorem 6.7.3.}  If $A$ is a matrix, then\nl
$A + -A = \underline{0} = -A + A.$  (Inverse of Addition)\nl

\textbf{Definition 6.7.4.}  If $A$ is a matrix and $k$ is a number, then there is a scalar product $kA$ whose $ij^{\text{th}}$ element is defined by $(kA)_{ij} = ka_{ij}.$\nl


\textbf{Example 6.7.4.}
\begin{equation*}
3\left[
\begin{array}{rrr}
5 & -2 & 4\\
-1&6&2\\
7&0&3\\
\end{array}
\right]=
\left[
\begin{array}{rrr}
15 & -6 & 12\\
-3&18&6\\
21&0&9\\
\end{array}
\right]
\end{equation*}\nl

\textbf{Theorem 6.7.4.}  If $h$ and $k$ are numbers and $A$ and $B$ are matrices with the same dimensions, then \nl
$k(A\pm B) = kA \pm kB$\nl
$(hk)A = h(kA)$\nl
$(h + k)A = hA + kA$\nl
$1A = A$\nl
$-1A = -A$\nl

\textbf{Definition 6.7.5.} If $A$ is an $m$-by-$p$ matrix and $B$ is a $p$-by-$n$ matrix, then there is a $m$-by-$n$ product matrix $AB$ whose $ij^{\text{th}}$ element is defined by
\[
(AB)_{ij} = (\sum_{k=1}^{p}a_{ik}b_{kj}).
\]\nl

\textbf{Example 6.7.5.}\nl
\begin{equation*}
\left[
\begin{array}{rrr}
2 & -1 & 4\\
-3&4&0\\
%7&0&3\\
\end{array}
\right]
\left[
\begin{array}{rr}
2 & 2\\
1&0\\
-3&-1\\
\end{array}
\right]=
\left[
\begin{array}{rr}
-9&0\\
-2&-6\\
\end{array}
\right]
\end{equation*}\nl

\textbf{Theorem 6.7.5.}  If $A,B,$ and $C$ are $n$-by-$n$ matrices and $k$ is a number, then \nl
\begin{tabular}{ll}
$A(BC) = (AB)C$&(Associative Law of Multiplication)\\
$A(B\pm C) = AB \pm AC$&(Distributive Law)\\
$(B \pm C)A = BA \pm CA$&\\
$(kA)B = k(AB) = A(kB)$&\\
$\underline{0}A = \underline{0} = A\underline{0}$\\
\end{tabular}\nl

Note : The Commutative Law does not hold for matrix multiplication.\nl

\textbf{Definition 6.7.6.} There is a matrix $I$ called the identity matrix for multiplication whose $ij^{\text{th}}$ element is defined by $I_{ij} = 1$ if $i = j$ and $0$ if $i \neq j.$\nl

Note : The dimensions of $I$ will be assumed to be the dimensions necessary for the operation to take place.\nl

\textbf{Example 6.7.6.}\nl
\begin{equation*}
\left[
\begin{array}{rrr}
6 & 1 & 0\\
-2&7&-3\\
3&-5&4\\
\end{array}
\right]
\left[
\begin{array}{rrr}
1 & 0 & 0\\
0&1&0\\
0&0&1\\
\end{array}
\right]=
\left[
\begin{array}{rrr}
6 & 1 & 0\\
-2&7&-3\\
3&-5&4\\
\end{array}
\right]
\end{equation*}\nl
%\newpage
%page 116
\textbf{Theorem 6.7.6.} If $A$ is a square matrix, then $IA = A = AI.$  (Identity for Multiplication)\nl

\textbf{Definition 6.7.7.}  If $A$ and $B$ are square matrices such that $AB = I,$ then $B$ is called the multiplicative inverse of $A$ and denoted by $A^{-1}.$\nl
There are row operations which will transform a matrix into a form where the unknowns can easily be determined.\nl
The allowable row operations for matrices are :\nl
3) $E_{ij}(a)$ -- add to row $i$ the result of multiplying $a$ times row $j$ of the matrix.\nl
2) $Ej(a)$ -- multiply row $j$ by $a$ to form a new row $j$ of the matrix.\nl
1) $P_{ij}$ -- interchange rows $i$ and $j$ of the matrix.\nl

\textbf{Theorem 6.7.7(a)}  If $A$ is a $m$-by-$m$ matrix, $I$ is the identity matrix for multiplication and there exists a sequence of allowable row operations that will transform the $m$-by-$2m$ matrix $[A | I]$ into a $m$-by-$2m$ matrix of the form $[I | C],$ then $A^{-1} = C$ is the multiplicative inverse for $A$.\nl

\textbf{Example 6.7.7.}  We start with a matrix of the form $[A | I]$ where $A$ occupies the first three columns and $I$ occupies the next three columns.  Between the two 3-by-6 matrices is a list of the row operations used to transform the first 3-by-6 matrix into the second 3-by-6 matrix.\nl

%\begin{tabular}{|rrrcrrr|c|rrrcrrr|}
%2&0&4&$|$&1&0&0&$\ra$&1&0&2&$|$&1/2&0&0\\
%$-1$&3&$1$&$|$&0&1&0&$E(.5)_{1}$&0&3&3&$|$&1/2&1&0\\
%0&2&$-1$&$|$&0&0&1&$E(1)_{21}$&0&2&$-1$&$|$&0&0&1\\
%\end{tabular}\nl

\begin{equation*}
\left[\begin{array}{rrrcrrr}
2&0&4&|&1&0&0\\
-1&3&1&|&0&3&3\\
0&2&-1&|&0&0&1\\
\end{array}
\right]
\begin{array}{c}
\ra\\
E(.5)_{1}\\
E(1)_{21}\\
\end{array}
\left[
\begin{array}{rrrcrrr}
1&0&2&|&1/2&0&0\\
0&3&3&|&1/2&1&0\\
0&2&-1&|&0&0&1
\end{array}
\right]
\end{equation*}\nl










We continue using row operations until we obtain a 3-by-6 matrix of the form $[I | C]$ where $C$ is the last three columns on the right.\nl
\end{large}
%\begin{tabular}{c|rrrcrrr|c|rrrcrrr|}
%$\ra$&E(-2)_{32}&$\ra 0&1&1&|&1/6&1/3&0\\$&1&0&0&$|$&5/18&$-4/9$&2/3\\
%$E(1/3)_{2}$&0&1&1&$|$&1/6&1/3&0&$E(-2)_{13}$&0&1&0&$|$&1/18&1/9&1/3\\
%$E(-2)_{32}$&0&0&$-3$&$|$&$-1/3$&$-2/3$&1&$E(-1)_{23}$&0&0&1&$|$&1/9&2/9&$-1/3$\\
%\end{tabular}\nl
\begin{equation*}
\begin{array}{c}
\ra\\
E(1/3)_{2}\\
E(-2)_{32}\\
\end{array}
\left[
\begin{array}{rrrcrrr}
1&0&2&|&1/2&0&0\\
0&1&1&|&1/6&1/3&0\\
0&0&-3&|&-1/3&-2/3&1\\
\end{array}
\right]
\begin{array}{c}
\ra \newline E(-1/3)_{3}\\
E(-2)_{13}\\
E(-1)_{23}\\
\end{array}
\left[
\begin{array}{rrrcrrr}
1&0&0&|&5/18&-4/9&2/3\\
0&1&0&|&1/18&1/9&1/3\\
0&0&1&|&1/9&2/9&-1/3\\
\end{array}
\right]
\end{equation*}\nl










\begin{large}

By Theorem 6.7.7(a), we have \nl
 \begin{equation*}
A^{-1} = C = \left[
\begin{array}{rrr}
5/18 & -4/9 & 2/3\\
1/18&1/9&1/3\\
1/9&2/9&-1/3\\
\end{array}
\right]
\end{equation*}\nl

\textbf{Theorem 6.7.7(b)}  If $A$ and $B$ are $m$-by-$n$ matrices, $k$ is a number and $n$ a positive integer, then\nl
$(AB)^{-1} = B^{-1}A^{-1}$ (Inverse of Multiplication)\nl
$AA^{-1} = A^{-1}A = I$\nl
$(kA)^{-1} = k^{-1}A^{-1} = (1/k)A^{-1}$\nl
$(A^{-1})^{-1} = A$\nl


\textbf{Definition 6.7.8.}  If $A$ is a matrix, then there is a transpose matrix $A^{t}$ whose $ij^{\textbf{th}}$ element is defined by $A_{ij} = a_{ji}.$\nl

\textbf{Example 6.7.8.}\nl
\begin{equation*}
\left[
\begin{array}{rr}
5 &1 \\
-3&6\\
2&-4\\
\end{array}
\right]^{t}
=
\left[
\begin{array}{rrr}
5 &-3&2 \\
1&6&-4\\
\end{array}
\right]
\end{equation*}\nl
%\newpage
%page 117

\textbf{Theorem 6.7.8.}  If $A$ and $B$ are matrices with the same dimension, then\nl
$(A + C)^{t} = A^{t} + B^{t}$    (Matrix Transposition)\nl
$(AB)^{t} = B^{t}A^{t}$\nl
$(kA)^{t} = kA^{t}$\nl
$(A^{t})^{t}= A$\nl
$(A^{t})^{-1} = (A^{-1})^{t}$\nl

\textbf{Definition 6.7.9.}  If $A$ is a square matrix and $n$ is a positive integer, then $A$ to the $n^{\text{th}}$ power is denoted by $A^{n}$ and defined by $A^{n} = AAAA \cdots A$ ($n$ factors.)\nl

\textbf{Theorem 6.7.9.}  If each of $A$ and $B$ is a $m$-by-$m$ matrix with an inverse and $n$ is a positive integer, then \nl
$(a^{n})^{-1} = (A^{-1})^{n}$\nl
$(A^{n})^{t} = (A^{t})^{n}$\nl

\newpage
\lhead{Section 6.8}

%\end{large}
\textbf{Section 6.8 Systems of Linear Equations}\nl

For a system of linear equations, there are three possibilities:\nl
1) If there are more unknowns than equations, then the system is called under-defined and there are either infinitely many solutions or no solution.\nl
2) If you have the same number of unknowns as equations, then there is either a unique solution or no solution.  If there is a solution, it can be found by either Theorem 6.8.1 or Theorem 6.8.2.\nl
3) If you have more equations than unknowns, then the system is called over-defined and there is no unique solution, but there is a best approximation given by Theorem 6.8.3.\nl

We will discuss two methods of solving systems of linear equations with the same number of unknowns as equations.  The \textit{first method} of solving a system of linear equations is to apply the following theorem:\nl

\textbf{Theorem 6.8.1.}  If $A$ is an $m$-by-$m$ matrix, $Y$ is an $m$-by-$1$ matrix, $I$ is the identity matrix for multiplication and there is a sequence of allowable row operations that will transform the $m$-by-$(m + 1)$ matrix $[A | Y]$ into a $m$-by-$(m + 1)$ matrix of the form $[I | C],$ then $X = C$ is the solution of the system of linear equations $AX = Y.$\nl

\textbf{Example 6.8.1.}  We will consider the following system of linear equations:\nl
\begin{tabular}{r r r c r}
$2x_{1}$&&$+~4x_{3}$&=& 2\\
$-x_{1}$&$+3x_{2}$&$+x_{3}$&=& $-1$\\
&$+2x_{2}$&$-x_{3}$&=& 3\\
\end{tabular}\nl

This system of linear equations can be abbreviated using matrices as $AX = Y$ where \nl
\begin{equation*}
\textit{A =}\left[
\begin{array}{rrr}
2 &0&4 \\
-1&3&1\\
0&2&-1\\
\end{array}
\right]
\text{~and~} \textit{~Y~=~}\left[
\begin{array}{rrr}
2 \\
-1\\
3\\
\end{array}
\right]
\end{equation*}\nl

We start with a matrix in the form $[A | Y]$ and then, using allowable row operations, transform it into a matrix of the form $[I | C].$\nl

%\begin{tabular}{|rrrcr|c|rrrcr|}
%2&0&4&$|$&2&$\ra$&1&0&2&$|$&1\\
%$-1$&3&1&$|$&$-1$&$E(.5)_{1}$&0&3&3&$|$&0\\
%0&2&$-1$&$|$&3&$E(1)_{21}$&0&2&$-1$&$|$&3\\
%\end{tabular}\nl

\begin{equation*}
\left[
\begin{array}{rrrcr}
2&0&4&|&2\\
-1&3&1&|&-1\\
0&2&-1&|&3\\
\end{array}
\right]
\begin{array}{c}
\ra\\
E(.5)_{1}\\
E(1)_{21}\\
\end{array}
\left[
\begin{array}{rrrcr}
1&0&2&|&1\\
0&3&3&|&0\\
0&2&-1&|&3\\
\end{array}
\right]
\end{equation*}\nl
\begin{center}
and\nl
\end{center}

%\begin{tabular}{c|rrrcr|c|rrrcr|}
%$\ra$&1&0&2&$|$&1&$\ra E(-1/3)_{3}$&1&0&0&$|$&3\\
%$E(1/3)_{2}$&0&1&1&$|$&0&$E(-2)_{13}$&0&1&0&$|$&1\\
%$E(-2)_{32}$&0&0&$-3$&$|$&3&$E(-1)_{23}$&0&0&1&$|$&$-1$\\
%\end{tabular}\nl

\begin{equation*}
\begin{array}{c}
\ra\\
E(1/3)_{2}\\
E(-2)_{32}\\
\end{array}
\left[
\begin{array}{rrrcr}
1&0&2&|&1\\
0&1&1&|&0\\
0&0&-3&|&3\\
\end{array}
\right]
\begin{array}{c}
\ra E(-1/3)_{3}\\
E(-2)_{13}\\
E(-1)_{23}\\
\end{array}
\left[
\begin{array}{rrrcr}
1&0&0&|&3\\
0&1&0&|&1\\
0&0&1&|&-1\\
\end{array}
\right]
\end{equation*}\nl

%\newpage
We continue applying row operations until we obtain a matrix of the form $[I | C].$\nl
By Theorem 6.8.1, the solution of the system of linear equations $AX =Y$ is \nl

\begin{equation*}
\textit{X = C =~}\left[
\begin{array}{r}
3\\
1\\
-1\\
\end{array}
\right]
\end{equation*}\nl

Note : The $A$ used in this example is the same $A$ used in Example 6.7.7 and the sequence of row operations
is also the same.\nl

The \textit{second method} of solving this system of linear equations is to apply the following theorem:\nl

\textbf{Theorem 6.8.2.}  If $A$ is an $m$-by-$m$ matrix which has an inverse, $I$ is the identity matrix for multiplication and $Y$ is an $m$-by-1 matrix, then $X = A^{-1}Y$ is a solution to the system of linear equations $AX = Y.$\nl
Proof.\nl
Multiplying both sides of the equation $AX = Y$ by $A^{-1},$ we get $A^{-1}AX = A^{-1}Y.$\nl
Since $A^{-1}A = I,$ the solution is $X = A^{-1}Y.$\nl

\textbf{Example 6.8.2.}  Suppose $A$ and $Y$ are defined as in Example 6.8.1, then the solution of the system of linear equations $AX = Y$ is $A^{-1}Y =$\nl

\begin{equation*}
\textit{X =}\left[
\begin{array}{rrr}
5/18 & -4/9 & 2/3\\
1/18&1/9&1/3\\
1/9&2/9&-1/3\\
\end{array}
\right]
\left[
\begin{array}{r}
2\\
-1\\
3\\
\end{array}
\right]=
\left[
\begin{array}{r}
3\\
1\\
-1\\
\end{array}
\right]
\end{equation*}\nl

\textbf{Definition 6.8.1.}  If $B$ is an $m$-by-1 matrix, then the norm of $B$ denoted by $|B|$ is defined to be $\sqrt{B^{t}B}.$\nl

Note : $B^{t}B$ is a number greater than or equal to 0.\nl

\textbf{Definition 6.8.2.}  If $A$ is an $m$-by-$n$ matrix, $Y$ is an $m$-by-1 matrix, then the matrix $X$ is said to be the least squares solution to $AX = Y$ if and only if for each $m$-by-1 matrix $Z$ not equal to $X, ~|AZ - Y | > |AX - Y|.$\nl 

\textbf{Theorem 6.8.3.} If $A$ is an $m$-by-$n$ matrix, $Y$ is an $m$-by-1 matrix, $m > n$ and there exists a least squares solution for $AX = Y,$ then the least squares solution of $AX = Y$ is   $A^{t}AX = A^{t}Y.$\nl

Note : Since $A^{t}A$ is an $n$-by-$n$ matrix and $A^{t}Y$ is an $n$-by-1 matrix, then the equation $A^{t}AX = A^{t}Y$ can be solved using Theorem 6.8.1 or Theorem 6.8.2.\nl










%\end{document}





















%\end{document}




%\include{Roach}


%\documentclass[12pt]{amsart}
%\documentclass{report}
%\usepackage{amssymb,latexsym}
%\usepackage{amsmath}
%\usepackage{amsfonts}
%\usepackage{graphics,graphicx}
\setlength{\topmargin}{-10pt}
\setlength{\voffset}{-40pt}
\setlength{\textheight}{700pt}
%\usepackage{amsscr}
%\usepackage[amsscr]{eucal}

%\newtheorem{theorem}{Theorem}
%\newtheorem{lemma}{Lemma}
%\newtheorem{proposition}{Proposition}
%\newtheorem{definition}{Definition}
%\newtheorem{corollary}{Corollary}
%\newtheorem{notation}{Notation}
%\newtheorem{remark}{Remark}
%\newtheorem{assertion}{Assertion}
%\newtheorem{claim}{Claim}

%\newcommand{\abs}[1]{\left|#1\right|}
%\newcommand{\norm}[1]{\left\|#1\right\|}
%\newcommand{\p}[1]{\left(#1\right)}
%\newcommand{\tail}[1]{#1^{\p{N}}}
%\newcommand{\trunk}[1]{#1^{\left[N\right]}}
%\newcommand{\dg}{^{\circ}}
%\newcommand{\nl}{\newline}
%\newcommand{\fl}{\flushleft}
%\newcommand{\tild}{\texttt{\symbol{126}}}
%\newcommand{\tild}{\sim}
%\newcommand{\up}{\wedge}
%\newcommand{\dn}{\lor}
%\newcommand{\lra}{\leftrightarrow}
%\newcommand{\ra}{\to}
%\newcommand{\qqu}{\qquad \qquad}
%\newcommand{\onto}{\dashv}
%\newcommand{\power}{\Pi}
%\newcommand{\univ}{\textbf{U}}
%\newcommand{\ns}{\emptyset}
%\newcommand{\of}{$(F, +, \cdot , <)\,$}
%\newcommand{\zp}{Z^{+}}
%\newcommand{\fld}{ $(F,+,\cdot)$}
%\newcommand{\po}[1]{\dfrac{\pi}{#1}}
%\newcommand{\lpo}[1]{\frac{\pi}{#1}}
%\newcommand{\q}{\indent}


%\pagestyle{plain}


%\begin{document}


\setcounter{page}{120}
\lhead{ }
%\begin{large}

\begin{center}
\textbf{Appendix 1}
\end{center}
$ $\nl

The purpose of this section is to prove Theorem 4.2.1 using only theorems from calculus which are intuitively obvious.
% Suppose $a,b, c \in P$ and $a < b.$  If $f$ is a positive valued function defined on $[a,b],$ then the integral of $f$ on $[a, b]$ is denoted by $\int_{a}^{b} f$ and is equal to the area bounded by $f$ and the lines $x = a,\, x = b,$ and $y = 0.$\nl

%\textbf{Definition 1.}  If $0 < a < b$ and $f$ is a positive valued function integrable on $[a, b],$ then let
%$\int_{b}^{a}$ denote $-\int_{a}^{b}$ and let $\int_{a}^{a} f = 0.$\nl

\textbf{Definition 1.}  Let $h$ denote the function $\{(x, y) \mid x \in P \text{~and~}y = \dfrac{1}{x}\}.$\nl

\textbf{Definition 2.}  Suppose $a, b, c \in P$ and $a < b.$  The integral of $h$ from $a$ to $b$ is denoted by $\int_{a}^{b} h$ and is equal to the area bounded by $h$ and the lines $x = a,\,\, x = b,$ and $y = 0.$\nl
\indent Let $L$ denote the function $\{(x, y)\mid x \in P \text{~and~} y = -\int_{x}^{1} h \text{~if~} x < 1,~y = 0 \text{~if~} x = 1, \text{~and~} y = \int_{1}^{x} h \text{~if~} x > 0\}.$\nl

Figure 1 on page 115 is a graph of the function $L.$\nl

\textbf{Definition 3.}  If $0 < a < b,$ then let $\int_{b}^{a} h$ denote $-\int_{a}^{b} h$ and let $\int_{a}^{a} h = 0.$\nl

%Suppose $a < b.$ We will assume that (1) there is a collection $A_{[a,b]}$ of nonnegative valued functions defined on $[a,b]$ and (2) there is a meaning for the word area such that
%if $f\in A_{[a,b]},$ then $\{(x,y)\mid x\in [a,b] \text{~and~} 0 \leq  f(x)\}$ has area and its value is denoted by $\int_{a}^{b}f.$
%If $f\in A_{[a,b]},$ then $f$ is said to be integrable on $[a,b].$ If $f\in A_{[a,b]}$ and $c \in [a,b],$ let $\int_{b}^{a} = -\int_{a}^{b}f$ and $\int_{c}^{c} f$  denote $0.$\nl



\textbf{Theorem 1.} If $a,b,c \in P,~a < c < b,$ and $f$ is a function integrable on $[a, b],$ then
$\int_{a}^{b} f = \int_{a}^{c} f + \int_{c}^{b} f.$   That is, the area of the whole is the sum of the areas of the parts.  It follows from Definition 3 that this theorem is true even if $c$ is not between $a$ and $b.$\nl

\textbf{Theorem 2.} If $0 < x < y,$ then $1 - \dfrac{x}{y} \leq L(y) - L(x) \leq \dfrac{y}{x} - 1.$\nl
Proof.  Suppose $0 < x < y.$  Since $h$ is decreasing, $\dfrac{1}{y} \leq h(t) \leq \dfrac{1}{x}$ for each
$t\in [x, y].$  $\int_{x}^{y} h$ is less than the area of a rectangle which contains the region whose area $\int_{x}^{y} h$ represents;\nl
and $\int_{x}^{y} h$ is greater than the area of a rectangle which is enclosed by the region whose area
$\int_{x}^{y} h$ represents.   Thus,\nl
\begin{tabular}{l l}
&\\
$\dfrac{1}{y}(y - x) < \int_{x}^{y} h < \dfrac{1}{x}(y - x),$&\\
&\\
$\dfrac{1}{y}(y - x) < \int_{x}^{1} h + \int_{1}^{y} h < \dfrac{1}{x}(y - x),$&(Th. 1)\\
&\\
$1 - \dfrac{x}{y} < \int_{1}^{y} h - \int_{1}^{x} h < \dfrac{y}{x} - 1,$ and&(Def. 3)\\
&\\
$1 - \dfrac{x}{y} < L(y) - L(x) < \dfrac{y}{x} -1.$&(Def. 2)\\
\end{tabular}\nl

Figure 2 on page 115 is a graph illustrating the rectangles discussed in the proof.\nl

\textbf{Theorem 3.} If $b > 0,$ then $L(b) \leq b - 1.$ \nl

Proof.\nl
\begin{tabular}{l l}
Case 1: $b > 1.$  If $a = 1,$ then $0 < a < b$ and $L(a) = 0.$&By Th. 2, we get $L(b) < b - 1.$\\
Case 2: $b = 1.~L(b) = 0 = 1 - 1 = b - 1.$&By Def. 2.\\
Case 3: $b < 1.$ If $a = 1,$ then $0<b<a$ and $L(a)=0.$&By Th. 2, we get $1 - b < -L(b)$\\
&and $L(b) < b - 1$\\
\end{tabular}\nl

%\begin{figure}[h!]
%	\centerline{\includegraphics[scale=0.5]{Desert.eps}}
%	\caption{Desert Phoo.}
%	\caption{Top-Left}
%	\label{Fi: Appendix1}
%\end{figure}

%\newpage
%\lhead{}

\textbf{Theorem 4.} If $a, b, c \in P$ and $a < b,$ then $\int_{a}^{b}h = \int_{ca}^{cb} h.$\nl
The proof of this theorem is beyond the scope of this book.  Instead, I will show that $\int_{a}^{b} h$ is
approximately equal to $\int_{ca}^{cb} h.$  By approximating $\int_{a}^{b} h$ and $\int_{ca}^{cb} h$ with trapezoids, we get\nl
Case 1: $a < b.$ $\int_{a}^{b} h \cong \dfrac{\frac{1}{a} + \frac{1}{b}}{2}(b - a) = \dfrac{a + b}{2ab}(b-a)$
and\nl
$\int_{ca}^{cb} h \cong \dfrac{\frac{1}{ca} + \frac{1}{cb}}{2}(cb - ca) = \dfrac{a + b}{2ab}(b - a).$  Thus,
$\int_{a}^{b} h \cong \int_{ca}^{cb} h.$\nl
Case 2: $b < a.$  $\int_{a}^{b} h = - \int_{b}^{a} h \cong -\int_{cb}^{ca} h = \int_{ca}^{cb} h.$\hfill (By Case 1)\nl
Case 3: $a = b.$  $\int_{a}^{b} h = 0 = \int_{ca}^{cb} h.$\hfill(Def. 3)\nl

Figure 3 on page 115 is a graph illustrating the trapezoids discussed in the proof.\nl
%\newpage
%\lhead{}


%\textbf{Definition 1.} Let $h$ denote the function $\{(x,y) \mid x\in R \text{~and~} y = \frac{1}{x}\}.$ Let $L$ denote the function $\{(x,y) \mid x\in R \text{~and~} y = \int_{1}^{x} h\}.$\nl

%\textbf{Should $x, y, 1/x, 1/y,$ and $y - x$ all be underlined inside the integrals, as $x$ and $y$ are constants?}\nl

%\textbf{Theorem 5.} If $0 < x < y,$ then $1 - \frac{x}{y} \leq L(y) - L(x) \leq \frac{y}{x} - 1.$\nl

% Proof. Suppose $0 < x < y.$ Since $\frac{1}{y} \leq h(t ) \leq \frac{1}{x}$ for each $t\in [x,y],$ we have,
%$\int_{x}^{y}\frac{1}{y} \leq \int_{x}^{y} h \leq \int_{x}^{y}\frac{1}{x}.$  (Th. 3.4) Thus \nl $\frac{1}{y}(y %- x) \leq \int_{x}^{1} h + \int_{1}^{y} h \leq \frac{1}{x}(y - x),$ (Th. 2), and \nl $ \frac{1}{y}(y - x) \leq %\int_{1}^{y} h  - \int _{1}^{x} h \leq \frac{1}{x}(y - x),$ (Def.). So by Definition,\nl


%$1 - \dfrac{x}{y} \leq L(y) - L(x) \leq \dfrac{y}{x} -1.$\nl

%\textbf{Theorem 6.} If $x > 0,$ then $L(x) \leq (x-1).$ \nl
%\newpage
%\lhead{}

\textbf{Theorem 5.} If $x,y\in P,$ then $L(xy) = L(x) + L(y).$\nl
Proof.  Suppose $x,y \in P.$  $L(xy) = \int_{1}^{xy} h = \int_{1}^{x} h + \int_{x}^{xy} h = \int_{1}^{x} h + \int_{1}^{y} h$\nl
$= L(x) + L(y).$\nl

\newpage
\lhead{}

\begin{center}
These are the graphs for Appendix 1.
\end{center}

\begin{figure}[h!]
	\centerline{\includegraphics[scale=1.2]{Appendix1.eps}}
%	\caption{Desert Phoo.}
%	\caption{Top-Left}
	\label{Fi: Appendix1}
\end{figure}

%\end{document}

\newpage
\lhead{}
%Appendix2
\begin{center}
\textbf{Appendix 2}
\end{center}
The purpose of this Appendix is to provide a method of calculating the value of $L$ for any positive number to six significant figures using only properties which follow from Theorem 4.2.1 and to calculate $e$ to six significant figures.  To illustrate the procedure, we will calculate the value of $L$ of all positive integers less that or equal to 10 to six significant figures.\nl

From Theorem 4.21 and Theorem 4.2.2 we have $1 - 1/x < L(x) < x - 1.$\nl
If we use $x - 1$ as an approximation of $L(x),$ our error is less than $$x - 1 - (1 - 1/x) = x + 1/x - 2.$$ %\nl
If $x = 1.000001,$ then the error is equal to $x + 1/x - 2 = 0$ to six significant figures.\nl
We will first find the inverse of $L.$  That is, we will find an expression for $x$ as a function of $L(x)$ using Theorem 4.3.2.  Given a value of $L(x),$ we can find $x$ as follows:%\nl
\begin{eqnarray*}
L(x) &=& L(x) \, 1000000\cdot 000001 \\
&=& L(x) \, 1000000\cdot L(1.000001)\\
&=& L \left(1.000001^{L(x)1000000} \right).  
\end{eqnarray*}
Therefore, $x = 1.000001^{L(x)\cdot 1000000}.$

Since $e = E(1), L(e) = 1, e = 1.000001^{1000000} = 2.71828.$\nl
If $x < 1,$ use the formula $L(x) = - L(x)$ to convert the problem to one of finding the value of $L$ for a number $> 1.$\nl

The procedure for finding the value of $L(x)$ for some number $x > 1$ is as follows:\nl
1)  Find two numbers $u$ and $v$ which have prime factors with known values of $L$ such that $u < x < v.$  For example, to compute $L(2)$ we will use $x = 2, u = 1, L(1) = 0, v = 2.718282,$ and $L(2.718282) = 1.$\nl
2)  Then use the formula $f(x) = \left(\frac{L(v)- L(u)}{v - u}\right)(x - u) + L(u)$ to find an approximation for $L(x).$  This is the formula for the straight line between the points $(u, L(u))$ and $(v, L(v)).$  In this case $f(x) = .581977.$\nl
3)  Set $L(u) = f(x).$  Therefore, $u = 1.000001^{L(u)1000000}.$  In this case,  $L(u) = .581977$ and $u = 1.000001^{.581977*1000000} = 1.789572.$\nl
4)  Repeat Steps 2 and 3 until there is no change in $L(u)$ for six significant figures.  In this case, we get $L(2) = .6931474,$ which is accurate to six figures in seven or more iterations.  This is the end of the    procedure.\nl

$L(4) = L(2^{2}) = 2L(2) = 2\cdot.6931474 = 1.3862948.$\nl
To calculate $L(x)$ for $x = 3,$ we use $u = 2, L(2) =.6931474, v = 4,$ and $L(4) = 1.3862948.$  After 8 iterations, we get $L(3) = 1.098613.$\nl
$L(8) = L(2^{3}) = 3\cdot L(2) = 3\cdot.6931474 = 2.079442.$\nl
To calculate $L(x)$ for $x = 5,$ we use $u = 4, L(4) = 1.3862948, v = 8,$ and $L(8) = 2.079442.$  After 8 iterations, we get $L(5) = 1.609438.$\nl
$L(6) = L(2\cdot3) = L(2) + L(3) = .6931474 + 1.098613 = 1.791760.$\nl
To calculate $L(x)$ for $x = 7,$ we use $u = 4, L(4) = 1.3862948, v = 8,$ and $ L(8) = 2.079442.$  After 5 iterations, we get $L(7) = 1.945911.$\nl
$L(9) = L(3^{2}) = 2L(3) = 2\cdot1.098613 = 2.197226.$\nl
$L(10) = L(2\cdot5) = L(2) + L(5) = .6931474 + 1.609438 = 2.302585.$\nl

To find the log of a number base 10, we use the formula $Log_{10}(x) = \frac{L(x)}{L(10)}.$   For example,
$Log_{10}(2) = L(2)/L(10) = .6931474/2.302585 = .301030.$\nl
%\newpage

To illustrate the general procedure, we will use the prime number $x = 251.$  Let $u = 250 = 2\cdot5^{3}$ and $v = 252 = 7\cdot2^{2}\cdot3^{2}.$\nl

$L(250) = L(2\cdot5^{3}) = L(2) + 3L(5) = .6931474 + 3\cdot1.609438 = 5.521461$ and\nl
$L(252) = L(7\cdot2^{2}\cdot3^{2}) = L(7) + 2L(2) + 2L(3) = 1.945911 + 2\cdot.6931474 + 2\cdot1.098613 = 5.529432.$  \nl
After 2 iterations, we get $L(251) = 5.525455.$\nl

\newpage
\lhead{}
\begin{center}
\textbf{Appendix 3}
\end{center}
$ $\nl
Part a)\nl

The purpose of this Appendix is to calculate $\pi$ to 6 sugnificant figures and to calculate the value of any trigonometric function of any angle accurate to 6 significant figures using only properties that follow from Theorem 5.1.1.  By Theorem 5.5.1 and successive applications of identities ID37 through ID39 of Theorem 5.4.3, we have:\nl
$\cos \left(\po{6}\right) = \sqrt{3}/2 = .8660254$\nl
%$\sin\left(\po{24}\right) = \sqrt{\left(1 + \cos\left(\po{12}\right)\right)/2} = \sqrt{(1 - .8660254)/2} = %.2588191$\nl
$\cos\left(\po{12}\right) = \sqrt{\left(1 + \cos\left(\po{6}\right)\right)/2} = \sqrt{(1 + .8660254)/2} = .9659258$\nl
$\cos\left(\po{24}\right) = \sqrt{\left(1 + \cos\left(\po{12}\right)\right)/2} = \sqrt{(1 + .9659258)/2} = .99144486$\nl
$\cos\left(\po{48}\right) = \sqrt{\left(1 + \cos\left(\po{24}\right)\right)/2} = \sqrt{(1 + .9914449)/2} = .99785892$\nl
$\cos\left(\po{96}\right) = \sqrt{\left(1 + \cos\left(\po{48}\right)\right)/2} = \sqrt{(1 + .9978589)/2} = .999464587$\nl
$\cos\left(\po{192}\right) = \sqrt{\left(1 + \cos\left(\po{96}\right)\right)/2} = \sqrt{(1 + .9994646)/2} = .999866138$\nl
%$\sin\left(\po{384}\right) = \sqrt{\left(1 - \cos\left(\po{192}\right)\right)/2} = \sqrt{(1 - .9998661)/2} = %.008181140$\nl
$\cos\left(\po{384}\right) = \sqrt{\left(1 + \cos\left(\po{192}\right)\right)/2} = \sqrt{(1 + .9998661)/2} = .9999665339$\nl
$\cos\left(\po{768}\right) = \sqrt{\left(1 + \cos\left(\po{384}\right)\right)/2} = \sqrt{(1 + .9999665339)/2} = .999999163344$\nl
$\sin\left(\po{768}\right) = \sqrt{\left(1 - \cos\left(\po{384}\right)\right)/2} = \sqrt{(1 + .9999665)/2} = .004090604$\nl
%$\tan\left(\po{768}\right) = \dfrac{1 - \cos(\pi/384)}{\sin(\pi/384)} = \dfrac{1 - .9999665}{.008181140} = %.004090638$\nl

By Theorem 5.1.1, $\sin(x) < x < \sin(x)/\cos(x).$  Therefore, by using $x$ as an approximation for $\sin(x),$ our error is less that $\sin(x)/\cos(x) - \sin(x).$   For $x = \pi/768,$ our error is $\sin(x)/\cos(x) - \sin(x) = 0$ for seven decimal places.\nl
$\pi = 768\cdot(\pi/768) = 768\cdot\sin(\pi/768) = 768\cdot .004090604 = 3.14159.$\nl

%\textit{For small values if $\theta, \sin \theta$ and $\tan \theta$ are both good approximations of $\theta.$  %$\sin \theta$ is a little less than $\theta$ and $\tan \theta$ is a little more than $\theta.$   We can get an %even better approximation of $\theta$ by taking the average of the two.  We can, therefore, obtain a value of %$\pi$ accurate to 6 significant figures.}\nl

%\textit{$\pi = 768\cdot(\pi/768) \cong 768\cdot(\sin(\pi/768) + \tan(\pi/768))/2$\nl
% $\qquad= 768\cdot(.004090604 + .004090638)/2 = 3.14159.$}\nl

Since $\pi$ is an irrational number, an exact value of $\pi$ cannot be calculated, but by Theorem 5.5.1 and successive applications of identities ID37 through ID39 of Theorem 5.4.3, $\pi$ can be computed to any desired degree of accuracy, limited only by the accuracy of the device used.\nl

Part b)\nl

The procedure for calculating $\sin$ and $\cos$ of any angle to six significant figures is as follows:  First, we must reducae the angle to an angle in the first quadrant using the process described on Page 89.\nl

For example, we will use $\angle = 25\dg = 25\dg\cdot\pi/180 = .4363323$ radians.\nl
$\angle = (\angle/2^{7})\cdot2^{7} = (.4363323/2^{7})\cdot2^{7}.$  Let $y = .4363323/2^{7} = .003408846.$\nl
By Theorem 5.1.1, $\sin(y) < y < \sin(y)/\cos(y).$   Therefore, by using $y$ as an approximation for $\sin(y),$ our error is less that $\sin(y)/\cos(y) - 1\sin(y)$ $ = \sin(.003408846)/\cos(.003408846) = 0$ for seven decimal places.  We lose one significant figure due to accumulated round off error leaving six significant figures.\nl
$\sin(y) = y = .4363323/2^{7} = .003408846.$\nl
$\cos(.003408846) = \sqrt{1 - \sin^{2}(.003408846)} =  \sqrt{1 - .000011620232} = .999994189.$\nl
From Theorem 5.4.3, we have $\sin(2x) = 2\sin(x)\cos(x).$\nl
From Theorem 5.1.2, ID 4b, we have $\cos(x) = \sqrt{1 - \sin^{2}(x)}.$\nl
$\sin(2\cdot y) = 2\sin(.003408846)\cdot\cos(.003408846)$\nl
\indent $= 2\cdot(.003408846\cdot.99997675953) = .006817653.$\nl
$\cos(2\cdot y) = \sqrt{1 - \sin^{2}(2\cdot.003408846)} = \sqrt{1 - .0000464803908}$\nl
$ = .99997675953.$\nl
$\sin(2^{2}\cdot y) = 2\sin(2\cdot.003408846)\cdot\cos(2\cdot.003408846)$\nl
$ = 2\cdot(.006817653\cdot.99997675953) = .0013634989.$\nl
$\cos(2^{2}\cdot y) = \sqrt{1 - \sin^{2}(2^{2}\cdot.003408846)} = \sqrt{1 - .00018591292}$\nl
$ = .999907039218.$\nl
$\sin(2^{3}\cdot y) = 2\sin(2^{2}\cdot.003408846)\cdot\cos(2^{2}\cdot.003408846)$\nl
$ = 2\cdot(.0013634989\cdot.999907039218) = .0027267443.$\nl
$\cos(2^{3}\cdot y) = \sqrt{1 - \sin^{2}(2^{3}\cdot.003408846)} = \sqrt{1 - .00074351343}$\nl
$ = .9999628174.$\nl
%$\sin(2^{2}\cdot y) = 2\sin(2\cdot.003408846)\cdot\cos(2\cdot.003408846) = 2\cdot(.006817653\cdot.999994189) = %.0013634989.$\nl
%$\cos(2^{2}\cdot y) = \sqrt{1 - \sin^{2}(2^{2}\cdot.003408846)} = \sqrt{1 - .00018591292} = .999907039218.$\nl
%$\sin(2^{3}\cdot y) = 2\sin(2^{2}\cdot.003408846)\cdot\cos(2^{2}\cdot.003408846) = %2\cdot(.0013634989\cdot.9999070392184189) = .0027267443.$\nl
%$\cos(2^{3}\cdot y) = \sqrt{1 - \sin^{2}(2^{3}\cdot.003408846)} = \sqrt{1 - .00074351343} = .9999628174.$\nl
$\sin(2^{4}\cdot y) = 2\sin(2^{3}\cdot.003408846)\cdot\cos(2^{3}\cdot.003408846)$\nl
$ = 2\cdot(.0027267443\cdot.9999628174) = .0545146079.$\nl
$\cos(2^{4}\cdot y) = \sqrt{1 - \sin^{2}(2^{4}\cdot.003408846)} = \sqrt{1 - .002971842463}$\nl
$ = .998512973.$\nl
$\sin(2^{5}\cdot y) = 2\sin(2^{4}\cdot.003408846)\cdot\cos(2^{4}\cdot.003408846)$\nl
$ = 2\cdot(.0545146079\cdot.998512973) = .10886708.$\nl
$\cos(2^{5}\cdot y) = \sqrt{1 - \sin^{2}(2^{5}\cdot.003408846)} = \sqrt{1 - .000185204251}$\nl
$ = .994056315.$\nl
$\sin(2^{6}\cdot y) = 2\sin(2^{5}\cdot.003408846)\cdot\cos(2^{5}\cdot.003408846)$\nl
$ = 2\cdot(.10886708\cdot.994056315) = .21644003.$\nl
$\cos(2^{6}\cdot y) = \sqrt{1 - \sin^{2}(2^{6}\cdot.003408846)} = \sqrt{1 - .04684622864}$\nl
$ = .976295915.$\nl
$\sin(2^{7}\cdot y) = 2\sin(2^{6}\cdot.003408846)\cdot\cos(2^{6}\cdot.003408846)$\nl
$ = 2\cdot(.21644003\cdot.976295915) = .422619033.$\nl
$\cos(2^{7}\cdot y) = \sqrt{1 - \sin^{2}(2^{7}\cdot.003408846)} = \sqrt{1 - .178606847}$\nl
$ = .906307527.$\nl


Therefore, $\sin 25\dg = sin(\angle) = \sin(y\cdot2^{7}) = \sin(.003408846\cdot2^{7})$\nl
$ = .422619033.$\nl
$\cos 25\dg = \cos(\angle) = \cos(y\cdot2^{7}) = \cos(.003408846\cdot2^{7}) = .906307527.$\nl

By adding more iteration, $\sin$ and $\cos$ may be calculated to any desired degree of accuracy.  Definition 5.1.5 can be used to find any other trigonometric function.\nl




\newpage
%symbols
\lhead{}
\begin{center}
\textbf{Table of Symbols}
\end{center}
$ $\nl
%\thispagestyle{empty}
\begin{tabular}{l l c}
\textbf{symbol}&\textbf{meaning}&\textbf{page of first appearance}\\
$\sim p$& not $p$ &1\\
$p\up q$& $p$ and $q$ &1\\
$p\dn q$ &$p$ or $q$& 1 \\
$p\ra q$& $p$ implies $q$ &1\\
$p \lra q$&$p$ if and only if $q$& l\\
$t$ &a theorem &6 \\
$f$ &$\sim t$ &6 \\
$\forall $ &the universal quantifier (for each)& 11\\
$\exists$&the existential quantifier (there exists) & 11 \\
$x\in S$ &$x$ is an element of the set $S$& 16 \\
$x\notin S$ &$x$ is not an element of the set $S$& 16 \\
$\{x \mid p(x)\}$ &the set to which $x$ belongs if and only if $p(x)$& 16 \\
$x = y$ &$x$ equals $y$ &16 \\
$A\subseteq B$ &$A$ is a subset of $B$ &16 \\
$A \subset B$ &$A$ is a proper subset of $B$& 17 \\
$A \neq B$ &$A$ is not equal to $B$ &17 \\
$(a,b)$&an ordered pair &17 \\
$(a,b)$& open interval from a to b &37 \\
$(a,b)$&a complex number &99\\
$(a,b,c)$ &an ordered triple &17\\
$A\times B$ &the Cartesian product of $A$ and $B$ &17\\
$x\sim y$ &$(x,y)$ is in the binary relation $\sim$ &17 \\
$D_{f}$ &the domain of $f$& 19 \\
$R_{f}$ &the range of $f$& 19\\
$f(x)$ &the second term of the ordered pair &\\
& of the function $f$ whose first term is $x$ &19 \\
$f:A\ra B$ &$f$ is a function from $A$ into $B$ &19\\
$f:A\onto b$ &$f$ is a function from $A$ onto $B$& 19 \\
$x\ast y$ &the second term of the ordered pair of the &\\
&binary operation $\ast$ whose first term is $(x,y)$ &19 \\
$A\cap B$ &$A$ intersection $B$ &21 \\
$A\cup B$& $A$ union $B$& 21 \\
$\ns$&the empty set &21 \\
$\Pi A$& the power set of $A$ &21 \\
\end{tabular}

\newpage
\lhead{}
\begin{tabular}{l l r}
$A-B$& the difference of $A$ and $B$& 21 \\
$A^{'}$& the complement of $A$ &21 \\
\textbf{U} &the universal set .& 21\\
$\cup G$ &the union of the elements $G$& 21\\
$\cap G$ &the intersection of the elements of $G$ &21\\
$-x$ &the additive inverse of $x$& 27 \\
$x^{-1}$ &the multiplicative inverse of $x$& 27\\
$x - y$ &$x + (-y)$ &29\\
$x\div y$& $x\cdot y^{-1}$& 29 \\
$\frac{x}{y}$ &$x\cdot y^{-1}$& 29\\
$x > y$& $x$ is greater than $y$& 33 \\
$x < y$& $x$ is less than $y$& 33 \\
$x \geq y$ &$x$ is greater than or equal to $y$& 33\\
$x \leq y$ &$x$ is less than or equal to $y$& 33\\
$[a,b]$& the closed interval from $a$ to $b$ &37 \\
$[a,b$), $(a,b]$&half open intervals from $a$ to $b$ &37 \\
%& half open intervals from a to b &	37\\
$[a, \infty ), (a,\infty)$  &infinite intervals &38 \\
$(-\infty , a], (-\infty,a)$&infinite intervals&38\\
$(-\infty , \infty )$&infinite interval&38\\
$\abs{x}$&the absolute value of $x$ &39,98 \\
$\abs{x}$& the norm of $x$ &106 \\
$\zp$ &the set of positive integers& 44 \\
$Z$ &the set of integers& 46 \\
$Z_{a}$ &the set of integers greater than or equal to $a$ &46\\
$Z_{m}^{n}$ &the set of integers in the interval $[m,n]$& 48 \\
$\sum_{i=1}^{n}a_{i}$&the sum of $a_{1}, a_{2}, \dots ,a_{n}$ & 48\\
$Q$ &the set of rational numbers& 50 \\
\underline{c}&the constant function &53\\
\end{tabular}
\newpage
\lhead{}
\begin{tabular}{l l r}
$n!$&$n$ factorial& 59 \\
$\binom{n}{k}$& binomial coefficient& 59\\
$\prod_{i=1}^{n}a_{i}$&the product of $a_{1}, a_{2}, \dots  ,a_{n}$& 58\\
glb $S$ &greatest lower bound of $S$& 62 \\
lub $S$ &least upper bound of $S$& 62 \\
$R$ & the set of real numbers & 62\\
$P$ &the set of positive numbers& 65 \\
$log_{a}$ &logarithmic function to the base $a$& 65 \\
$L$ &logarithmic function to the base $e$ &65,101 \\
$E$& exponential function& 67,101 \\
$log_{a}x$& log of $x$ to the base $a$ &65,101\\
$a^{x}$& $a$ to the $x$ power& 67,101 \\
$\sqrt[n]{a}$  & the $n^{\text{th}}$ root of $a$ &69 \\
$d(u, v)$& the distance from $u$ to $v$& 78\\
\textbf{0}& the zero vector &78 \\
$c_{1}$& the unit circle &78 \\
$\alpha$& the wrapping function& 78 \\
$\pi$ &pi& 78 \\
$x^{\circ} $& $x$ degrees& 90\\
$Im~x$& the imaginary part of $x$ &98 \\
$Re~x$& the real part of $x$& 98 \\
$\overline{x}$ &the complex conjugate of $x$& 98 \\
$\overleftrightarrow{xy}$& the line xy& 106 \\
$\overrightarrow{xy}$ &the ray $xy$& 106 \\
$\overline{xy}$ &the segment $xy$ &106\\
$\angle xyz$& the angle $xyz$& 106\\
$\triangle xyz$& the triangle $xyz$& 106 \\
$m \angle xyz$&the measure of the angle $xyz$ &107\\
\end{tabular}
\newpage
\lhead{}
%Index
\begin{center}
\textbf{Index}
\end{center}
$ $\nl
\begin{tabular}{l r l r}
Absolute value &39,98 &Equivalence, justification &5,13\\
And (connective) &  1 &Existential quantifier &11\\
Angle  &   106 & Exponent & 32,44,67\\
\indent right & 107 &  Factorial &59\\
Arc Functions& 95& Field &26\\
Axioms && \indent complete ordered &62\\
\indent predicate calculus& 11,12& \indent ordered& 33\\
\indent statement calculus& 4,5 &  Finite Set& 48\\
Biconditional &1 &Formula& 11,12\\
Binary Operation& 19 &Function &19\\
Binary Relation &17 &\indent algebra of &51\\
Binomial Theorem &59& \indent arc& 95\\
Bound &&\indent constant& 53\\
\q greatest lower &62 &\q cycle of &85\\
\q least upper &62 &\q decreasing &53,56\\
\q lower &62 &\q domain of &19\\
\q upper &62 &\q increasing &53,56\\
Bounded Set& 62 &\q inverse &55\\
Cartesian product &17 &\q inverse trigonometric &95\\
Collinear points &106& \q one-to-one &48\\
Complement of a set& 21 &\q logarithmic &65,101\\
Complete Ordered Field& 62&\q period of &85\\
Complex number system &98 &\q range of& 19\\
Complex conjugate &98 &\q trigonometric& 78\\
Composite statement& 1& Generalization, justification &13 \\
Conclusion &5& Greatest lower bound &62\\
Conjunction& 1& Hypothesis &5\\
Connectives& 1 &Identity &27,53\\
Cosecant& 78 &Implication &1\\
Cosine &78& Inductive set &44\\
\q graph of& 93 &Infinite sequence &48\\
\q Law of &108 &Infinite set &48\\
Cotangent& 78 &Inner product space &104\\
Cycle of a Function& 85 &Integer& 46\\
Decimal &57 &Intercept &111\\
Definition& 1 &Intersection of sets &21\\
Degree &90 &Intervals &37\\
Difference of sets& 21 &Inverse &27,55,96\\
Disjunction& 1 &Irrational number &64\\
Divisibility& 47& Justification &5,12\\
Domain of a function& 19 &Law of Cosines &108\\
Empty set &21 &Law of Sines &108\\
Equation of a set& 110 &Least upper bound &62\\
Equality &&Line &106\\
\q of objects &16 &\q slope of &110\\
%\q of sets &16&Logarithmic functions& 65,101\\
\q of sets &16&Linear equations&117\\
&&Logaritmic functions&65,101\\

\end{tabular}

\newpage
%2nd page of Index
\lhead{}
\begin{tabular}{l r l r}
Lower Bound &62&\q empty &21\\
Matrices&113&\q equality of& 16\\
Measure of an angle& 107& \q equation for& 110\\
Modens Ponens& 5,12&\q finite& 48\\
Multiplicity of a root of a&& \q infinite &48\\
\q polynomial& 76&\q  intersection &21\\
Negation& 1& \q notation &16\\
Normed linear space& 106&\q power& 21\\
Not free & 11&\q proper subset& 17\\
Number& & \q subset& 16\\
\q irrational &64& \q truth &16\\
\q rational& 50&\q union of& 21\\
\q real& 62 &\q universal &21\\
Or (connective)& 1& Sine& 78\\
Ordered field& 33 &\q graph of &93\\
Ordered pair &17& \q Law of &108\\
Ordered triple& 17 &Slope of a line &110\\
Parallel lines& 111& Space&\\
Partial fractions& 75& \q inner product& 104\\
Period of a function &85& \q  normed linear& 106\\
Perpendicular lines &111& \q vector &104\\
Pi $(\pi)$ &78& Statement, definition of& 1\\
Polynomial& 73& Subset& 16\\
Positive integer& 44& Substatement &2\\
Power set& 21 &Summation notation& 48\\
Product && Synthetic division &75\\
\q Cartesian& 17 &Tangent &78\\
\q inner product space& 104& \q graph of& 93\\
Proof of a theorem &5,12& Tautology& 3\\
Proper subset& 17 &Theorem, definition of& 5,12\\
Quadratic Formula&76&Triangle&106\\
%Quantifiers& 11 &Triangle &106\\ %out
Quantifiers&11&Trigonometric Functions&78\\
%Range of a function &19 &Trigonometric functions& 78\\ %out
Range of a Function&19&\q table of values&88\\
%Rational number &50 &\q table of values& 88\\ %out
Rational Number&50&Truth Set&16\\
%Ray& 106& Truth set& 16\\ %out
Ray&106&Union of Sets&21\\
%Real number system& 62& Union of sets& 21\\ %out
Real number system&62&Universal quantifiers&11\\
%&&Universal quantifiers &11\\ %out
Root &&Universal set &21\\
\q of a polynomial& 75& Upper bound &62\\
\q of a number& 102& Variable& 11\\
Secant& 78& Vector space& 104\\
Segment& 106& $x$-axis &111\\
Sequence, infinite& 48& $x$-intercept& 111\\
Set && $y$-axis& 111\\
\q bounded &62& $y$-intercept &111\\
\q complement of &21&&\\
\q difference of &21&&\\
\end{tabular}


%\end{large}







%\end{document}




\end{large}

\end{document}


